% Leçon 41 — Expression orale : rédaction d’une lettre formelle — Chinois 1ère A — Cours signé OnBuch+
% Programme officiel MINESEC. Compiler avec : tectonic ZHA41-oral-redaction-dune-lettre-formelle.tex
\newcommand{\DOCMATIERE}{Chinois}
\newcommand{\DOCNIVEAU}{1ère A}
\newcommand{\DOCMODULE}{Module 5 : Mass médias et télécommunication}
\newcommand{\DOCLECON}{ZHA41}
\newcommand{\DOCTITRE}{Leçon 41 — Expression orale : rédaction d’une lettre formelle}
% OnBuch+ — préambule « cours signés » ENRICHI (compilé avec tectonic / XeLaTeX)
% Système visuel « Pop » d'OnBuch+ : fond crème (jamais blanc), bordures encre
% épaisses, ombres « dures » décalées, titres Archivo Black en capitales.
% Version MATHÉMATIQUES : graphes (pgfplots/tikz), boîtes pédagogiques
% (définition, propriété, méthode, attention, le savais-tu, exercice résolu,
% exercice, TP, bilan), page de garde et signature OnBuch+ ancrée en bas.
%
% Macros attendues dans main.tex AVANT \input{preamble} :
%   \DOCMATIERE  \DOCNIVEAU  \DOCMODULE  \DOCLECON (numéro)  \DOCTITRE
\documentclass[11pt]{article}

\usepackage{fontspec}
\usepackage[french]{babel} % français = langue principale ; le chinois via \zh{..} / textebox (xeCJK)
\usepackage{xeCJK}
\usepackage{pifont} % \ding{51} \ding{55} utilisés par les modèles
\usepackage[a4paper,margin=2cm,top=3.4cm,bottom=2.8cm,headheight=1.4cm,footskip=1.3cm]{geometry}
\usepackage{amsmath,amssymb,amsfonts}
\usepackage{mathtools} % \xmapsto, \coloneqq... souvent utilisés par les modèles
\usepackage{cancel} % simplifications barrées (bilans de forces)
% \abs{z} (module, valeur absolue) et \norm{u} (norme), utilisés par les modèles
\providecommand{\abs}{}\let\abs\relax\DeclarePairedDelimiter{\abs}{\lvert}{\rvert}
\providecommand{\norm}{}\let\norm\relax\DeclarePairedDelimiter{\norm}{\lVert}{\rVert}
\usepackage[table]{xcolor} % [table] : fournit aussi \rowcolors (alternance de lignes)
\usepackage{pagecolor}
\usepackage{tikz}
\usetikzlibrary{arrows.meta,positioning,calc,shapes.geometric,shapes.misc,shapes.arrows,decorations.pathmorphing,decorations.pathreplacing,decorations.markings,patterns,shadows,mindmap,trees,fit,backgrounds,matrix,intersections,angles,quotes}
\usepackage{pgfplots}
\usepackage[version=4]{mhchem} % équations nucléaires \ce{^{235}_{92}U}
\usepackage[european,siunitx]{circuitikz} % schémas électriques
\pgfplotsset{compat=1.18}
\usepgfplotslibrary{fillbetween,groupplots} % aires entre courbes (calcul intégral)
\usepackage{enumitem}
\usepackage{fancyhdr}
% [most] charge toutes les bibliothèques tcolorbox (listings, minted, raster,
% documentation...) et gonfle inutilement la mémoire TeX sur un document long
% (~300 boîtes) : on ne charge que ce qui est réellement utilisé.
\usepackage[skins,breakable]{tcolorbox}
\usepackage{graphicx}
\usepackage{colortbl}
\usepackage{booktabs}
\usepackage{tabularx}
\usepackage{diagbox} % case d'en-tête diagonale des tableaux à double entrée
% Colonnes extensibles centrée / gauche / droite (C, L, R), souvent utilisées par les modèles
\newcolumntype{C}{>{\centering\arraybackslash}X}
\newcolumntype{L}{>{\raggedright\arraybackslash}X}
\newcolumntype{R}{>{\raggedleft\arraybackslash}X}
\usepackage{array}
\usepackage{multicol}
\usepackage{multirow}
\usepackage{siunitx}
% parse-numbers=false : accepte aussi \SI{2,5 \times 10^{-3}}{...}
\sisetup{locale=FR,output-decimal-marker={,},group-minimum-digits=4,parse-numbers=false}
\DeclareSIUnit\ohm{\ensuremath{\symup{\Omega}}} % U+2126 absent de la police
\DeclareSIUnit\division{div} % réglages d'oscilloscope (ms/div, V/div)
\DeclareSIUnit\tr{tr} % tours (vitesse de rotation, ex. \SI{1500}{\tr\per\minute})
\DeclareSIUnit\molar{M} % concentration molaire (ex. \SI{3}{\molar}, \SI{1}{\micro\molar})
\DeclareSIUnit\an{an} % année (le modèle confond parfois avec \year, un compteur TeX) — ex. \SI{5}{\kilo\meter\per\an}
\DeclareSIUnit\FCFA{FCFA} % franc CFA (ex. \SI{500}{\FCFA\per\kilo\gram})
\DeclareSIUnit\degreeBrix{\textdegree Brix} % teneur en sucre d'un sirop/jus (réfractométrie)
\DeclareSIUnit\month{mois} % le modèle utilise parfois \month, un compteur TeX, comme unité de durée
\usepackage{qrcode}
% \degree : commande fréquemment utilisée par les modèles pour un angle ou une
% température en dehors de \SI{}{} (non fournie par les paquets chargés).
\providecommand{\degree}{^{\circ}}
% \qed (fin de démonstration), utilisé par les modèles sans amsthm
\providecommand{\qed}{\hfill$\square$}
% \barre : double barre des valeurs interdites dans un tableau de variations (tkz-tab)
\providecommand{\barre}{\ensuremath{\|}}
% environnement proof (amsthm non chargé) utilisé par les modèles
\ifdefined\proof\else\newenvironment{proof}[1][Démonstration]{\par\noindent\textit{#1.}\ }{\qed\par}\fi
\usepackage{lastpage}
\usepackage{hyperref}
\hypersetup{hidelinks}

% ============================================================
% Palette « Pop » OnBuch+ — mêmes hex que la classe P (plus_ui.dart)
% ============================================================
\definecolor{poporange}{HTML}{F59321}
\definecolor{popdark}{HTML}{E07A0C}
\definecolor{popcream}{HTML}{FFF6E8}
\definecolor{popink}{HTML}{1C1714}
\definecolor{poppurple}{HTML}{7A5AE0}
\definecolor{popgreen}{HTML}{2E9E6A}
\definecolor{poppink}{HTML}{E0457B}
\definecolor{popblue}{HTML}{2F7DE1}
\definecolor{popgold}{HTML}{C98A12}
\definecolor{popmuted}{HTML}{8A7F76}
\definecolor{popline}{HTML}{EADFD2}
\definecolor{popgray}{HTML}{8A7F76}
\colorlet{popgrey}{popgray}
% Alias tolérants (le modèle nomme parfois les couleurs différemment)
\colorlet{popwhite}{white}
\colorlet{popblack}{popink}
\colorlet{popred}{poppink}
\colorlet{popyellow}{popgold}
% Teintes claires (fonds de tableaux/encadrés) — le modèle nomme parfois
% les couleurs avec un suffixe L (ex. popblueL) sans les avoir définies.
\colorlet{poporangeL}{poporange!20}
\colorlet{popdarkL}{popdark!20}
\colorlet{poppurpleL}{poppurple!20}
\colorlet{popgreenL}{popgreen!20}
\colorlet{poppinkL}{poppink!20}
\colorlet{popblueL}{popblue!20}
\colorlet{popgoldL}{popgold!20}
\colorlet{popredL}{popred!20}
\colorlet{popgrayL}{popgray!20}
% Teintes claires pour fonds de tableaux / zones
\colorlet{popblueL}{popblue!10!white}
\colorlet{poporangeL}{poporange!14!white}
\colorlet{popgreenL}{popgreen!12!white}
\colorlet{poppurpleL}{poppurple!12!white}
\colorlet{poppinkL}{poppink!12!white}
\colorlet{popgoldL}{popgold!14!white}

\pagecolor{popcream}

% ============================================================
% Typographie
% ============================================================
\defaultfontfeatures{Ligatures=TeX}
\setmainfont{PlusJakartaSans-Medium.ttf}[Path=../../../fonts/,BoldFont=PlusJakartaSans-Bold.ttf]
% Symboles, API et emojis absents de la police principale : repli sur DejaVu Sans (caractères actifs)
\newfontface\symfont{DejaVuSans.ttf}[Path=../../../fonts/]
\makeatletter
\newcommand\symactive[1]{\catcode#1=13 \begingroup\lccode`\~=#1 \lowercase{\endgroup\def~}{{\symfont\char#1}}}
\newcommand\symblank[1]{\catcode#1=13 \begingroup\lccode`\~=#1 \lowercase{\endgroup\def~}{}}
\newcommand\symhyph[1]{\catcode#1=13 \begingroup\lccode`\~=#1 \lowercase{\endgroup\def~}{\mbox{-}}}
\newcount\symc
\newcommand\symrange[2]{\symc=#1 \@whilenum\symc<#2 \do{\expandafter\symactive\expandafter{\the\symc}\advance\symc by 1 }}
\newcommand\symblankrange[2]{\symc=#1 \@whilenum\symc<#2 \do{\expandafter\symblank\expandafter{\the\symc}\advance\symc by 1 }}
\makeatother
\symrange{"0250}{"02FF}\symactive{"2713}\symactive{"2714}\symactive{"2717}\symactive{"2718}\symactive{"2610}\symactive{"2611}\symactive{"2612}
\symrange{"2600}{"27BF}
\catcode"2705=13 \begingroup\lccode`\~="2705 \lowercase{\endgroup\def~}{{\symfont\char"2713}}\symblank{"26BD}\symblank{"26A1}
\symrange{"2460}{"257F}\symactive{"223C}\symactive{"2248}\symactive{"2260}
\symactive{"01F8}\symactive{"01F9}\symactive{"1E3E}\symactive{"1E3F}
\symhyph{"2011}\symhyph{"2E1A}\symblankrange{"1F300}{"1FAFF}\symblank{"FE0F}
% Caractères chinois : WenQuanYi Zen Hei (sous-ensemble GB2312 + pinyin), gras/italique simulés
\setCJKmainfont{WenQuanYiZenHei-subset.ttf}[Path=../../../fonts/,AutoFakeBold=3,AutoFakeSlant=0.2]
\xeCJKsetup{CJKspace=false}
% Pinyin : voyelles à caron (ǎ ǐ ǒ ǔ ǖ ǘ ǚ ǜ) absentes de la police principale -> caractères actifs
\newfontface\pyfont{WenQuanYiZenHei-subset.ttf}[Path=../../../fonts/]
\newcommand\pyactive[1]{\catcode#1=13 \begingroup\lccode`\~=#1 \lowercase{\endgroup\def~}{{\pyfont\char#1}}}
\pyactive{"01CD}\pyactive{"01CE}\pyactive{"01CF}\pyactive{"01D0}\pyactive{"01D1}\pyactive{"01D2}\pyactive{"01D3}\pyactive{"01D4}\pyactive{"01D5}\pyactive{"01D6}\pyactive{"01D7}\pyactive{"01D8}\pyactive{"01D9}\pyactive{"01DA}\pyactive{"01DB}\pyactive{"01DC}
% Maths en sans-serif (Fira Math) avec chiffres et lettres droites de Plus
% Jakarta, pour que les formules se fondent dans le texte.
\usepackage[math-style=ISO,bold-style=ISO]{unicode-math}
\setmathfont{FiraMath-Regular.otf}[Path=../../../fonts/]
\setmathfont{PlusJakartaSans-Medium.ttf}[Path=../../../fonts/,range={up/num,up/{Latin,latin},\mathup}]
\usepackage{icomma} % virgule décimale sans espace en mode math
% Caractères mathématiques Unicode absents des polices de texte (Plus Jakarta,
% Archivo Black) : rendus par leur équivalent mathématique, en texte comme en
% formule — sinon ils s'affichent en carré vide (ex. « Arithmétique dans ℤ »).
\usepackage{newunicodechar}
\newunicodechar{ℝ}{\ensuremath{\mathbb{R}}}
\newunicodechar{ℤ}{\ensuremath{\mathbb{Z}}}
\newunicodechar{ℂ}{\ensuremath{\mathbb{C}}}
\newunicodechar{ℕ}{\ensuremath{\mathbb{N}}}
\newunicodechar{ℚ}{\ensuremath{\mathbb{Q}}}
\newunicodechar{𝔻}{\ensuremath{\mathbb{D}}}
\newunicodechar{α}{\ensuremath{\alpha}}
\newunicodechar{β}{\ensuremath{\beta}}
\newunicodechar{γ}{\ensuremath{\gamma}}
\newunicodechar{δ}{\ensuremath{\delta}}
\newunicodechar{ε}{\ensuremath{\varepsilon}}
\newunicodechar{θ}{\ensuremath{\theta}}
\newunicodechar{λ}{\ensuremath{\lambda}}
\newunicodechar{μ}{\ensuremath{\mu}}
\newunicodechar{σ}{\ensuremath{\sigma}}
\newunicodechar{τ}{\ensuremath{\tau}}
\newunicodechar{φ}{\ensuremath{\varphi}}
\newunicodechar{ψ}{\ensuremath{\psi}}
\newunicodechar{ω}{\ensuremath{\omega}}
\newunicodechar{Σ}{\ensuremath{\Sigma}}
% majuscules produites par \MakeUppercase dans les titres (α-aminés -> Α)
\newunicodechar{Α}{\ensuremath{\alpha}}
\newunicodechar{Β}{\ensuremath{\beta}}
\newunicodechar{Γ}{\ensuremath{\Gamma}}
\newunicodechar{Δ}{\ensuremath{\Delta}}
\newunicodechar{Ε}{\ensuremath{\varepsilon}}
\newunicodechar{Θ}{\ensuremath{\Theta}}
\newunicodechar{Λ}{\ensuremath{\Lambda}}
\newunicodechar{Μ}{\ensuremath{\mu}}
\newunicodechar{Τ}{\ensuremath{\tau}}
\newunicodechar{Φ}{\ensuremath{\Phi}}
\newunicodechar{Ψ}{\ensuremath{\Psi}}
\newunicodechar{Ω}{\ensuremath{\Omega}}
\newunicodechar{↦}{\ensuremath{\mapsto}}
\newunicodechar{∈}{\ensuremath{\in}}
\newunicodechar{∉}{\ensuremath{\notin}}
\newunicodechar{∧}{\ensuremath{\wedge}}
\newunicodechar{⇔}{\ensuremath{\Leftrightarrow}}
\newunicodechar{⟺}{\ensuremath{\Longleftrightarrow}}
\newunicodechar{⇒}{\ensuremath{\Rightarrow}}
\newunicodechar{⟹}{\ensuremath{\Longrightarrow}}
\newunicodechar{∘}{\ensuremath{\circ}}
\newunicodechar{∪}{\ensuremath{\cup}}
\newunicodechar{∩}{\ensuremath{\cap}}
\newunicodechar{≡}{\ensuremath{\equiv}}
\newunicodechar{⊂}{\ensuremath{\subset}}
\newunicodechar{⊥}{\ensuremath{\perp}}
\newunicodechar{∃}{\ensuremath{\exists}}
\newunicodechar{∀}{\ensuremath{\forall}}
\newunicodechar{∛}{\ensuremath{\sqrt[3]{\,}}}
\newunicodechar{∜}{\ensuremath{\sqrt[4]{\,}}}
\newunicodechar{⃗}{\ensuremath{{}^{\to}}}
\newunicodechar{ⁿ}{\ifmmode^{n}\else\textsuperscript{\ensuremath{n}}\fi}
\newunicodechar{ˣ}{\ifmmode^{x}\else\textsuperscript{\ensuremath{x}}\fi}
\newunicodechar{ᵇ}{\ifmmode^{b}\else\textsuperscript{\ensuremath{b}}\fi}
\newunicodechar{ʳ}{\ifmmode^{r}\else\textsuperscript{\ensuremath{r}}\fi}
\newunicodechar{ᵘ}{\ifmmode^{u}\else\textsuperscript{\ensuremath{u}}\fi}
\newunicodechar{ʸ}{\ifmmode^{y}\else\textsuperscript{\ensuremath{y}}\fi}
\newunicodechar{ᵃ}{\ifmmode^{a}\else\textsuperscript{\ensuremath{a}}\fi}
\newunicodechar{ᵉ}{\ifmmode^{e}\else\textsuperscript{\ensuremath{e}}\fi}
\newunicodechar{ᵏ}{\ifmmode^{k}\else\textsuperscript{\ensuremath{k}}\fi}
\newunicodechar{ᵖ}{\ifmmode^{p}\else\textsuperscript{\ensuremath{p}}\fi}
\newunicodechar{ᴺ}{\ifmmode^{N}\else\textsuperscript{\ensuremath{N}}\fi}
\newunicodechar{ᶜ}{\ifmmode^{c}\else\textsuperscript{\ensuremath{c}}\fi}
\newunicodechar{ʰ}{\ifmmode^{h}\else\textsuperscript{\ensuremath{h}}\fi}
\newunicodechar{ᵠ}{\ifmmode^{q}\else\textsuperscript{\ensuremath{q}}\fi}
\newunicodechar{ₙ}{\ifmmode_{n}\else\textsubscript{\ensuremath{n}}\fi}
\newunicodechar{ₐ}{\ifmmode_{a}\else\textsubscript{\ensuremath{a}}\fi}
\newunicodechar{ᵢ}{\ifmmode_{i}\else\textsubscript{\ensuremath{i}}\fi}
\newunicodechar{ᵣ}{\ifmmode_{r}\else\textsubscript{\ensuremath{r}}\fi}
\newunicodechar{ₖ}{\ifmmode_{k}\else\textsubscript{\ensuremath{k}}\fi}
\newunicodechar{ᵦ}{\ifmmode_{\beta}\else\textsubscript{\ensuremath{\beta}}\fi}
\newunicodechar{ₓ}{\ifmmode_{x}\else\textsubscript{\ensuremath{x}}\fi}
\newfontfamily\popdisplay{ArchivoBlack-Regular.ttf}[Path=../../../fonts/]
\newfontfamily\poplogo{SpaceGrotesk-Bold.ttf}[Path=../../../fonts/]
\newfontfamily\popsemi{PlusJakartaSans-SemiBold.ttf}[Path=../../../fonts/]
\newfontfamily\popxbold{PlusJakartaSans-ExtraBold.ttf}[Path=../../../fonts/]
\newfontfamily\popxboldls{PlusJakartaSans-ExtraBold.ttf}[Path=../../../fonts/,LetterSpace=3.0]

\setlist[enumerate]{leftmargin=1.7em,itemsep=5pt,topsep=4pt}
\setlist[itemize]{leftmargin=1.7em,itemsep=4pt,topsep=4pt,label={\color{poporange}\textbullet}}
\setlength{\parindent}{0pt}
\setlength{\parskip}{5pt}
\renewcommand{\arraystretch}{1.25}

% pgfplots : style Pop par défaut
\pgfplotsset{
  every axis/.append style={axis line style={popink,line width=1pt},tick style={popink},
    grid style={popline,line width=0.5pt},label style={font=\small},tick label style={font=\footnotesize}},
  popcurve/.style={poporange,line width=1.8pt,no markers,smooth},
}
% Même style disponible dans un \draw TikZ simple (hors pgfplots)
\tikzset{popcurve/.style={poporange,line width=1.8pt}}
\tikzset{popgrid/.style={popline,very thin}}
\colorlet{popcurve}{poporange} % le nom du style est parfois utilisé comme couleur

% ============================================================
% Pill / squiggle
% ============================================================
\newcommand{\poppill}[3]{%
  \tikz[baseline=-0.55ex]\node[draw=popink,line width=1.2pt,rounded corners=2.6pt,%
    fill=#2,inner xsep=6.5pt,inner ysep=3pt]{%
    \popxboldls\fontsize{7}{7}\selectfont\color{#3}\MakeUppercase{#1}};%
}
\newcommand{\popsquiggle}[1]{%
  \tikz[baseline=-0.4ex]{
    \draw[#1,line width=2.6pt,line cap=round]
      plot [smooth] coordinates {(0,0.09) (0.34,0.01) (0.68,0.21) (1.02,0.01) (1.36,0.10)};
  }%
}
% Mot-clé surligné (marqueur orange sous le texte)
\newcommand{\cle}[1]{{\popxbold\color{popdark}#1}}

% ============================================================
% En-tête / pied
% ============================================================
\pagestyle{fancy}
\fancyhf{}
\renewcommand{\headrulewidth}{0pt}
\renewcommand{\footrulewidth}{0pt}
\lhead{\raisebox{-0.15ex}{\poplogo\fontsize{13}{13}\selectfont\color{popink}OnBuch\color{poporange}+}}
\rhead{\poppill{\DOCMATIERE\ \textbullet\ \DOCNIVEAU\ \textbullet\ Leçon \DOCLECON}{white}{popink}}
\lfoot{\popxbold\fontsize{7.6}{7.6}\selectfont\color{popmuted}\MakeUppercase{Cours signé OnBuch+}\ \textbullet\ {\popsemi\DOCTITRE}}
\rfoot{\tikz[baseline=-0.6ex]\node[rounded corners=8.5pt,fill=popink,minimum height=17pt,minimum width=17pt,inner xsep=5pt,inner sep=0]{\popxbold\fontsize{8}{8}\selectfont\color{popcream}\thepage\,/\,\pageref*{LastPage}};}

% ============================================================
% Titres
% ============================================================
\newcommand{\chaptitre}[2]{%
  \begin{tcolorbox}[enhanced,colback=poporange,colframe=popink,boxrule=2.6pt,arc=11pt,
      drop shadow=popink,left=15pt,right=15pt,top=12pt,bottom=11pt,before skip=3pt,after skip=18pt]
    \poppill{#2}{popink}{popcream}\par\vspace{9pt}
    {\raggedright\hyphenpenalty=10000\exhyphenpenalty=10000\popdisplay\fontsize{22}{24}\selectfont\color{popcream}\MakeUppercase{#1}\par}
  \end{tcolorbox}
}
% Section numérotée : pastille orange avec le numéro + titre Archivo
\newcounter{popsec}
\newcommand{\coursec}[1]{%
  \stepcounter{popsec}\setcounter{popsub}{0}%
  \par\vspace{16pt}\noindent
  \begin{minipage}{\linewidth}
  \tikz[baseline=-0.7ex]\node[draw=popink,line width=1.6pt,fill=poporange,rounded corners=4pt,minimum size=22pt,inner sep=0]{\popdisplay\fontsize{12}{12}\selectfont\color{popcream}\thepopsec};%
  \hspace{8pt}{\popdisplay\fontsize{15}{17}\selectfont\color{popink}\MakeUppercase{#1}}\par
  \vspace{4pt}\noindent{\color{poporange}\rule{36pt}{3.6pt}}
  \end{minipage}\par\nopagebreak\vspace{8pt}
}
\newcounter{popsub}
\newcommand{\courssub}[1]{%
  \stepcounter{popsub}%
  \par\vspace{9pt}\noindent{\popxbold\fontsize{11.5}{13}\selectfont\color{popink}{\color{poporange}\thepopsec.\thepopsub}\hspace{6pt}#1}\par\nopagebreak\vspace{4pt}
}

% ============================================================
% Boîtes pédagogiques — même recette : fond blanc, bordure épaisse colorée,
% ombre dure encre, badge plein. Titre optionnel [#1].
% ============================================================
\tcbset{popbase/.style={breakable,enhanced,colback=white,boxrule=2.3pt,arc=9pt,
  shadow={3pt}{-3pt}{0pt}{popink},left=14pt,right=14pt,top=11pt,bottom=11pt,before skip=11pt,after skip=15pt,
  fontupper=\popsemi\fontsize{10.6}{14.5}\selectfont\color{popink},
  fontlower=\popsemi\fontsize{10.6}{14.5}\selectfont\color{popink}}}
% Coche dessinée (le glyphe ✓ n'existe pas dans les polices du document)
\newcommand{\popcheck}[1]{\tikz[baseline=-0.5ex]\draw[#1,line width=1.6pt,line cap=round,line join=round](-0.13,0.0)--(-0.04,-0.09)--(0.13,0.1);}
\providecommand{\coche}{\popcheck{popgreen}}
\providecommand{\resultat}[1]{\par\smallskip\noindent\fcolorbox{popgreen}{popgreenL}{\parbox{\dimexpr\linewidth-2\fboxsep-2\fboxrule\relax}{\textbf{Résultat.} #1}}\par\smallskip}
\newcommand{\popboxhead}[3]{\poppill{#1}{#2}{white}\ifx\relax#3\relax\else\hspace{6pt}{\popxbold\fontsize{10.6}{12}\selectfont\color{popink}#3}\fi\par\vspace{7pt}}

\newtcolorbox{aretenir}[1][]{popbase,colframe=popgreen,before upper={\popboxhead{À retenir}{popgreen}{#1}}}
% Texte en chinois (document de lecture, dialogue, extrait) : cadre
% d'étude. Un titre optionnel (source, auteur) s'affiche dans l'en-tête.
\newtcolorbox{textebox}[1][]{popbase,colframe=popblue,before upper={\popboxhead{文本 · Texte}{popblue}{#1}},after upper={}}
% Marques ✓ / ✗ (forme correcte / fautive) souvent utilisées par les modèles
\providecommand{\cross}{\textcolor{poppink}{\ensuremath{\times}}}
\providecommand{\xmark}{\cross}\providecommand{\crossmark}{\cross}\providecommand{\cmark}{\ensuremath{\checkmark}}
% Ligne de réponse / justification à compléter (inventée par les modèles)
\providecommand{\justification}[1]{\par\noindent\textit{Justification~:} \dotfill\par}
% \textipa (package tipa non chargé) : la police ne couvre pas l'API -> simple passage
\providecommand{\textipa}[1]{#1}
% Soulignement (ulem non chargé) : \uline, \ul
\providecommand{\uline}[1]{\underline{#1}}\providecommand{\ul}[1]{\underline{#1}}
% \glossaire{\item ...} inventée par les modèles : liste de glossaire
\providecommand{\glossaire}[1]{\begin{itemize}#1\end{itemize}}
% \alert (beamer) inventée par les modèles : texte mis en évidence
\providecommand{\alert}[1]{\textbf{\textcolor{poppink}{#1}}}
% variantes ASCII d'ordinaux écrites en macro
\providecommand{\iere}{\textsuperscript{re}}\providecommand{\ieme}{\textsuperscript{e}}\providecommand{\ere}{\textsuperscript{re}}\providecommand{\eme}{\textsuperscript{e}}\providecommand{\er}{\textsuperscript{er}}\providecommand{\ie}{\textsuperscript{e}}
% Ordinaux français écrits en macro par les modèles : 1\ière, 3\ième
\newcommand{\ière}{\textsuperscript{re}}\newcommand{\ième}{\textsuperscript{e}}
% \exemple (inventée par les modèles) : étiquette « Exemple : »
\providecommand{\exemple}{\textbf{Exemple~:}\ }
% \square utilisable aussi hors mode math (cases à cocher)
\AtBeginDocument{\let\squareMath\square\renewcommand{\square}{\ensuremath{\squareMath}}}
% Mot ou phrase en chinois à l'intérieur d'un paragraphe français
\newcommand{\zh}[1]{\ifmmode\text{#1}\else{#1}\fi}
\let\eng\zh \let\esp\zh \let\all\zh % alias tolérés
% flèches d'intonation inventées par les modèles
\providecommand{\textsearrow}{}\renewcommand{\textsearrow}{\ensuremath{\searrow}}\providecommand{\textnearrow}{}\renewcommand{\textnearrow}{\ensuremath{\nearrow}}
% \fr{..} : mot français (inventé par les modèles)
\providecommand{\fr}[1]{\textit{#1}}
% zhbox : encadré de texte chinois inventé par les modèles
\newenvironment{zhbox}{\begin{quote}}{\end{quote}}
% \courssubsub : titre de niveau 3 inventé par les modèles
\providecommand{\courssubsub}[1]{\par\medskip\noindent\textbf{#1}\par\smallskip}
% \zhbf : texte chinois en gras inventé par les modèles
\providecommand{\zhbf}[1]{\textbf{#1}}
% \vs : « versus » inventé par les modèles
\providecommand{\vs}{\textit{vs}\ }
% \blank : case à compléter inventée par les modèles
\providecommand{\blank}[1][]{\underline{\hspace{1.6em}}}
\newtcolorbox{exemplebox}[1][]{popbase,colframe=poporange,before upper={\popboxhead{Exemple}{poporange}{#1}}}
\newtcolorbox{definition}[1][]{popbase,colframe=popblue,before upper={\popboxhead{Définition}{popblue}{#1}}}
\newtcolorbox{propriete}[1][]{popbase,colframe=poppurple,before upper={\popboxhead{Propriété}{poppurple}{#1}}}
\newtcolorbox{methode}[1][]{popbase,colframe=popgold,before upper={\popboxhead{Méthode}{popgold}{#1}}}
\newtcolorbox{attention}[1][]{popbase,colframe=poppink,before upper={\popboxhead{Attention}{poppink}{#1}}}
\newtcolorbox{experience}[1][]{popbase,colframe=popblue,colback=popblueL,before upper={\popboxhead{Expérience}{popblue}{#1}}}
% « Le savais-tu ? » : carte violette pleine (contexte, culture, Cameroun)
\newtcolorbox{savaistu}[1][]{popbase,colback=poppurple,colframe=popink,
  fontupper=\popsemi\fontsize{10.4}{14.2}\selectfont\color{white},
  before upper={\poppill{Le savais-tu\,?}{white}{poppurple}\ifx\relax#1\relax\else\hspace{6pt}{\popxbold\color{white}#1}\fi\par\vspace{7pt}}}
% Exercice résolu : énoncé en haut, \tcblower puis la solution sur fond crème
\newcounter{popexo}\newcounter{popexr}
\newtcolorbox{exoresolu}[1][]{popbase,colframe=poporange,colbacklower=popcream,
  segmentation style={popink,line width=1.2pt,dashed},
  before upper={\stepcounter{popexr}\popboxhead{Exercice résolu \thepopexr}{poporange}{#1}},
  before lower={\poppill{Solution}{popink}{popcream}\par\vspace{6pt}}}
% Exercice (non résolu en place) — la correction est dans la partie Corrigés
\newtcolorbox{exercice}[1][]{popbase,colframe=popink,
  before upper={\stepcounter{popexo}\popboxhead{Exercice \thepopexo}{popink}{#1}}}
% Corrigé
\newtcolorbox{corrige}[1][]{popbase,colframe=popgreen,colback=popgreenL,
  before upper={\popboxhead{Corrigé}{popgreen}{#1}}}
% Objectifs / prérequis / bilan
\newtcolorbox{objectifs}{popbase,colframe=popink,colback=poporangeL,
  before upper={\popboxhead{Objectifs de la leçon}{popink}{}}}
\newtcolorbox{prerequis}{popbase,colframe=popink,colback=popblueL,
  before upper={\popboxhead{Prérequis}{popblue}{}}}
\newtcolorbox{bilan}{popbase,colframe=popink,colback=white,boxrule=2.8pt,
  before upper={\popboxhead{Fiche bilan}{poporange}{L'essentiel à connaître pour l'examen}}}
% Figure encadrée avec légende
\newtcolorbox{popfigure}[1][]{enhanced,breakable=false,colback=white,colframe=popline,boxrule=1.4pt,arc=8pt,
  left=10pt,right=10pt,top=10pt,bottom=8pt,before skip=10pt,after skip=12pt,halign=center,
  lower separated=false,halign lower=center,
  fontlower=\popsemi\fontsize{9}{11.5}\selectfont\color{popmuted},
  #1}
\newcommand{\legende}[1]{\par\vspace{6pt}{\popsemi\fontsize{9}{11.5}\selectfont\color{popmuted}#1}}

% ============================================================
% Signature OnBuch+
% ============================================================
\newcommand{\poplogoinline}[1]{{\poplogo\fontsize{#1}{#1}\selectfont\color{popink}OnBuch\color{poporange}+}}
\newcommand{\signatureonbuch}{%
  \begin{tcolorbox}[enhanced,colback=popink,colframe=popink,boxrule=2.2pt,arc=10pt,
      drop shadow=poporange,left=14pt,right=14pt,top=10pt,bottom=10pt,before skip=0pt,after skip=0pt]
    \tikz[baseline=-0.7ex]\node[circle,fill=popgreen,draw=popcream,line width=1.3pt,minimum size=22pt,inner sep=0]{\popcheck{white}};%
    \hspace{9pt}\begin{minipage}[c]{0.62\linewidth}
      {\popxbold\fontsize{12}{13}\selectfont\color{popcream}Signé\ \textbullet\ {\poplogo OnBuch\color{poporange}+}}\par\vspace{2pt}
      {\raggedright\popsemi\fontsize{8}{10.5}\selectfont\color{popline}\MakeUppercase{Équipe pédagogique OnBuch+}\\\MakeUppercase{Conforme au programme officiel MINESEC}\par}
    \end{minipage}\hfill
    {\poplogo\fontsize{22}{22}\selectfont\color{popcream}OnBuch\color{poporange}+}
  \end{tcolorbox}%
}

% ============================================================
% Page de garde
% #1 titre · #2 matière · #3 niveau · #4 module · #5 n° leçon · #6 durée ·
% #7 description / objectifs (texte ou itemize)
% ============================================================
\newcommand{\pagegarde}[7]{%
  \thispagestyle{empty}
  \newgeometry{a4paper,margin=1.7cm,top=1.6cm,bottom=1.4cm}%
  \begin{tikzpicture}[remember picture,overlay]
    \fill[poporange!18] ([xshift=-3.2cm,yshift=-3.6cm]current page.north east) circle (4.6cm);
    \fill[poppurple!14] ([xshift=2.4cm,yshift=7.6cm]current page.south west) circle (3.8cm);
  \end{tikzpicture}%
  {\poplogo\fontsize{30}{30}\selectfont\color{popink}OnBuch\color{poporange}+}\hfill
  \raisebox{0.6ex}{\poppill{Cours signé \textbullet\ Premium}{popink}{popcream}}\par
  \vspace{1pt}\popsquiggle{poppurple}\par
  \vspace{8pt}
  \begin{tcolorbox}[enhanced,colback=poporange,colframe=popink,boxrule=2.8pt,arc=13pt,
      drop shadow=popink,left=18pt,right=18pt,top=12pt,bottom=13pt,after skip=10pt]
    \poppill{#2 \textbullet\ #3}{popink}{popcream}\hspace{5pt}\poppill{#4}{white}{popink}\par\vspace{8pt}
    {\popdisplay\fontsize{14}{14}\selectfont\color{popink}LEÇON #5}\par\vspace{4pt}
    {\raggedright\hyphenpenalty=10000\exhyphenpenalty=10000\popdisplay\fontsize{25}{27}\selectfont\color{popcream}\MakeUppercase{#1}\par}
    \vspace{8pt}
    {\popxbold\fontsize{9.5}{9.5}\selectfont\color{popink}\MakeUppercase{Durée conseillée : #6}}
  \end{tcolorbox}
  \begin{tcolorbox}[enhanced,colback=white,colframe=popink,boxrule=2.2pt,arc=10pt,
      drop shadow=popink,left=16pt,right=16pt,top=10pt,bottom=10pt,after skip=8pt]
    \poppill{Dans cette leçon}{poppurple}{white}\par\vspace{5pt}
    \popsemi\fontsize{9.3}{12.6}\selectfont\color{popink}#7
  \end{tcolorbox}
  \noindent\begin{minipage}[t]{0.487\linewidth}
    \begin{tcolorbox}[enhanced,colback=popgreen,colframe=popink,boxrule=2.2pt,arc=10pt,
        drop shadow=popink,left=11pt,right=11pt,top=10pt,bottom=10pt,height=3.1cm,valign=center]
      \tcbox[colback=white,colframe=popink,boxrule=1.3pt,arc=5pt,boxsep=0pt,nobeforeafter,
        left=3pt,right=3pt,top=3pt,bottom=3pt,valign=center]{\qrcode[height=1.9cm]{https://play.google.com/store/apps/details?id=com.luvvix.onbuch&hl=fr}}%
      \hspace{8pt}\begin{minipage}[c]{0.5\linewidth}
        {\popdisplay\fontsize{11}{12.5}\selectfont\color{white}\MakeUppercase{Télécharge OnBuch}}\par\vspace{4pt}
        {\raggedright\popsemi\fontsize{8.4}{11}\selectfont\color{white}Gratuit \textbullet\ Google Play\\Tuteur IA, annales, concours, résultats\par}
      \end{minipage}
    \end{tcolorbox}
  \end{minipage}\hfill
  \begin{minipage}[t]{0.487\linewidth}
    \begin{tcolorbox}[enhanced,colback=poppurple,colframe=popink,boxrule=2.2pt,arc=10pt,
        drop shadow=popink,left=11pt,right=11pt,top=10pt,bottom=10pt,height=3.1cm,valign=center]
      \tikz[baseline=-0.7ex]\node[circle,fill=white,draw=popink,line width=1.3pt,minimum size=1.6cm,inner sep=0]{\popdisplay\fontsize{24}{24}\selectfont\color{poppurple}+};%
      \hspace{8pt}\begin{minipage}[c]{0.6\linewidth}
        {\popdisplay\fontsize{11}{12.5}\selectfont\color{white}\MakeUppercase{Passe à OnBuch+}}\par\vspace{4pt}
        {\raggedright\popsemi\fontsize{8.4}{11}\selectfont\color{white}Tous les cours signés, Tuteur IA illimité, fiches PDF — onglet OnBuch+ de l'app\par}
      \end{minipage}
    \end{tcolorbox}
  \end{minipage}
  \vspace{8pt}
  \popcontenu
  \vfill
  \signatureonbuch
  \restoregeometry
  \newpage
}

% Rangée « Ce cours contient » (page de garde)
\newcommand{\poptile}[3]{% #1 couleur #2 gros texte #3 légende
  \begin{tcolorbox}[enhanced,colback=#1,colframe=popink,boxrule=2pt,arc=9pt,shadow={2.5pt}{-2.5pt}{0pt}{popink},
      left=6pt,right=6pt,top=9pt,bottom=9pt,halign=center,valign=center,height=1.75cm,width=\linewidth,nobeforeafter]
    {\popdisplay\fontsize{12.5}{13}\selectfont\color{white}\MakeUppercase{#2}}\par\vspace{3pt}
    {\centering\popsemi\fontsize{7.8}{9.5}\selectfont\color{white}#3\par}
  \end{tcolorbox}}
\newcommand{\popcontenu}{%
  \par\noindent{\popxboldls\fontsize{7.5}{7.5}\selectfont\color{popmuted}CE COURS SIGNÉ CONTIENT}\par\vspace{7pt}
  \noindent\begin{minipage}[t]{0.235\linewidth}\poptile{popblue}{Cours}{complet, illustré, pas à pas}\end{minipage}\hfill
  \begin{minipage}[t]{0.235\linewidth}\poptile{poporange}{Méthodes}{et exercices résolus}\end{minipage}\hfill
  \begin{minipage}[t]{0.235\linewidth}\poptile{poppink}{Exercices}{type examen + corrigés}\end{minipage}\hfill
  \begin{minipage}[t]{0.235\linewidth}\poptile{popgreen}{Fiche bilan}{l'essentiel à retenir}\end{minipage}\par}

% Sommaire
\newtcolorbox{sommaire}{enhanced,colback=white,colframe=popink,boxrule=2.2pt,arc=10pt,
  drop shadow=popink,left=18pt,right=18pt,top=13pt,bottom=13pt,before skip=6pt,after skip=20pt,
  before upper={\poppill{Sommaire}{poppurple}{white}\par\vspace{9pt}\popsemi\fontsize{10.6}{14.5}\selectfont\color{popink}}}

% Signature de fin de cours — poussée tout en bas de la dernière page
\newcommand{\findecours}{%
  \par\vfill\nopagebreak
  \begin{center}\popsquiggle{poporange}\end{center}\vspace{4pt}
  \signatureonbuch
}
\AtBeginDocument{\def\llbracket{\mathopen{[\mkern-3.5mu[}}\def\rrbracket{\mathclose{]\mkern-3.5mu]}}} % intervalles d'entiers (glyphe ⟦ absent de la police)
% Glyphes absents de Fira Math : redéfinis à partir de glyphes disponibles
\AtBeginDocument{%
  \def\setminus{\mathbin{\backslash}}\let\smallsetminus\setminus
  \def\vdots{\mathord{\vbox{\baselineskip3.2pt\lineskiplimit0pt\kern4pt\hbox{.}\hbox{.}\hbox{.}}}}%
  \def\ddots{\mathinner{\mkern1mu\raise6.5pt\hbox{.}\mkern2mu\raise3.6pt\hbox{.}\mkern2mu\raise0.7pt\hbox{.}\mkern1mu}}%
  \def\triangle{\mathord{\scalebox{0.9}{$\symup{\Delta}$}}}%
  \def\nabla{\mathord{\raisebox{\depth}{\scalebox{1}[-1]{$\symup{\Delta}$}}}}%
  \def\checkmark{\popcheck{popdark}}%
  \def\star{\mathbin{\ast}}\def\bigstar{\mathord{\ast}}%
  \def\diamond{\mathbin{\scalebox{0.6}{$\Diamond$}}}%
  \def\bigcup{\mathop{\mathchoice{\raisebox{-0.3em}{\scalebox{1.7}{$\cup$}}}{\scalebox{1.25}{$\cup$}}{\cup}{\cup}}\displaylimits}%
  \def\bigcap{\mathop{\mathchoice{\raisebox{-0.3em}{\scalebox{1.7}{$\cap$}}}{\scalebox{1.25}{$\cap$}}{\cap}{\cap}}\displaylimits}%
  \def\lesseqqgtr{\mathrel{\lessgtr}}\def\bigoplus{\mathop{\oplus}\displaylimits}%
}
\providecommand{\ssi}{si et seulement si} % abréviation fréquente
\AtBeginDocument{\providecommand{\wideparen}{\overparen}} % arc de cercle

%% FIGFIX-BEGIN v4 — correctifs globaux des figures (docs/onbuch-generation/tools/figfix)
\usepackage{adjustbox}\usepackage{etoolbox}
% toute figure TikZ/pgfplots plus large que la ligne est réduite à la largeur disponible
\BeforeBeginEnvironment{tikzpicture}{\begin{adjustbox}{max width=\linewidth}}
\AfterEndEnvironment{tikzpicture}{\end{adjustbox}}
% légendes pgfplots lisibles par-dessus les courbes
\pgfplotsset{every axis/.append style={legend style={fill=white,fill opacity=0.92,text opacity=1,draw=black!15,inner sep=3pt}}}
% étiquettes d'arêtes (syntaxe edge["..."]) avec fond blanc
\tikzset{every edge quotes/.append style={fill=white,inner sep=1.2pt}}
%% FIGFIX-END
\begin{document}

\pagegarde{Leçon 41 — Expression orale : rédaction d’une lettre formelle}{Chinois}{1ère A}{Module 5 : Mass médias et télécommunication}{ZHA41}{2 h}{Cette leçon d'expression orale apprend aux élèves de 1ère A à présenter oralement la structure et le contenu d'une lettre formelle en chinois, à utiliser les formules fonctionnelles de politesse et à débattre sur la place de la lettre formelle face aux nouveaux médias.
\vspace{4pt}
\begin{multicols}{2}\small
\begin{itemize}
\item À la fin de cette leçon, je sais présenter oralement la structure d'une lettre formelle chinoise (appel, salutations, corps, formule finale, signature)
\item je sais utiliser les formules fonctionnelles pour présenter un sujet : 我要介绍... / 首先... 然后... 最后...
\item je sais donner mon avis dans un débat : 我认为... / 我觉得... / 我同意 / 我不同意
\item je sais relancer un échange : 你觉得呢？/ 你说得对吗？/ 还有呢？
\item je sais prononcer correctement les tons des formules de politesse (zūnjìng, cǐzhì, jìnglǐ)
\item je sais échanger avec mes camarades dans un petit débat guidé sur la lettre formelle et les réseaux sociaux
\end{itemize}\end{multicols}}

\begin{sommaire}
\begin{multicols}{2}
\begin{enumerate}
\item Découverte : la lettre formelle chinoise
\item Formules fonctionnelles : présenter, donner son avis, relancer
\item Dialogues modèles avec pinyin
\item Les tons à travailler
\item Jeux de rôle et débat guidé
\item Activité d'intégration
\item Méthodes et exercices résolus
\item Exercices
\item Corrigés des exercices
\item Fiche bilan
\end{enumerate}
\end{multicols}
\end{sommaire}

\begin{objectifs}
\begin{itemize}
\item À la fin de cette leçon, je sais présenter oralement la structure d'une lettre formelle chinoise (appel, salutations, corps, formule finale, signature)
\item je sais utiliser les formules fonctionnelles pour présenter un sujet : 我要介绍... / 首先... 然后... 最后...
\item je sais donner mon avis dans un débat : 我认为... / 我觉得... / 我同意 / 我不同意
\item je sais relancer un échange : 你觉得呢？/ 你说得对吗？/ 还有呢？
\item je sais prononcer correctement les tons des formules de politesse (zūnjìng, cǐzhì, jìnglǐ)
\item je sais échanger avec mes camarades dans un petit débat guidé sur la lettre formelle et les réseaux sociaux
\item je sais faire un exposé simple de 1 à 2 minutes en chinois sur la rédaction d'une lettre formelle
\item je sais écouter un exposé et poser une question de compréhension
\end{itemize}
\end{objectifs}

\begin{prerequis}
\begin{itemize}
\item Salutations et formules de politesse de base (Seconde) : 你好, 谢谢, 请, 对不起
\item Structures de présentation simples : 我叫..., 我是...
\item Vocabulaire des médias et télécommunications du Module 5 : 信, 邮件, 手机, 网络
\item Prononciation des quatre tons et du ton neutre
\item Connecteurs simples : 因为...所以..., 但是
\end{itemize}
\end{prerequis}

\newpage

\coursec{Découverte : la lettre formelle chinoise}

\courssub{Situation d'accroche et problématique}

Imagine : tu veux demander une bourse d'études en Chine ou écrire au directeur d'une entreprise de télécommunications chinoise implantée à Douala. Au Cameroun, on rédige la lettre en français ou en anglais ; en Chine, il faut connaître les formules de politesse chinoises : \zh{尊敬的...} (Honoré...), \zh{此致敬礼！} (Salutations distinguées !). Aujourd'hui, tu vas apprendre à présenter oralement le contenu d'une lettre formelle et à débattre de son importance à l'ère des réseaux sociaux.

\textbf{Problématique :} comment une lettre formelle chinoise est-elle construite, et quelles formules faut-il connaître pour la présenter oralement devant la classe ?

\courssub{Observation : une lettre formelle modèle}

Avant toute règle, observons. L'enseignant lit à haute voix la lettre suivante ; la classe écoute d'abord sans regarder, puis répète en chœur phrase par phrase. C'est une demande de stage adressée à une entreprise de télécommunications.

\begin{textebox}[Une demande de stage — lettre modèle]
尊敬的王经理：您好！\\
Zūnjìng de Wáng jīnglǐ : Nín hǎo !\\
Honoré Directeur Wang : Bonjour !\\
我叫李明，是雅温得中学的学生。\\
Wǒ jiào Lǐ Míng, shì Yǎwēndé zhōngxué de xuésheng.\\
Je m'appelle Li Ming, je suis élève du lycée de Yaoundé.\\
我想申请贵公司的实习。\\
Wǒ xiǎng shēnqǐng guì gōngsī de shíxí.\\
Je souhaite demander un stage dans votre honorable entreprise.\\
此致敬礼！\\
Cǐzhì jìnglǐ !\\
Salutations distinguées !\\
李明　2024年5月10日\\
Lǐ Míng, 2024 nián 5 yuè 10 rì.\\
Li Ming, le 10 mai 2024.
\end{textebox}

\textbf{Petit glossaire :} \zh{贵公司} (guì gōngsī) = « votre honorable société » (préfixe de respect \zh{贵}) ; \zh{实习} (shíxí) = stage ; \zh{申请} (shēnqǐng) = demander officiellement, poser sa candidature.

\textbf{Questions de compréhension (oral) :}
\begin{enumerate}
\item Qui écrit cette lettre, et à qui ? (Qui est \zh{王经理} ?)
\item Que demande Li Ming ?
\item Où sont placés le nom et la date dans la lettre chinoise ? Compare avec une lettre française.
\end{enumerate}

\courssub{Les cinq parties d'une lettre formelle chinoise}

En observant la lettre modèle, on distingue cinq blocs, toujours dans le même ordre. C'est la structure fixe de toute lettre formelle chinoise, du courrier administratif à l'e-mail professionnel.

\begin{definition}[Les cinq parties de la lettre formelle]
\begin{enumerate}
\item \cle{称呼} (chēnghu) : l'\textbf{appel}, en haut à gauche, suivi des deux-points \zh{：}. Ex. \zh{尊敬的王经理：}.
\item \cle{问候语} (wènhòuyǔ) : les \textbf{salutations} d'ouverture. Ex. \zh{您好！}.
\item \cle{正文} (zhèngwén) : le \textbf{corps} de la lettre : qui on est, ce qu'on demande, pourquoi.
\item \cle{祝颂语} (zhùsòngyǔ) : la \textbf{formule finale} de politesse. Ex. \zh{此致敬礼！}.
\item \cle{署名和日期} (shǔmíng hé rìqī) : la \textbf{signature et la date}, placées en bas \textbf{à droite}.
\end{enumerate}
\end{definition}

\begin{popfigure}
\begin{tikzpicture}[node distance=0.55cm,
  bloc/.style={rectangle, rounded corners, draw=popblue, fill=popblueL, text width=6.2cm, align=center, minimum height=0.9cm, font=\small},
  lbl/.style={font=\small\itshape, text=popdark}]
\node[bloc, align=center] (a){\textbf{1. 称呼} chēnghu\\ 尊敬的王经理：};
\node[bloc, below=of a, fill=poporangeL, draw=poporange, align=center] (b){\textbf{2. 问候语} wènhòuyǔ\\ 您好！};
\node[bloc, below=of b, fill=popgreenL, draw=popgreen, align=center] (c){\textbf{3. 正文} zhèngwén\\ 我叫李明……我想申请……};
\node[bloc, below=of c, fill=poppurpleL, draw=poppurple, align=center] (d){\textbf{4. 祝颂语} zhùsòngyǔ\\ 此致敬礼！};
\node[bloc, below=of d, fill=popgoldL, draw=popgold, align=center] (e){\textbf{5. 署名和日期} shǔmíng hé rìqī\\ 李明　2024年5月10日};
\draw[-{Stealth}, thick, popmuted] (a) -- (b);
\draw[-{Stealth}, thick, popmuted] (b) -- (c);
\draw[-{Stealth}, thick, popmuted] (c) -- (d);
\draw[-{Stealth}, thick, popmuted] (d) -- (e);
\node[lbl, right=0.3cm of a] {appel};
\node[lbl, right=0.3cm of b] {salutations};
\node[lbl, right=0.3cm of c] {corps};
\node[lbl, right=0.3cm of d] {formule finale};
\node[lbl, right=0.3cm of e] {signature + date};
\end{tikzpicture}
\legende{Organigramme vertical des cinq parties de la lettre formelle chinoise}
\end{popfigure}

\begin{popfigure}
\begin{tikzpicture}[mindmap, grow cyclic, every node/.style={concept, font=\small},
  root concept/.append style={concept color=popblue, text=white, font=\small\bfseries},
  level 1/.append style={level distance=3.1cm, sibling angle=72},
  level 1 concept/.append style={concept color=poporange, font=\small}]
\node[root concept, align=center]{正式的信\\ zhèngshì de xìn}
  child { node[align=center]{称呼\\ appel} }
  child { node[align=center]{问候语\\ salutations} }
  child { node[align=center]{正文\\ corps} }
  child { node[align=center]{祝颂语\\ formule finale} }
  child { node[align=center]{署名\\ signature} };
\end{tikzpicture}
\legende{Carte mentale : les composantes de la lettre formelle}
\end{popfigure}

\courssub{Vocabulaire clé de la leçon}

\begin{center}
\begin{tabularx}{\linewidth}{|X|X|X|X|}
\hline
\rowcolor{popblueL} Caractères & Pinyin & Nature & Traduction \\
\hline
信 & xìn & nom & lettre, courrier \\
\hline
正式 & zhèngshì & adjectif & formel, officiel \\
\hline
申请 & shēnqǐng & verbe / nom & demander officiellement ; candidature \\
\hline
尊敬 & zūnjìng & verbe / adjectif & respecter ; honoré, respectable \\
\hline
此致敬礼 & cǐzhì jìnglǐ & formule figée & salutations distinguées \\
\hline
称呼 & chēnghu & nom & appel, formule d'adresse \\
\hline
经理 & jīnglǐ & nom & directeur, gérant \\
\hline
实习 & shíxí & nom / verbe & stage ; faire un stage \\
\hline
贵公司 & guì gōngsī & expression & votre honorable société \\
\hline
署名 & shǔmíng & nom / verbe & signature ; signer \\
\hline
\end{tabularx}
\end{center}

\begin{definition}[此致敬礼 (cǐzhì jìnglǐ)]
Formule finale \textbf{obligatoire} des lettres formelles chinoises. Elle équivaut au français « Veuillez agréer, Monsieur le Directeur, l'expression de mes salutations distinguées ». Littéralement : « par ceci (\zh{此致}), je présente mon respect (\zh{敬礼}) ». Elle ne se traduit jamais mot à mot et ne s'emploie pas à l'oral.
\end{definition}

\courssub{Comparaison avec la lettre formelle française}

\begin{center}
\begin{tabularx}{\linewidth}{|X|X|X|}
\hline
\rowcolor{popblueL} Élément & Français & Chinois \\
\hline
Appel & Monsieur le Directeur, & 尊敬的经理： (Zūnjìng de jīnglǐ :) \\
\hline
Salutation d'ouverture & (aucune en général) & 您好！ (Nín hǎo !) \\
\hline
Présentation de soi & Je me permets de... & 我叫……，是……的学生。 \\
\hline
Formule finale & Veuillez agréer... & 此致敬礼！ \\
\hline
Place de la date & en haut à droite & en bas, sous la signature \\
\hline
\end{tabularx}
\end{center}

Deux différences majeures à retenir : en chinois, on salue explicitement le destinataire avec \zh{您好！} après l'appel ; et la date se place \textbf{en bas}, avec la signature, et non en haut de la page.

\begin{attention}[Erreur fréquente des francophones]
Ne traduis jamais mot à mot « Cher Monsieur » par \zh{亲爱的先生} ! Le mot \zh{亲爱的} (qīn'ài de, « cher/chérie ») est réservé aux lettres \textbf{intimes} (famille, amis proches, amoureux). Dans une lettre officielle, on dit toujours \zh{尊敬的} + fonction ou nom : \zh{尊敬的王经理} (Honoré Directeur Wang), \zh{尊敬的老师} (Honoré Professeur). Dire \zh{亲爱的经理} à un directeur serait aussi choquant que d'écrire « Mon chéri » à un employeur !
\end{attention}

\begin{exoresolu}[Identifier les cinq parties]
Énoncé : voici une mini-lettre. Indique à quelle partie correspond chaque ligne.\\
\zh{尊敬的李老师：您好！我是杜阿拉中学的学生。我想申请去中国的奖学金。此致敬礼！学生小王　2024年6月1日。}
\tcblower
Solution détaillée :
\begin{enumerate}
\item \zh{尊敬的李老师：} → \cle{称呼} (appel : \zh{尊敬的} + nom + fonction).
\item \zh{您好！} → \cle{问候语} (salutation d'ouverture).
\item \zh{我是杜阿拉中学的学生。我想申请去中国的奖学金。} → \cle{正文} (corps : présentation + demande). Traduction : « Je suis élève du lycée de Douala. Je souhaite demander une bourse pour aller en Chine. »
\item \zh{此致敬礼！} → \cle{祝颂语} (formule finale).
\item \zh{学生小王　2024年6月1日。} → \cle{署名和日期} (signature et date, en bas).
\end{enumerate}
La lettre est donc bien construite : les cinq parties sont présentes et dans le bon ordre.
\end{exoresolu}

\begin{savaistu}[此致 à la ligne, 敬礼 en tête de ligne]
En Chine, la mise en page de la formule finale obéit à une règle précise : on écrit \zh{此致} à la fin d'une ligne (souvent avec un léger retrait), puis \zh{敬礼！} \textbf{en tête de la ligne suivante, sans aucun retrait}, car placer le mot « respect » au sommet de la ligne honore le destinataire. Cette règle typographique est encore respectée aujourd'hui, y compris dans les e-mails formels chinois !
\end{savaistu}

\begin{experience}[Écoute et répétition chorale]
\begin{enumerate}
\item \textbf{Écoute} : l'enseignant lit la lettre modèle deux fois ; les élèves ferment les yeux et repèrent les tons.
\item \textbf{Répétition chorale} : la classe répète phrase par phrase, en insistant sur \zh{尊敬的} (zūn-jìng, 1\textsuperscript{er} ton puis 4\textsuperscript{e} ton) et \zh{此致敬礼} (3\textsuperscript{e}-4\textsuperscript{e}-4\textsuperscript{e}-4\textsuperscript{e} tons).
\item \textbf{Lecture en demi-groupes} : un groupe lit la partie \zh{称呼} + \zh{问候语}, l'autre le \zh{正文}, puis tous ensemble le \zh{此致敬礼！}.
\item \textbf{Présentation orale individuelle} : un volontaire présente la lettre à voix haute : « \zh{这是一封正式的信。李明想申请实习。} » (C'est une lettre formelle. Li Ming veut demander un stage.)
\end{enumerate}
\end{experience}

\begin{exercice}[À toi de jouer]
\begin{enumerate}
\item Remets dans l'ordre les cinq parties : \zh{正文} — \zh{称呼} — \zh{署名} — \zh{问候语} — \zh{祝颂语}.
\item Traduis en français : \zh{我想申请贵公司的实习。}
\item Complète l'appel d'une lettre adressée à la directrice Zhang : \zh{\_\_\_\_\_\_ 的张经理：}
\end{enumerate}
\end{exercice}

\begin{corrige}[Exercice « À toi de jouer »]
\begin{enumerate}
\item Ordre correct : \zh{称呼} → \zh{问候语} → \zh{正文} → \zh{祝颂语} → \zh{署名}.
\item « Je souhaite demander un stage dans votre honorable entreprise. » (\zh{贵} marque le respect envers la société.)
\item \zh{尊敬的张经理：} (Zūnjìng de Zhāng jīnglǐ :) — jamais \zh{亲爱的} dans une lettre officielle !
\end{enumerate}
\end{corrige}

\begin{aretenir}[L'essentiel de la découverte]
Une lettre formelle chinoise comporte cinq parties fixes : \zh{称呼} (appel avec \zh{尊敬的}), \zh{问候语} (\zh{您好！}), \zh{正文} (corps), \zh{祝颂语} (\zh{此致敬礼！}) et \zh{署名和日期} (signature et date en bas). La formule \zh{此致敬礼} est obligatoire et \zh{亲爱的} est interdit dans un courrier officiel. Savoir nommer ces parties en chinois est la première étape pour présenter oralement une lettre devant la classe.
\end{aretenir}

\coursec{Leçon 41 — Expression orale : rédaction d'une lettre formelle}

\courssub{Découverte : la lettre formelle chinoise}

Avant de prendre la parole pour présenter ou débattre, il faut connaître l'objet dont on parle. La lettre formelle chinoise \zh{正式信} (\emph{zhèngshì xìn}) suit une structure codifiée, immuable, qui diffère de la lettre administrative française (pas d'en-tête expéditeur en haut à gauche, pas d'objet \og Objet : \fg{} distinct). La maîtriser est la condition \emph{sine qua non} pour en parler correctement à l'oral.

\begin{definition}[Les quatre parties obligatoires]
Une lettre formelle chinoise standard comporte quatre blocs, dans cet ordre strict :
\begin{enumerate}
\item \textbf{L'appel \zh{称呼} (\emph{chēnghu})} : en haut à gauche, suivi de deux-points pleins \zh{：}. Ex. \zh{尊敬的校长：} (\emph{Zūnjìng de xiàozhǎng:}) « Monsieur le Proviseur, ».
\item \textbf{Le corps \zh{正文} (\emph{zhèngwén})} : commence deux lignes plus bas, avec un retrait de deux caractères (espace \zh{空两格} \emph{kòng liǎng gé}). On y expose le motif, les faits, la demande.
\item \textbf{La formule de politesse finale \zh{此致敬礼} (\emph{cǐ zhì jìnglǐ})} : alignée à gauche, sur deux lignes : \zh{此致} (ligne 1) / \zh{敬礼} (ligne 2, retrait de deux caractères). Litt. « En voici l'essentiel / Salutations respectueuses ». Équivalent de \og Je vous prie d'agréer... \fg{}.
\item \textbf{La signature \zh{落款} (\emph{luòkuǎn})} : en bas à droite, sur deux lignes : \zh{写信人：姓名} (ligne 1) / \zh{日期} (ligne 2, année-mois-jour). Ex. \zh{写信人：李华} / \zh{2025年10月15日}.
\end{enumerate}
\end{definition}

\begin{textebox}[Modèle de lettre formelle : demande de stage]
\zh{尊敬的经理：}\\
\zh{\hspace{2em}我是喀麦隆雅温得第一高中的学生。我写信申请贵公司的暑期实习岗位。我认为贵公司的管理经验很丰富，能让我学到很多。}\\
\zh{\hspace{2em}随信附上我的简历。请查收。}\\
\zh{此致}\\
\zh{\hspace{2em}敬礼}\\
\zh{写信人：李华}\\
\zh{2025年10月15日}\\
\\
\emph{Zūnjìng de jīnglǐ:}\\
\emph{Wǒ shì Kāmàilóng Yǎwēndé Dìyī Gāozhōng de xuéshēng. Wǒ xiě xìn shēnqǐng guì gōngsī de shǔqī shíxí gǎngwèi. Wǒ rènwéi guì gōngsī de guǎnlǐ jīngyàn hěn fēngfù, néng ràng wǒ xué dào hěn duō.}\\
\emph{Suí xìn fù shàng wǒ de jiǎnlì. Qǐng cháshōu.}\\
\emph{Cǐ zhì}\\
\emph{Jìnglǐ}\\
\emph{Xiě xìn rén: Lǐ Huá}\\
\emph{2025 nián 10 yuè 15 rì}\\
\\
« Monsieur le Directeur, Je suis élève au Lycée de Yaoundé 1 (Cameroun). J'écris pour solliciter un stage d'été dans votre entreprise. Je pense que votre expérience de gestion est riche et m'apprendra beaucoup. Je joins mon CV à cette lettre. Merci de le réceptionner. Salutations respectueuses. Li Hua. Le 15 octobre 2025. »
\end{textebox}

\begin{attention}[Différences clés Cameroun / France / Chine]
\begin{itemize}
\item \textbf{Français (Cameroun/France) :} En-tête expéditeur (haut gauche), Destinataire (haut droite), Objet, Pièces jointes (bas gauche).
\item \textbf{Chinois :} Pas d'en-tête expéditeur au début. L'appel \zh{称呼} est tout en haut. La signature \zh{落款} (nom + date) est \textbf{en bas à droite}. La formule \zh{此致敬礼} est fixe et obligatoire dans l'administration/éducation.
\item \textbf{Politesse :} \zh{尊敬的} (\emph{zūnjìng de}) « respecté » s'utilise pour les supérieurs, aînés, clients. \zh{亲爱的} (\emph{qīn'ài de}) « cher » est réservé aux proches (informel). Ne jamais confondre.
\end{itemize}
\end{attention}

\begin{center}
\begin{tabularx}{\linewidth}{|X|X|X|X|}
\hline
\rowcolor{popblueL} Caractères & Pinyin & Nature & Traduction \\
\hline
正式信 & zhèngshì xìn & nom & lettre formelle \\
\hline
称呼 & chēnghu & nom & appel, formule d'appel \\
\hline
正文 & zhèngwén & nom & corps du texte \\
\hline
此致敬礼 & cǐ zhì jìnglǐ & locution & salutations respectueuses (formule de clôture) \\
\hline
落款 & luòkuǎn & nom & signature (nom + date) \\
\hline
尊敬的 & zūnjìng de & adj. & respecté (devant titre/nom) \\
\hline
空两格 & kòng liǎng gé & verbe + compl. & laisser deux cases (retrait) \\
\hline
申请 & shēnqǐng & verbe & demander, postuler \\
\hline
实习 & shíxí & nom/verbe & stage / faire un stage \\
\hline
附上 & fù shàng & verbe & joindre (pièce) \\
\hline
查收 & cháshōu & verbe & vérifier et réceptionner \\
\hline
\end{tabularx}
\end{center}

\courssub{Formules fonctionnelles : présenter, donner son avis, relancer}

Dans la section précédente, nous avons découvert la structure de la lettre formelle chinoise : l'appel \zh{称呼} (chēnghu), le corps \zh{正文} (zhèngwén), la salutation finale \zh{此致敬礼} (cǐ zhì jìnglǐ). Mais à l'oral, savoir décrire une lettre ne suffit pas : il faut encore savoir \cle{présenter} son sujet, \cle{organiser} son propos, \cle{donner son avis}, \cle{réagir} aux camarades et \cle{relancer} la discussion. C'est précisément l'objet de cette section : vous donner une boîte à outils de formules toutes faites, utilisables immédiatement dans un exposé ou un débat en classe.

\courssub{Observer d'abord : un extrait d'exposé modèle}

Lisons d'abord comment une élève de Première, Aïcha, présente son exposé sur la lettre formelle. Observons les expressions soulignées par leur répétition : ce sont nos formules cibles.

\begin{textebox}[Exposé d'Aïcha : présenter une lettre formelle]
大家好！\zh{今天我讲的是}一封正式的信。\zh{首先}，信的开头有称呼，比如“尊敬的老师”。\zh{然后}，信的正文要说清楚事情。\zh{最后}，写信的人要写名字和日期。\zh{我认为}写信比发邮件更正式，也更礼貌。我的介绍完了，谢谢大家！\\
\emph{Dàjiā hǎo! Jīntiān wǒ jiǎng de shì yì fēng zhèngshì de xìn. Shǒuxiān, xìn de kāitóu yǒu chēnghu, bǐrú “zūnjìng de lǎoshī”. Ránhòu, xìn de zhèngwén yào shuō qīngchu shìqing. Zuìhòu, xiě xìn de rén yào xiě míngzi hé rìqī. Wǒ rènwéi xiě xìn bǐ fā yóujiàn gèng zhèngshì, yě gèng lǐmào. Wǒ de jièshào wán le, xièxie dàjiā!}\\
« Bonjour à tous ! Aujourd'hui, je parle d'une lettre formelle. D'abord, le début de la lettre contient l'appel, par exemple “Monsieur le Professeur”. Ensuite, le corps de la lettre doit expliquer clairement les choses. Enfin, celui qui écrit doit mettre son nom et la date. Je pense qu'écrire une lettre est plus formel et plus poli qu'envoyer un e-mail. Mon exposé est terminé, merci à tous ! »
\end{textebox}

Que remarque-t-on ? L'exposé suit toujours le même chemin : \textbf{annoncer} le sujet (\zh{今天我讲的是...}), \textbf{ordonner} les idées (\zh{首先... 然后... 最后...}), \textbf{donner un avis} (\zh{我认为...}), \textbf{remercier} (\zh{谢谢大家}). Cette architecture est universelle en chinois scolaire : apprenez-la une fois, elle vous servira pour tous vos exposés.

\courssub{Présenter un sujet : annoncer de quoi on va parler}

\begin{propriete}[Formules pour présenter un sujet]
Pour annoncer le thème d'un exposé, le chinois utilise des structures verbales simples, sans inversion ni subordonnée compliquée :
\begin{itemize}
\item \zh{今天我讲的是 + sujet.} (\emph{Jīntiān wǒ jiǎng de shì...}) — « Aujourd'hui, ce dont je parle, c'est... » Structure : \cle{讲的是} + nom. Le \zh{的} nominalise le verbe \zh{讲} « parler » : littéralement « ce que je présente aujourd'hui, c'est... ».
\item \zh{我要介绍 + sujet.} (\emph{Wǒ yào jièshào...}) — « Je vais présenter... » Structure : sujet + \cle{要} (futur proche / intention) + \zh{介绍} « présenter ».
\item \zh{我要介绍一封正式的信。} (\emph{Wǒ yào jièshào yì fēng zhèngshì de xìn.}) — « Je vais présenter une lettre formelle. » Notez le classificateur \cle{封} (fēng), spécifique aux lettres.
\end{itemize}
\end{propriete}

\begin{attention}[Le classificateur 封 pour les lettres]
En français on dit « une lettre » sans mot-outil. En chinois, impossible de dire \zh{一信} : il faut le classificateur \zh{封} : \zh{一封信} (\emph{yì fēng xìn}) « une lettre », \zh{两封信} (\emph{liǎng fēng xìn}) « deux lettres ». Erreur fréquente des francophones : oublier le classificateur ou utiliser le classificateur général \zh{个}. À l'oral d'examen, \zh{一封信} est la forme attendue.
\end{attention}

\courssub{Organiser l'exposé : 首先... 然后... 最后...}

\begin{propriete}[Les connecteurs d'ordre]
Pour ordonner un propos, le chinois emploie une série de connecteurs fixes, placés en tête de phrase et suivis d'une virgule :
\begin{itemize}
\item \zh{首先} (\emph{shǒuxiān}) — « d'abord, premièrement »
\item \zh{然后} (\emph{ránhòu}) — « ensuite, puis »
\item \zh{最后} (\emph{zuìhòu}) — « enfin, pour finir »
\end{itemize}
Schéma : \textbf{首先，+ phrase 1。然后，+ phrase 2。最后，+ phrase 3。}\\
Contrairement au français qui varie (« d'abord, ensuite, puis, enfin, pour conclure »), le chinois scolaire accepte volontiers la répétition de \zh{然后} pour plusieurs étapes intermédiaires : \zh{首先... 然后... 然后... 最后...} est tout à fait correct.
\end{propriete}

\begin{exemplebox}[Exemples traduits]
\zh{首先，信的开头有称呼。}\\
\emph{Shǒuxiān, xìn de kāitóu yǒu chēnghu.}\\
« D'abord, le début de la lettre contient l'appel. »\\[4pt]
\zh{然后，我要介绍信的正文。}\\
\emph{Ránhòu, wǒ yào jièshào xìn de zhèngwén.}\\
« Ensuite, je vais présenter le corps de la lettre. »\\[4pt]
\zh{最后，写信的人要写日期和名字。}\\
\emph{Zuìhòu, xiě xìn de rén yào xiě rìqī hé míngzi.}\\
« Enfin, celui qui écrit la lettre doit écrire la date et son nom. »
\end{exemplebox}

\courssub{Donner son avis : 认为 ou 觉得 ?}

Trois formules permettent d'exprimer une opinion :

\begin{center}
\begin{tabularx}{\linewidth}{|X|X|X|X|}
\hline
\rowcolor{popblueL} Caractères & Pinyin & Nature & Traduction / nuance \\
\hline
我认为... & wǒ rènwéi & verbe d'opinion & je pense que... (avis réfléchi, argumenté, registre formel) \\
\hline
我觉得... & wǒ juéde & verbe d'opinion & je trouve que... (avis spontané, personnel, registre courant) \\
\hline
在我看来，... & zài wǒ kànlái & locution & à mon avis, selon moi... (introduit une prise de position) \\
\hline
我同意你的看法。 & wǒ tóngyì nǐ de kànfǎ & phrase & je suis d'accord avec ton point de vue. \\
\hline
我不同意，因为... & wǒ bù tóngyì, yīnwèi & phrase & je ne suis pas d'accord, parce que... \\
\hline
\end{tabularx}
\end{center}

\begin{attention}[觉得 ou 认为 : choisir selon le registre]
Les deux verbes se traduisent par « penser / trouver », mais ils ne sont pas interchangeables en contexte formel :
\begin{itemize}
\item \zh{觉得} (\emph{juéde}) exprime une opinion \textbf{personnelle et spontanée}, souvent liée à une impression ou un sentiment : \zh{我觉得今天很热。} (\emph{Wǒ juéde jīntiān hěn rè.}) « Je trouve qu'il fait chaud aujourd'hui. »
\item \zh{认为} (\emph{rènwéi}) exprime une opinion \textbf{réfléchie et argumentée}, celle d'un raisonnement : \zh{我认为写信比发邮件更正式。} (\emph{Wǒ rènwéi xiě xìn bǐ fā yóujiàn gèng zhèngshì.}) « Je pense qu'écrire une lettre est plus formel qu'envoyer un e-mail. »
\end{itemize}
\textbf{Règle pratique pour l'exposé oral :} dans un contexte formel (exposé, débat, lettre officielle), préférez toujours \zh{认为} ou \zh{在我看来}. Réservez \zh{觉得} aux conversations informelles entre amis. Utiliser \zh{我觉得} dans un exposé sur la lettre formelle donnerait une impression de légèreté.
\end{attention}

\begin{exemplebox}[Donner son avis et réagir]
\zh{我认为写信比发邮件更正式。}\\
\emph{Wǒ rènwéi xiě xìn bǐ fā yóujiàn gèng zhèngshì.}\\
« Je pense qu'écrire une lettre est plus formel qu'envoyer un e-mail. »\\[4pt]
\zh{在我看来，正式的信很重要。}\\
\emph{Zài wǒ kànlái, zhèngshì de xìn hěn zhòngyào.}\\
« À mon avis, la lettre formelle est très importante. »\\[4pt]
\zh{我同意你的看法，写信更礼貌。}\\
\emph{Wǒ tóngyì nǐ de kànfǎ, xiě xìn gèng lǐmào.}\\
« Je suis d'accord avec ton point de vue, écrire une lettre est plus poli. »\\[4pt]
\zh{我不同意，因为发邮件更快。}\\
\emph{Wǒ bù tóngyì, yīnwèi fā yóujiàn gèng kuài.}\\
« Je ne suis pas d'accord, parce qu'envoyer un e-mail est plus rapide. »
\end{exemplebox}

\courssub{La comparaison avec 比 : « plus... que »}

Dans l'exemple \zh{写信比发邮件更正式}, apparaît une structure grammaticale essentielle : la comparaison avec \cle{比}.

\begin{propriete}[La structure comparative avec 比]
Pour dire « A est plus... que B », le chinois utilise la structure rigide :\\
\textbf{A + 比 + B + (更) + adjectif}\\
\begin{itemize}
\item \zh{比} (\emph{bǐ}) correspond à « que » de la comparaison française, mais il se place \textbf{avant} B, et l'adjectif vient \textbf{après} B.
\item \zh{更} (\emph{gèng}) « encore plus » renforce l'adjectif ; il est facultatif mais très fréquent à l'oral.
\item Attention : on ne dit jamais \zh{很} dans une comparaison avec \zh{比} (pas de \zh{比... 很正式} : faute grave).
\end{itemize}
Exemples :\\
\zh{写信比发邮件更正式。} \emph{Xiě xìn bǐ fā yóujiàn gèng zhèngshì.} « Écrire une lettre est plus formel qu'envoyer un e-mail. »\\
\zh{雅温得比杜阿拉凉快。} \emph{Yǎwēndé bǐ Dù'ālā liángkuai.} « Yaoundé est plus frais que Douala. »\\
\zh{汉语比英语难一点儿。} \emph{Hànyǔ bǐ Yīngyǔ nán yìdiǎnr.} « Le chinois est un peu plus difficile que l'anglais. »
\end{propriete}

\begin{attention}[Piège de traduction : l'ordre des mots]
Le français dit « la lettre est \textbf{plus formelle que} l'e-mail » : le « que » suit l'adjectif. Le chinois fait l'inverse : \zh{信 \textbf{比} 邮件 \textbf{更正式}} — littéralement « lettre / que / e-mail / plus formelle ». L'adjectif se place \textbf{après} \zh{比 + B}, jamais avant. Erreur typique du francophone : \zh{信更正式比邮件} (calque du français) — à bannir absolument.
\end{attention}

\courssub{Relancer la parole et conclure}

Un bon orateur ne parle pas seul : il fait participer son auditoire, puis conclut proprement.

\begin{propriete}[Relancer la parole et conclure]
\textbf{Pour relancer un camarade} (questions ouvertes, avec le point d'interrogation chinois ？) :
\begin{itemize}
\item \zh{你觉得呢？} (\emph{Nǐ juéde ne?}) — « Et toi, qu'en penses-tu ? » Le \zh{呢} final reprend la question à l'interlocuteur.
\item \zh{你有什么看法？} (\emph{Nǐ yǒu shénme kànfǎ?}) — « Quel est ton point de vue ? »
\item \zh{还有别的问题吗？} (\emph{Hái yǒu bié de wèntí ma?}) — « Y a-t-il encore d'autres questions ? »
\end{itemize}
\textbf{Pour conclure un exposé}, la formule rituelle est :\\
\zh{我的介绍完了，谢谢大家！} (\emph{Wǒ de jièshào wán le, xièxie dàjiā!}) — « Mon exposé est terminé, merci à tous ! »\\
Notez \zh{完了} (\emph{wán le}) « terminé » : \zh{完} est le complément de résultat « finir », suivi du \zh{了} d'accompli.
\end{propriete}

\begin{center}
\begin{tabularx}{\linewidth}{|X|X|X|}
\hline
\rowcolor{popblueL} Fonction & Formule (caractères + pinyin) & Traduction \\
\hline
Présenter & 今天我讲的是... \emph{Jīntiān wǒ jiǎng de shì...} & Aujourd'hui, je parle de... \\
\hline
Présenter & 我要介绍一封信。 \emph{Wǒ yào jièshào yì fēng xìn.} & Je vais présenter une lettre. \\
\hline
Organiser & 首先... 然后... 最后... \emph{shǒuxiān... ránhòu... zuìhòu...} & d'abord... ensuite... enfin... \\
\hline
Donner son avis & 我认为... / 在我看来... \emph{wǒ rènwéi / zài wǒ kànlái} & je pense que... / à mon avis... \\
\hline
Être (pas) d'accord & 我同意。/ 我不同意，因为... \emph{wǒ tóngyì / bù tóngyì, yīnwèi} & je suis d'accord / pas d'accord, car... \\
\hline
Relancer & 你觉得呢？ \emph{Nǐ juéde ne?} & Et toi, qu'en penses-tu ? \\
\hline
Conclure & 我的介绍完了，谢谢大家！ \emph{Wǒ de jièshào wán le, xièxie dàjiā!} & Mon exposé est terminé, merci à tous ! \\
\hline
\end{tabularx}
\end{center}

\begin{popfigure}
\begin{tikzpicture}
\node[draw=popblue, fill=popblueL, rounded corners, align=center, minimum width=4.2cm, minimum height=2.8cm] (pres) at (0,0) {\textbf{Présenter}\\[2pt] 今天我讲的是...\\ \emph{Jīntiān wǒ jiǎng de shì}\\[2pt] 我要介绍一封信。\\ \emph{Wǒ yào jièshào yì fēng xìn}};
\node[draw=poporange, fill=poporangeL, rounded corners, align=center, minimum width=4.2cm, minimum height=2.8cm] (avis) at (5.2,0) {\textbf{Donner son avis}\\[2pt] 我认为... \emph{wǒ rènwéi}\\ 在我看来... \emph{zài wǒ kànlái}\\[2pt] 我不同意，因为...\\ \emph{wǒ bù tóngyì, yīnwèi}};
\node[draw=popgreen, fill=popgreenL, rounded corners, align=center, minimum width=4.2cm, minimum height=2.8cm] (rel) at (10.4,0) {\textbf{Relancer}\\[2pt] 你觉得呢？\\ \emph{Nǐ juéde ne?}\\[2pt] 你有什么看法？\\ \emph{Nǐ yǒu shénme kànfǎ?}};
\draw[-{Stealth}, thick, popmuted] (pres) -- (avis);
\draw[-{Stealth}, thick, popmuted] (avis) -- (rel);
\end{tikzpicture}
\legende{Les trois familles de formules fonctionnelles pour l'exposé et le débat.}
\end{popfigure}

\begin{methode}[Structurer un mini-exposé en cinq étapes]
\begin{enumerate}
\item \textbf{Saluer et annoncer le sujet} : \zh{大家好！今天我讲的是一封正式的信。} « Bonjour ! Aujourd'hui je parle d'une lettre formelle. »
\item \textbf{Développer avec les connecteurs} : \zh{首先... 然后... 最后...} — une idée par étape, des phrases courtes.
\item \textbf{Donner son avis argumenté} : \zh{我认为写信比发邮件更正式，因为...} — utilisez \zh{认为} et la comparaison avec \zh{比}.
\item \textbf{Relancer l'auditoire} : \zh{你觉得呢？} ou \zh{还有别的问题吗？} — montrez que vous écoutez les autres.
\item \textbf{Remercier et conclure} : \zh{我的介绍完了，谢谢大家！} — formule obligatoire, à prononcer avec des tons nets.
\end{enumerate}
\end{methode}

\begin{exoresolu}[Compléter un exposé avec les formules fonctionnelles]
Énoncé : complétez le mini-exposé de Serge avec les formules étudiées : \zh{首先} / \zh{然后} / \zh{最后} / \zh{我认为} / \zh{你觉得呢}.\\
\zh{大家好！今天我讲的是一封正式的信。(1) ，信的开头有称呼。(2) ，我要介绍信的正文。(3) ，写信的人要写名字和日期。(4) 写信比发邮件更正式。(5) ？我的介绍完了，谢谢大家！}
\tcblower
Solution détaillée :\\
(1) \zh{首先} — on ouvre l'énumération des parties de la lettre : « d'abord ».\\
(2) \zh{然后} — deuxième étape, le corps de la lettre : « ensuite ».\\
(3) \zh{最后} — dernière étape, le nom et la date : « enfin ».\\
(4) \zh{我认为} — l'élève donne un avis argumenté avec une comparaison (\zh{比}) : le registre formel impose \zh{认为} et non \zh{觉得}.\\
(5) \zh{你觉得呢} — question de relance adressée à l'auditoire avant la conclusion.\\
Texte complet : \zh{大家好！今天我讲的是一封正式的信。首先，信的开头有称呼。然后，我要介绍信的正文。最后，写信的人要写名字和日期。我认为写信比发邮件更正式。你觉得呢？我的介绍完了，谢谢大家！}\\
\emph{Dàjiā hǎo! Jīntiān wǒ jiǎng de shì yì fēng zhèngshì de xìn. Shǒuxiān, xìn de kāitóu yǒu chēnghu. Ránhòu, wǒ yào jièshào xìn de zhèngwén. Zuìhòu, xiě xìn de rén yào xiě míngzi hé rìqī. Wǒ rènwéi xiě xìn bǐ fā yóujiàn gèng zhèngshì. Nǐ juéde ne? Wǒ de jièshào wán le, xièxie dàjiā!}\\
« Bonjour à tous ! Aujourd'hui je parle d'une lettre formelle. D'abord, le début de la lettre contient l'appel. Ensuite, je vais présenter le corps de la lettre. Enfin, celui qui écrit doit mettre son nom et la date. Je pense qu'écrire une lettre est plus formel qu'envoyer un e-mail. Et toi, qu'en penses-tu ? Mon exposé est terminé, merci à tous ! »
\end{exoresolu}

\begin{savaistu}[La modestie chinoise dans l'expression de l'opinion]
En Chine, dans un cadre formel (exposé, réunion, lettre), on évite souvent d'affirmer « J'ai raison ». On préfère nuancer : \zh{我觉得... 但是...} (\emph{Wǒ juéde... dànshì...}) « Je pense que... mais... » ou \zh{这只是我的个人看法} (\emph{Zhè zhǐ shì wǒ de gèrén kànfǎ}) « Ce n'est que mon avis personnel ». Cela montre l'humilité \zh{谦虚} (\emph{qiānxū}), valeur confucéenne centrale. Au Cameroun, dans un débat scolaire, on attend au contraire une prise de position claire (\zh{我认为}). Adaptez votre registre au contexte : en classe de chinois, \zh{我认为} est la norme attendue.
\end{savaistu}

\begin{exercice}[Entraînement : Mon exposé en 5 phrases]
Rédigez en caractères (avec pinyin) un mini-exposé de 5 phrases sur le thème : \zh{“手机的利与弊”} (Les avantages et inconvénients du smartphone). Utilisez : 1. Salutation + annonce sujet. 2. \zh{首先} (avantage). 3. \zh{然后} (inconvénient). 4. \zh{我认为} + \zh{比} (avis comparatif). 5. Relance + Conclusion.
\end{exercice}

\begin{corrige}[Exercice : Mon exposé en 5 phrases]
\zh{大家好！今天我讲的是手机的利与弊。首先，手机很方便，可以联系朋友。然后，玩手机太久对眼睛不好。我认为学习比玩手机更重要，因为学习决定未来。你觉得呢？我的介绍完了，谢谢大家！}\\
\emph{Dàjiā hǎo! Jīntiān wǒ jiǎng de shì shǒujī de lì yǔ bì. Shǒuxiān, shǒujī hěn fāngbiàn, kěyǐ liánxì péngyou. Ránhòu, wán shǒujī tài jiǔ duì yǎnjing bù hǎo. Wǒ rènwéi xuéxí bǐ wán shǒujī gèng zhòngyào, yīnwèi xuéxí juédìng wèilái. Nǐ juéde ne? Wǒ de jièshào wán le, xièxie dàjiā!}
\end{corrige}

\courssub{Dialogues modèles avec pinyin}

Cette section propose deux dialogues types pour s'entraîner à l'interaction orale : correction d'une lettre (professeur/élève) et débat informel (élèves).

\begin{textebox}[Dialogue 1 : Correction d'une lettre (Professeur -- Élève)]
\zh{老师：} 李华，你的信写得很好。\zh{但是}，称呼那里少了一个冒号。\\
\zh{李华：} \zh{真的吗？} 我以为不用冒号。\\
\zh{老师：} 正式信\zh{必须}加冒号。\zh{而且}，正文要空两格。\\
\zh{李华：} \zh{我明白了。} 那“此致敬礼”写在哪儿？\\
\zh{老师：} “此致”顶格写，“敬礼”空两格。\zh{最后}，落款在右下角。\\
\zh{李华：} \zh{谢谢老师指正。} 我回去改一下。\\
\\
\emph{Lǎoshī: Lǐ Huá, nǐ de xìn xiě de hěn hǎo. Dànshì, chēnghu nàli shǎo le yí ge màohào.}\\
\emph{Lǐ Huá: Zhēn de ma? Wǒ yǐwéi bù yòng màohào.}\\
\emph{Lǎoshī: Zhèngshì xìn bìxū jiā màohào. Érqiě, zhèngwén yào kòng liǎng gé.}\\
\emph{Lǐ Huá: Wǒ míngbái le. Nà “cǐ zhì jìnglǐ” xiě zài nǎr?}\\
\emph{Lǎoshī: “Cǐ zhì” dǐnggé xiě, “jìnglǐ” kòng liǎng gé. Zuìhòu, luòkuǎn zài yòu xiàjiǎo.}\\
\emph{Lǐ Huá: Xièxie lǎoshī zhǐzhèng. Wǒ huíqù gǎi yí xià.}\\
\\
« Professeur : Li Hua, ta lettre est bien écrite. Mais il manque les deux-points à l'appel. Li Hua : Vraiment ? Je croyais qu'il n'en fallait pas. Professeur : La lettre formelle doit avoir les deux-points. De plus, le corps doit avoir un retrait de deux cases. Li Hua : Je comprends. Et "Cizhi Jingli", où l'écrire ? Professeur : "Cizhi" en début de ligne, "Jingli" avec deux cases de retrait. Enfin, la signature en bas à droite. Li Hua : Merci pour la correction, professeur. Je vais corriger en rentrant. »
\end{textebox}

\begin{textebox}[Dialogue 2 : Débat — Lettre vs WeChat (Élèves)]
\zh{小张：} \zh{大家好，今天我讲的是}写信和发微信的区别。\zh{首先}，写信慢，但正式。\zh{然后}，发微信快，但不正式。\zh{我认为}，找工作要写信。\\
\zh{小李：} \zh{我不同意。} \zh{因为}现在很多公司用微信招聘，也很快。\\
\zh{小张：} \zh{你有道理。} \zh{但是}正式录用通知还是发信。\zh{你觉得呢？}\\
\zh{小王：} \zh{在我看来}，视情况而定。给长辈写信更礼貌。\\
\zh{小张：} \zh{好的，谢谢大家！} 我的介绍完了。\\
\\
\emph{Xiǎo Zhāng: Dàjiā hǎo, jīntiān wǒ jiǎng de shì xiě xìn hé fā Wēixìn de qūbié. Shǒuxiān, xiě xìn màn, dàn zhèngshì. Ránhòu, fā Wēixìn kuài, dàn bù zhèngshì. Wǒ rènwéi, zhǎo gōngzuò yào xiě xìn.}\\
\emph{Xiǎo Lǐ: Wǒ bù tóngyì. Yīnwèi xiànzài hěn duō gōngsī yòng Wēixìn zhāopìn, yě hěn kuài.}\\
\emph{Xiǎo Zhāng: Nǐ yǒu dàolǐ. Dànshì zhèngshì lùyòng tōngzhī hái shì fā xìn. Nǐ juéde ne?}\\
\emph{Xiǎo Wáng: Zài wǒ kànlái, shì qíngkuàng ér dìng. Gěi zhǎngbèi xiě xìn gèng lǐmào.}\\
\emph{Xiǎo Zhāng: Hǎo de, xièxie dàjiā! Wǒ de jièshào wán le.}\\
\\
« Xiao Zhang : Bonjour, aujourd'hui je parle de la différence entre écrire une lettre et envoyer un WeChat. D'abord, la lettre est lente mais formelle. Ensuite, WeChat est rapide mais pas formel. Je pense que pour chercher du travail, il faut écrire une lettre. Xiao Li : Je ne suis pas d'accord. Parce que maintenant beaucoup d'entreprises recrutent par WeChat, c'est aussi rapide. Xiao Zhang : Tu as raison. Mais la notification d'embauche formelle est toujours par lettre. Qu'en penses-tu ? Xiao Wang : À mon avis, ça dépend. Écrire une lettre aux aînés est plus poli. Xiao Zhang : D'accord, merci à tous ! Mon exposé est terminé. »
\end{textebox}

\courssub{Les tons à travailler}

La prononciation correcte des tons est cruciale pour l'intelligibilité et la note d'oral. Cette leçon concentre des tons 3, 4 et le ton neutre. Travaillez les paires minimales et les enchaînements ci-dessous.

\begin{propriete}[Points de vigilance tonaux de la leçon]
\begin{itemize}
\item \textbf{Ton 3 (bas-ascendant) :} \zh{写} (\emph{xiě}), \zh{信} (\emph{xìn}), \zh{正} (\emph{zhèng}), \zh{式} (\emph{shì}), \zh{比} (\emph{bǐ}), \zh{认} (\emph{rèn}), \zh{为} (\emph{wéi}), \zh{看} (\emph{kàn}), \zh{法} (\emph{fǎ}), \zh{哪} (\emph{nǎ}), \zh{好} (\emph{hǎo}).
\item \textbf{Ton 4 (descendant) :} \zh{讲} (\emph{jiǎng}), \zh{介} (\emph{jiè}), \zh{绍} (\emph{shào}), \zh{称} (\emph{chēng}), \zh{呼} (\emph{hū}), \zh{落} (\emph{luò}), \zh{款} (\emph{kuǎn}), \zh{申} (\emph{shēn}), \zh{请} (\emph{qǐng}), \zh{附} (\emph{fù}), \zh{收} (\emph{shōu}), \zh{同} (\emph{tóng}), \zh{意} (\emph{yì}), \zh{问} (\emph{wèn}), \zh{题} (\emph{tí}), \zh{完} (\emph{wán}).
\item \textbf{Ton neutre (léger, court) :} \zh{的} (\emph{de}), \zh{了} (\emph{le}), \zh{呢} (\emph{ne}), \zh{吗} (\emph{ma}), \zh{里} (\emph{li} dans \zh{那里}).
\item \textbf{Sandhi du 3e ton :} \zh{写信} (\emph{xié xìn} $\rightarrow$ \emph{xié xìn} : 3+3 $\rightarrow$ 2+3), \zh{认为} (\emph{rén wéi} $\rightarrow$ \emph{rén wéi} : 3+2 pas de sandhi), \zh{比较} (\emph{bǐ jiào} : 3+4 pas de sandhi). Attention : \zh{很好} (\emph{hěn hǎo} $\rightarrow$ \emph{hén hǎo}).
\end{itemize}
\end{propriete}

\begin{experience}[Atelier de prononciation : "Le relais des tons"]
\begin{enumerate}
\item \textbf{Écoute-Modèle :} L'enseignant lit le dialogue 1 phrase par phrase. La classe répète en chœur (chorale), puis individuellement (3 élèves).
\item \textbf{Repérage :} Sur le texte du dialogue 1, surlignez en rouge tous les tons 3, en bleu les tons 4, en vert les tons neutres.
\item \textbf{Relais :} En binôme, A lit la réplique du professeur, B celle de Li Hua. Inversez. Chronométrez : fluidité + justesse des tons.
\item   \textbf{Enregistrement :} Chaque élève s'enregistre (téléphone) en lisant l'exposé complet d'Aïcha (section 2). Auto-évaluation avec la grille : Tons justes / Liaisons / Intonation finale.
\end{enumerate}
\end{experience}

\courssub{Jeux de rôle et débat guidé}

Mise en situation finale : réinvestir tout le lexique, les structures et les formules fonctionnelles dans une tâche complexe.

\begin{experience}[Jeu de rôle : "La lettre de motivation pour le stage"]
\textbf{Contexte :} Vous êtes en Terminale (ou 1ère) au Lycée de Buea / Douala / Yaoundé. Vous postulez pour un stage d'été dans une entreprise chinoise implantée au Cameroun (ex. \zh{华为} Huáwéi, \zh{中建} Zhōngjiàn) ou au consulat de Chine.
\textbf{Rôles :}
\begin{itemize}
\item \textbf{Candidat (A) :} Présente sa lettre à l'oral (2 min) : structure, choix du vocabulaire, argumentaire (\zh{我认为... 比...}).
\item \textbf{Recruteur (B) :} Pose 3 questions précises (\zh{你为什么选我们公司？}, \zh{你的优势是什么？}, \zh{你觉得实习最重要是什么？}). Note les réponses.
\item \textbf{Observateur (C) :} Grille d'évaluation : Tons (5 pts), Structure exposé (5 pts), Formules fonctionnelles utilisées (5 pts), Vocabulaire lettre (5 pts), Interaction/Relance (5 pts). Total /25.
\end{itemize}
\textbf{Consigne :} Rotation des rôles. Chaque élève passe par les 3 postes. Restitution collective : meilleurs extraits entendus.
\end{experience}

\begin{experience}[Débat guidé : "La lettre formelle a-t-elle encore sa place à l'ère numérique ?"]
\textbf{Thèse :} \zh{正式信依然重要} (La lettre formelle reste importante). \textbf{Antithèse :} \zh{电子邮件/微信完全可以取代} (Email/WeChat peuvent totalement remplacer).
\textbf{Déroulement (20 min) :}
\begin{enumerate}
\item \textbf{Préparation (5 min) :} Par groupes de 4. Chaque groupe tire une carte (Pour / Contre). Liste 3 arguments avec \zh{首先/然后/最后} et \zh{因为...所以...}.
\item \textbf{Confrontation (10 min) :} Deux groupes (Pour vs Contre) débattent devant la classe. Règles : 1) On commence par \zh{我认为...} ou \zh{在我看来...}. 2) On répond par \zh{我同意/不同意，因为...}. 3) On relance par \zh{你觉得呢？}.
\item \textbf{Synthèse (5 min) :} L'enseignant note au tableau les arguments clés en caractères. Vote à main levée. Conclusion de l'enseignant en chinois : \zh{总之，正式信在严肃场合不可替代，但电子通讯更高效。} (\emph{Zǒng'ér, zhèngshì xìn zài yánsù chǎnghé bù kě tìdài, dàn diànzǐ tōngxùn gèng xiàolǜ.})
\end{enumerate}
\end{experience}

\begin{aretenir}[L'essentiel de la leçon 41 — Synthèse globale]
\begin{itemize}
\item \textbf{Structure lettre :} \zh{称呼：} (haut gauche) $\rightarrow$ \zh{正文} (retrait 2 cases) $\rightarrow$ \zh{此致/敬礼} (gauche, 2 lignes) $\rightarrow$ \zh{落款} (nom/date, bas droite).
\item \textbf{Vocabulaire clé :} \zh{尊敬的}, \zh{申请}, \zh{实习}, \zh{附上}, \zh{查收}, \zh{空两格}, \zh{落款}.
\item \textbf{Exposé oral :} \zh{今天我讲的是...} / \zh{首先... 然后... 最后...} / \zh{我认为...} (formel) / \zh{我的介绍完了，谢谢大家！}
\item \textbf{Grammaire :} Comparatif \zh{A 比 B 更 + Adj} (pas de \zh{很}). Classificateur \zh{一封信}.
\item \textbf{Interaction :} \zh{你觉得呢？} / \zh{你有什么看法？} / \zh{我同意/不同意，因为...}
\item \textbf{Tons :} Maîtrise des tons 3/4/neutre + sandhi 3+3 sur mots fréquents (\zh{写信}, \zh{很好}, \zh{可以}).
\end{itemize}
\end{aretenir}

\coursec{Leçon 41 — Expression orale : rédaction d'une lettre formelle}

\courssub{Découverte : structure et codes de la lettre formelle chinoise}

\begin{textebox}[Situation de communication : Écrire au proviseur]
Dans un contexte scolaire ou administratif (demande de stage, de bourse, d'autorisation d'absence), la lettre formelle \cle{正式信函} (zhèngshì xìnhán) suit un canevas immuable. Contrairement à la lettre française, l'en-tête ne comporte pas l'adresse de l'expéditeur en haut à gauche, mais seulement le destinataire (aligné à gauche), suivi de la salutation \cle{您好} à la ligne. Le corps du texte utilise des connecteurs logiques (\cle{首先}, \cle{然后}, \cle{最后}). La clôture \cle{此致敬礼} s'écrit sur deux lignes décalées. La signature mentionne le statut (\cle{学生}, \cle{教师}) et la date au format \cle{AAAA年MM月DD日}.
\end{textebox}

\begin{propriete}[Schéma de la lettre formelle chinoise (结构)]
\begin{center}
\begin{tabularx}{\linewidth}{|X|X|X|}
\hline
\rowcolor{popblueL} \textbf{Partie} & \textbf{Contenu chinois} & \textbf{Rôle / Remarques} \\
\hline
En-tête (抬头) & 尊敬的校长：\newline 您好！ & Destinataire + Titre honorifique. \cle{尊敬的} (zūnjìng de) + Fonction. \cle{您好} (nín hǎo) à la ligne, \cle{您} = vous de politesse. \\
\hline
Corps (正文) & 首先，我申请参加夏令营。\newline 然后，我介绍我的情况。\newline 最后，恳请批准。 & Exposition structurée. Connecteurs temporels/logiques : \cle{首先} (shǒuxiān), \cle{然后} (ránhòu), \cle{最后} (zuìhòu). \\
\hline
Formule de fin (落款) & \hspace{2cm} 此致\newline \hspace{4cm} 敬礼 & Formule figée (chengyu). \cle{此致} (cǐzhì) aligné à gauche (ou léger retrait). \cle{敬礼} (jìnglǐ) sur ligne suivante, retrait plus grand (environ 2 caractères). \\
\hline
Signature (署名) & 学生：王明\newline 2024年10月15日 & Statut (学生, 教师) + Nom. Date : Année 年, Mois 月, Jour 日. Pas de lieu. \\
\hline
\end{tabularx}
\end{center}
\end{propriete}

\begin{exemplebox}[Variantes d'ouverture et de clôture selon le destinataire]
\begin{itemize}
    \item \zh{敬爱的老师：} (Jìng'ài de lǎoshī :) — Chère Professeure / Cher Professeur (affectueux, respectueux, pour un enseignant).
    \item \zh{尊敬的经理：} (Zūnjìng de jīnglǐ :) — Monsieur le Directeur / Responsable (entreprise, administration).
    \item \zh{尊敬的招生办主任：} (Zūnjìng de zhāoshēngbàn zhǔrèn :) — Monsieur le Directeur des Admissions (contexte universitaire).
    \item \zh{此致\newline 敬礼} (Cǐzhì \newline Jìnglǐ) — Standard universel (administration, école, entreprise).
    \item \zh{顺颂商祺} (Shùn sòng shāng qí) — Très formel, affaires (Souhaite prospérité commerciale). Sur une ligne.
    \item \zh{祝\newline 工作顺利} (Zhù \newline Gōngzuò shùnlì) — Souhaite bon travail (moins formel que \cle{此致敬礼}, usage courant inter-services).
\end{itemize}
\end{exemplebox}

\coursec{Dialogues modèles avec pinyin}

Après avoir étudié la structure de la lettre, nous passons à sa présentation orale et au débat sur son usage actuel. Cette section propose trois dialogues progressifs : un exposé d'élève, une séance de questions-réponses et un mini-débat. L'objectif est d'observer les formules fonctionnelles en contexte, de travailler la prononciation des tons (notamment le \cle{不} (bù) → \cle{bú} devant ton 4, et les enchaînements tons 3-4) et d'automatiser les marqueurs de politesse et de liaison.

\courssub{Dialogue 1 : Exposé d'un élève — Présentation d'une lettre formelle}

\begin{textebox}[Dialogue 1 — Présentation en classe]
A: 大家好！今天我要介绍一封正式的信。\\
Dàjiā hǎo ! Jīntiān wǒ yào jièshào yì fēng zhèngshì de xìn.\\
Bonjour à tous ! Aujourd'hui je vais présenter une lettre formelle.

B: 请问，这封信写给谁？\\
Qǐngwèn, zhè fēng xìn xiě gěi shéi ?\\
Excusez-moi, à qui est adressée cette lettre ?

A: 写给校长，申请参加夏令营。\\
Xiě gěi xiàozhǎng, shēnqǐng cānjiā xiàlìngyíng.\\
Elle est adressée au proviseur, pour demander à participer au camp d'été.

B: 信的开头怎么写？\\
Xìn de kāitóu zěnme xiě ?\\
Comment écrit-on le début de la lettre ?

A: 首先写«尊敬的校长»，然后换行写«您好»。\\
Shǒuxiān xiě « Zūnjìng de xiàozhǎng », ránhòu huànháng xiě « Nín hǎo ».\\
D'abord on écrit « Respecté Proviseur », puis on va à la ligne et on écrit « Bonjour (politesse) ».

B: 最后写什么？\\
Zuìhòu xiě shénme ?\\
Qu'écrit-on à la fin ?

A: 最后写«此致敬礼», 然后写名字和日期。\\
Zuìhòu xiě « Cǐzhì jìnglǐ », ránhòu xiě míngzi hé rìqī.\\
À la fin on écrit « Salutations distinguées », puis le nom et la date.
\end{textebox}

\begin{aretenir}[Observation linguistique — Dialogue 1]
\begin{itemize}
    \item \textbf{Structure de présentation :} \cle{今天我要介绍…} (Aujourd'hui je vais présenter…) / \cle{首先… 然后… 最后…} (D'abord… Ensuite… Enfin…). Ces connecteurs temporels structurent l'exposé oral.
    \item \textbf{Formule de politesse :} \cle{尊敬的} (zūnjìng de) + Titre/Fonction (校长, 老师, 经理) est l'ouverture standard formelle. Suivi de \cle{您好} (nín hǎo) à la ligne.
    \item \textbf{Formule de clôture :} \cle{此致敬礼} (cǐzhì jìnglǐ) est une formule figée (chengyu) signifiant « Je m'arrête ici et vous présente mes respects ». Elle s'écrit sur deux lignes : \cle{此致} aligné à gauche, \cle{敬礼} en retrait sur la ligne suivante.
    \item \textbf{Classificateur :} \cle{一封信} (yì fēng xìn) — \cle{封} (fēng) est le classificateur pour les lettres, mails, enveloppes.
    \item \textbf{Tons à surveiller :} \cle{介绍} (jièshào, 4-4), \cle{正式} (zhèngshì, 4-4), \cle{夏令营} (xiàlìngyíng, 4-4-2), \cle{换行} (huànháng, 4-2), \cle{此致敬礼} (cǐzhì jìnglǐ, 3-4 4-3). Attention au \cle{不} dans \cle{不是} (bú shì) si utilisé.
\end{itemize}
\end{aretenir}

\courssub{Dialogue 2 : Échange questions/réponses — Interaction après l'exposé}

\begin{textebox}[Dialogue 2 — Questions de la classe]
C: 王明，你的信里提到«夏令营», 是在哪里举办的？\\
Wáng Míng, nǐ de xìn lǐ tí dào « xiàlìngyíng », shì zài nǎlǐ jǔbàn de ?\\
Wang Ming, ta lettre mentionne le « camp d'été », où est-il organisé ?

A: 在北京语言大学，为期两周。\\
Zài Běijīng Yǔyán Dàxué, wéiqī liǎng zhōu.\\
À l'Université de langue et culture de Pékin, durée deux semaines.

D: 听起来很有意思。请问你怎么知道这个机会的？\\
Tīng qǐlái hěn yǒu yìsi. Qǐngwèn nǐ zěnme zhīdào zhège jīhuì de ?\\
Ça a l'air intéressant. Comment as-tu su pour cette opportunité ?

A: 我在学校的微信公众号上看到的通知。\\
Wǒ zài xuéxiào de Wēixìn gōngzhònghào shàng kàn dào de tōngzhī.\\
Je l'ai vu sur le compte officiel WeChat de l'école.

C: 你有什么看法？这次经历对你有什么帮助？\\
Nǐ yǒu shénme kànfǎ ? Zhè cì jīnglì duì nǐ yǒu shénme bāngzhù ?\\
Quelle est ton avis ? En quoi cette expérience t'aidera-t-elle ?

A: 我认为能提高中文水平，还能交中国朋友。\\
Wǒ rènwéi néng tígāo Zhōngwén shuǐpíng, hái néng jiāo Zhōngguó péngyou.\\
Je pense que ça peut améliorer mon niveau de chinois et me permettre de me faire des amis chinois.

D: 说得很好！祝你申请成功。\\
Shuō de hěn hǎo ! Zhù nǐ shēnqǐng chénggōng.\\
Très bien dit ! Je te souhaite succès pour ta candidature.
\end{textebox}

\begin{definition}[Vocabulaire clé — Dialogue 2]
\begin{center}
\begin{tabularx}{\linewidth}{|X|X|X|X|}
\hline
\rowcolor{popblueL} \textbf{Caractères} & \textbf{Pinyin} & \textbf{Nature} & \textbf{Traduction} \\
\hline
提到 & tí dào & V & mentionner, évoquer \\
\hline
举办 & jǔbàn & V & organiser (événement, réunion) \\
\hline
为期 & wéiqī & V & durer (période), avoir une durée de \\
\hline
机会 & jīhuì & N & opportunité, chance \\
\hline
微信公众号 & Wēixìn gōngzhònghào & N & Compte officiel WeChat (page publique) \\
\hline
通知 & tōngzhī & N/V & notification, avis ; notifier \\
\hline
看法 & kànfǎ & N & opinion, point de vue \\
\hline
经历 & jīnglì & N & expérience (vécue) \\
\hline
提高 & tígāo & V & améliorer, augmenter (niveau, qualité) \\
\hline
申请成功 & shēnqǐng chénggōng & VP & avoir sa candidature acceptée / réussite de la candidature \\
\hline
\end{tabularx}
\end{center}
\end{definition}

\begin{propriete}[Structures grammaticales observées — Dialogue 2]
\begin{center}
\begin{tabularx}{\linewidth}{|X|X|X|}
\hline
\rowcolor{poporangeL} \textbf{Structure} & \textbf{Exemple (Caractères + Pinyin)} & \textbf{Traduction / Emploi} \\
\hline
\cle{在…上} (Localisateur + \cle{上}) & \zh{在微信公众号上} (Zài Wēixìn gōngzhònghào shàng) & « Sur le compte WeChat ». \cle{上} (shàng) suit un nom désignant un support, un médium, une surface (journal, site, mur, corps). Indispensable. \\
\hline
\cle{是…的} (Accent sur le détail passé) & \zh{是在北京语言大学举办的} (Shì zài Běijīng Yǔyán Dàxué jǔbàn de) & C'est À Pékin qu'il est organisé. Focus sur le lieu/le mode/le temps d'une action passée connue. \\
\hline
\cle{能/可以 + V} (Capacité/Possibilité) & \zh{能提高中文水平} (Néng tígāo Zhōngwén shuǐpíng) & Peut améliorer le niveau. \cle{能} (néng) = capacité acquise / possibilité objective. \cle{可以} (kěyǐ) = permission / possibilité conditionnelle. \\
\hline
\cle{不但…还/也…} (Progression) & \zh{不但能提高…, 还能交…} (Búdàn néng tígāo…, hái néng jiāo…) & Non seulement… mais encore… Structure de coordination progressive (formel). \cle{不但} (búdàn) + \cle{还} (hái) ou \cle{也} (yě). \\
\hline
\cle{祝 + Sujet + VO} (Souhait) & \zh{祝你申请成功} (Zhù nǐ shēnqǐng chénggōng) & Je te souhaite de réussir ta candidature. Structure figée : \cle{祝} + [bénéficiaire] + [souhait]. Pas de \cle{要}, \cle{会}, \cle{能}. \\
\hline
\end{tabularx}
\end{center}
\end{propriete}

\begin{attention}[Pièges fréquents — Francophones]
\begin{itemize}
    \item \textbf{Erreur :} \zh{*我在微信公众号看见通知。} (Manque \cle{上} après le support virtuel). \textbf{Correct :} \zh{在微信公众号\textbf{上}看见通知。}
    \item \textbf{Erreur :} \zh{*我认为这很好帮助。} (\cle{帮助} bāngzhù est un nom/verbe, pas adjectif). \textbf{Correct :} \zh{我认为这\textbf{很有帮助} / \textbf{对我很有用}。}
    \item \textbf{Ton :} \cle{机会} (jīhuì) — 1er ton + 4e ton. Ne pas dire *jīhuí (2e ton). \cle{经历} (jīnglì) — 1er ton + 4e ton. \cle{提高} (tígāo) — 2e + 1er.
    \item \textbf{Ordre des mots :} \cle{祝你成功} (Souhaite-toi succès) et non *\cle{祝成功你}. Le bénéficiaire suit immédiatement \cle{祝}.
    \item \textbf{Ton \cle{不} (bù) :} \cle{不同意} → \cle{bú tóngyì} (ton 2 + ton 4). \cle{不是} → \cle{bú shì}. \cle{不对} → \cle{bú duì}. \textbf{Règle absolue à l'oral.}
\end{itemize}
\end{attention}

\courssub{Dialogue 3 : Mini-débat — « Lettre formelle ou message WeChat ? »}

\begin{textebox}[Dialogue 3 — Débat guidé]
A: 今天我们讨论一下：申请工作，写正式信好，还是发微信好？\\
Jīntiān wǒmen tǎolùn yíxià: shēnqǐng gōngzuò, xiě zhèngshì xìn hǎo, háishì fā Wēixìn hǎo ?\\
Aujourd'hui nous discutons : pour postuler un emploi, mieux vaut lettre formelle ou WeChat ?

B: 我认为写信比发微信更正式，更有诚意。\\
Wǒ rènwéi xiě xìn bǐ fā Wēixìn gèng zhèngshì, gèng yǒu chéngyì.\\
Je pense qu'écrire une lettre est plus formel, plus sincère.

C: 我不同意。因为现在大家都用手机，微信更方便，回复也快。\\
Wǒ bù tóngyì. Yīnwèi xiànzài dàjiā dōu yòng shǒujī, Wēixìn gèng fāngbiàn, huífù yě kuài.\\
Je ne suis pas d'accord. Parce que maintenant tout le monde utilise le téléphone, WeChat est plus pratique, la réponse est aussi rapide.

A: 但是，正式场合，比如申请国企或公务员，还是要写正式的信。\\
Dànshì, zhèngshì chǎnghé, bǐrú shēnqǐng guóqì huò gōngwùyuán, háishì yào xiě zhèngshì de xìn.\\
Mais, dans un cadre formel, par ex. postuler dans une entreprise d'État ou la fonction publique, il faut quand même écrire une lettre formelle.

B: 对了，而且信可以存档，有法律效力。\\
Duì le, érqiě xìn kěyǐ cúndàng, yǒu fǎlǜ xiàolì.\\
Exact, en plus la lettre peut être archivée, a une valeur juridique.

C: 那平时联系同事，用微信没问题吧？\\
Nà píngshí liánxì tóngshì, yòng Wēixìn méi wèntí ba ?\\
Alors pour contacter ses collègues au quotidien, WeChat pas de problème ?

A: 当然。区分场合最重要。\\
Dāngrán. Qūfēn chǎnghé zuì zhòngyào.\\
Bien sûr. Distinguer le contexte est le plus important.
\end{textebox}

\begin{aretenir}[Formules fonctionnelles repérées — Dialogue 3]
\begin{center}
\begin{tabularx}{\linewidth}{|X|X|X|}
\hline
\rowcolor{popgreenL} \textbf{Fonction} & \textbf{Formule chinoise} & \textbf{Contexte / Nuance} \\
\hline
Lancer le débat & \zh{今天我们讨论一下…} (Jīntiān wǒmen tǎolùn yíxià…) & Ouverture officielle, neutre. \\
\hline
Donner son avis & \zh{我认为…} (Wǒ rènwéi…) / \zh{我觉得…} (Wǒ juéde…) & \cle{认为} : avis réfléchi, formel (exposé, débat). \cle{觉得} : ressenti, subjectif. \\
\hline
Comparer & \zh{A 比 B 更…} (A bǐ B gèng…) & Comparatif de supériorité standard. \\
\hline
Exprimer désaccord & \zh{我不同意，因为…} (Wǒ bù tóngyì, yīnwèi…) & Poli mais ferme. \cle{我不太赞同} (plus doux). \\
\hline
Nuancer / Concéder & \zh{但是… 还是要…} (Dànshì… háishì yào…) & \cle{还是} (háishì) = « malgré tout, quand même » (insistance). \\
\hline
Ajouter un argument & \zh{对了，而且…} (Duì le, érqiě…) & Rebondir. \cle{对了} (duì le) = « Ah oui / D'ailleurs » (oral). \cle{而且} (érqiě) = « de plus / qui plus est ». \\
\hline
Conclure / Résumer & \zh{区分场合最重要} (Qūfēn chǎnghé zuì zhòngyào) & Synthèse finale. \cle{总之} (zǒngzhī) = « En résumé / Bref ». \\
\hline
\end{tabularx}
\end{center}
\end{aretenir}

\begin{savaistu}[Culture numérique : La lettre vs WeChat en Chine et au Cameroun]
En Chine, \cle{微信} (Wēixìn) est omniprésent (plus de 1,3 milliard d'utilisateurs actifs). Il remplace email, SMS, carte de visite, paiement, dossier administratif. Pourtant, pour les \textbf{démarches officielles} (démission \cle{辞职信}, admission universitaire \cle{录取通知书}, candidature fonction publique \cle{公务员考试}, contrats juridiques \cle{合同}), le \textbf{papier timbré} ou le \textbf{PDF signé/cacheté} (\cle{盖章} gāizhāng, cachet rouge officiel) reste la norme légale. La lettre formelle (\cle{正式信函}) n'a pas disparu ; elle a migré vers le format numérique \cle{PDF} mais conserve ses codes stricts (en-tête, \cle{此致敬礼}, cachet). Au Cameroun, l'usage du \cle{WhatsApp} professionnel se développe (groupes d'enseignants, administration), mais la \textbf{lettre administrative manuscrite ou dactylographiée} (avec en-tête du Ministère/Établissement, visa hiérarchique, cachet) reste le standard légal pour les demandes de stage, de congé, de bourse ou de recrutement.
\end{savaistu}

\courssub{Les tons à travailler : focus sur le tone sandhi et l'enchaînement}

\begin{methode}[Mémoriser et jouer un dialogue — 4 étapes]
\begin{enumerate}
    \item \textbf{Lecture silencieuse \& compréhension :} Lire le dialogue chinois + pinyin + traduction. Identifier les \cle{mots-outils} (连词, 助词) et les \cle{formules figées}. Surligner les tons difficiles : \cle{不 bù} → \cle{bú} (ton 2) devant ton 4 (\cle{不同意 bú tóngyì}, \cle{不是 bú shì}, \cle{不对 bú duì}) ; enchaînement ton 3 + ton 3 (\cle{很好 hěn hǎo} → \cle{hén hǎo}) ; ton 3 + ton 4 (\cle{申请 shēnqǐng}, \cle{经历 jīnglì}).
    \item \textbf{Répétition en chœur (跟读 gēndú) :} Professeur lit une réplique → Classe répète en imitant l'intonation, le rythme, la liaison. Insister sur les groupes rythmiques (pauses aux virgules) : \zh{今天我要 / 介绍一封 / 正式的信} \textbar \zh{我认为 / 写信比发微信 / 更正式}.
    \item \textbf{Jeu de rôle AVEC le texte (对练 duìliàn) :} En binômes (A/B ou A/B/C/D). Lire en se regardant, pas en lisant dans le livre. Travailler le contact visuel et les gestes (main ouverte pour \cle{首先}, index pour \cle{我认为}, main coupante pour \cle{但是}).
    \item \textbf{Jeu de rôle SANS le texte (脱稿 tuōgǎo) :} Fermer le livre. Rejouer la scène en changeant un élément (ex: Dialogue 1 → lettre pour stage à Douala ; Dialogue 3 → débat « Lettre de motivation manuscrite vs Email »).
\end{enumerate}
\end{methode}

\begin{experience}[Atelier « Repérage des formules » — 15 min]
\begin{enumerate}
    \item Distribuer le transcript des 3 dialogues \textbf{sans pinyin ni traduction}.
    \item En binômes, surligner : \textcolor{popblue}{\textbf{Bleu}} = Formules de politesse/ouverture/clôture (\cle{尊敬的}, \cle{您好}, \cle{此致敬礼}, \cle{请问}, \cle{祝你}). \textcolor{poporange}{\textbf{Orange}} = Connecteurs logiques (\cle{首先}, \cle{然后}, \cle{最后}, \cle{但是}, \cle{而且}, \cle{总之}, \cle{因为…所以…}). \textcolor{popgreen}{\textbf{Vert}} = Verbes d'action/démarche (\cle{申请}, \cle{介绍}, \cle{讨论}, \cle{区分}, \cle{提交}, \cle{存档}).
    \item Mise en commun au tableau : construire un \textbf{tableau de réemploi} (voir ci-dessous) que les élèves copieront dans leur cahier « Boîte à outils ».
\end{enumerate}
\end{experience}

\begin{exemplebox}[Tableau de réemploi — Ma « Boîte à outils » pour l'oral formel]
\begin{center}
\begin{tabularx}{\linewidth}{|X|X|}
\hline
\rowcolor{poppurpleL} \textbf{Fonction} & \textbf{Formules prêtes à l'emploi (Caractères + Pinyin)} \\
\hline
Ouvrir / Présenter & \zh{大家好，今天我要介绍…} (Dàjiā hǎo, jīntiān wǒ yào jièshào…) \\
\hline
Structurer & \zh{首先… 然后… 最后…} (Shǒuxiān… ránhòu… zuìhòu…) \\
\hline
Demander poliment & \zh{请问…} (Qǐngwèn…) / \zh{不知…} (Bùzhī…) / \cle{可否…} (Kěfǒu…) \\
\hline
Donner avis (formel) & \zh{我认为…} (Wǒ rènwéi…) / \zh{依我看…} (Yī wǒ kàn…) \\
\hline
Comparer & \zh{A 比 B 更/最…} (A bǐ B gèng/zuì…) \\
\hline
Désaccord poli & \zh{我不同意，因为…} (Wǒ bù tóngyì, yīnwèi…) / \zh{我不太赞同。} (Wǒ bú tài zàntóng.) \\
\hline
Nuancer / Relancer & \zh{但是… 还是…} (Dànshì… háishì…) / \zh{对了，而且…} (Duì le, érqiě…) \\
\hline
Conclure & \zh{总之…} (Zǒngzhī…) / \zh{区分场合最重要。} (Qūfēn chǎnghé zuì zhòngyào.) \\
\hline
Souhaiter / Clôturer & \zh{祝你…} (Zhù nǐ…) / \zh{此致敬礼} (Cǐzhì jìnglǐ) \\
\hline
\end{tabularx}
\end{center}
\end{exemplebox}

\courssub{Synthèse visuelle : Schéma de l'échange oral structuré}

\begin{popfigure}
\begin{tikzpicture}[
    node distance=1.2cm and 2.5cm,
    box/.style={draw=popblue, thick, rounded corners, fill=popblueL, text width=2.8cm, align=center, minimum height=1.6cm, font=\small},
    arrow/.style={-Stealth, thick, popdark},
    labelstyle/.style={font=\tiny, text=popmuted, midway, sloped, above}
]
% Nœuds
\node[box, align=center] (start){PRÉSENTER\normalfont\\(Exposé structuré ; 首选 然后 最后)};
\node[box, right=of start, align=center] (q){QUESTIONNER\normalfont\\(Clarifier ; 请问… 怎么… 谁)};
\node[box, right=of q, align=center] (opinion){DONNER AVIS\normalfont\\(Argumenter ; 我认为 比 较)};
\node[box, right=of opinion, align=center] (react){RÉAGIR/RELANCER\normalfont\\(Nuancer ; 但是 而且 原来)};
\node[box, below=of q, align=center] (conclude){CONCLURE\normalfont\\(Synthétiser ; 总之 区分场合)};

% Flèches principales
\draw[arrow] (start) -- node[labelstyle] {Écoute active} (q);
\draw[arrow] (q) -- node[labelstyle] {Réponse} (opinion);
\draw[arrow] (opinion) -- node[labelstyle] {Échange} (react);
\draw[arrow] (react) -- node[labelstyle, right] {Rebond} (conclude);
\draw[arrow] (conclude) -- node[labelstyle, below] {Clôture} (start.west);

% Flèches retour (interaction dynamique)
\draw[arrow, poporange, dashed] (q) to[bend left=20] node[labelstyle] {Relance} (start);
\draw[arrow, poporange, dashed] (opinion) to[bend left=20] node[labelstyle] {Contredire} (q);
\end{tikzpicture}
\legende{Schéma cyclique d'un échange oral structuré : Présenter → Questionner → Argumenter → Relancer → Conclure. Les flèches oranges pointillées montrent la nature dynamique (aller-retour) du débat.}
\end{popfigure}

\courssub{Exercice résolu : Transformer un message informel en échange formel}

\begin{exoresolu}[Réécriture : De WeChat à l'oral formel]
\textbf{Consigne :} Ton camarade t'envoie ce message WeChat : \zh{« 嘿，听说你要去北京夏令营？牛啊！怎么申请的？ »} (Héi, tīngshuō nǐ yào qù Běijīng xiàlìngyíng? Niú a! Zěnme shēnqǐng de?). Transforme cet échange en une \textbf{question orale formelle} (comme au Dialogue 2) et une \textbf{réponse structurée}.

\tcblower
\textbf{Solution détaillée :}

\textbf{1. Analyse du registre :} Message original = oral, familier (\cle{嘿} hēi « hé », \cle{牛啊} niú a « c'est génial/top » argotique, phrase nominale \cle{怎么申请的} sans sujet, ton direct).

\textbf{2. Version formelle (Question - Élève C) :}
\begin{quote}
\zh{王同学，请问我听说您计划参加北京夏令营，不知申请流程是怎样的？}\\
Wáng tóngxué, qǐngwèn wǒ tīngshuō nín jìhuà cānjiā Běijīng xiàlìngyíng, bùzhī shēnqǐng liúchéng shì zěnyàng de ?\\
Camarade Wang, permettez-moi de demander : j'ai entendu dire que vous prévoyez de participer au camp de Pékin, pourriez-vous m'indiquer la procédure de candidature ?
\end{quote}
\begin{itemize}
    \item \cle{王同学} (Wáng tóngxué) : Appellation polie camarade (Nom + Tongxue) remplace \cle{嘿}.
    \item \cle{请问… 不知…} (Qǐngwèn… bùzhī…) : Formules d'atténuation (hédging) remplacent l'interrogation directe brutale.
    \item \cle{您} (nín) : Vous de politesse au lieu de \cle{你} (nǐ).
    \item \cle{申请流程} (shēnqǐng liúchéng) : Vocabulaire précis (procédure de candidature) au lieu de \cle{怎么申请} (comment postuler).
    \item \cle{是怎样的} (shì zěnyàng de) : Structure complète avec \cle{是…的} pour focus sur la manière.
\end{itemize}

\textbf{3. Version formelle (Réponse - Élève A) :}
\begin{quote}
\zh{谢谢关心。我是在学校微信公众号上看到通知的。首先填写表格，然后提交成绩单，最后面试。我认为这个机会很宝贵。}\\
Xièxie guānxīn. Wǒ shì zài xuéxiào Wēixìn gōngzhònghào shàng kàn dào tōngzhī de. Shǒuxiān tiánxiě biǎogé, ránhòu tíjiāo chéngjīdān, zuìhòu miànshì. Wǒ rènwéi zhège jīhuì hěn bǎoguì.\\
Merci de votre intérêt. Je l'ai vu sur le compte officiel WeChat de l'école. D'abord remplir le formulaire, ensuite soumettre le relevé de notes, enfin passer l'entretien. Je pense que cette opportunité est précieuse.
\end{quote}
\begin{itemize}
    \item \cle{谢谢关心} (Xièxie guānxīn) : Remerciement standard pour l'intérêt porté.
    \item Structure \cle{首先… 然后… 最后…} pour l'explication claire et séquentielle.
    \item \cle{提交} (tíjiāo, soumettre/déposer officiellement) > \cle{交} (jiāo, donner/remettre - neutre).
    \item \cle{宝贵} (bǎoguì, précieux/précieux) > \cle{牛} (niú, génial - argot).
    \item \cle{我认为…} (Wǒ rènwéi…) pour l'avis personnel formel.
\end{itemize}
\end{exoresolu}

\coursec{Jeux de rôle et débat guidé}

\begin{exercice}[Entraînement : Créer votre mini-débat]
En binôme, choisissez un sujet ci-dessous. Rédigez un mini-dialogue (6-8 répliques) en utilisant \textbf{au moins 5 formules fonctionnelles} du tableau « Boîte à outils » (ex: \cle{我认为}, \cle{我不同意}, \cle{但是}, \cle{而且}, \cle{总之}, \cle{请问}, \cle{祝你}). Écrivez en caractères + pinyin. Préparez-vous à le jouer sans notes (étape \cle{脱稿}).
\begin{enumerate}
    \item Sujet : « Pour un stage en entreprise : Lettre de motivation manuscrite vs Email avec PDF. »
    \item Sujet : « Inviter un conférencier chinois : Courrier officiel vs Message WeChat. »
    \item Sujet : « Répondre à une convocation d'examen : SMS de confirmation vs Lettre de réception signée. »
\end{enumerate}
\end{exercice}

\begin{corrige}[Exercice — Mini-débat (Exemple Sujet 1 : Stage)]
\textbf{A:} 今天我们讨论：找实习，手写信好还是发邮件好？\\
Jīntiān wǒmen tǎolùn: zhǎo shíxí, shǒuxiě xìn hǎo háishì fā yóujiàn hǎo ?\\
Aujourd'hui on discute : pour un stage, lettre manuscrite ou email ?

\textbf{B:} 我认为发邮件更方便，而且HR看得快。\\
Wǒ rènwéi fā yóujiàn gèng fāngbiàn, érqiě HR kàn de kuài.\\
Je pense que l'email est plus pratique, et en plus les RH le lisent vite.

\textbf{A:} 我不同意。手写信更有诚意，能看到字迹和态度。\\
Wǒ bù tóngyì. Shǒuxiě xìn gèng yǒu chéngyì, néng kàn dào zìjì hé tàidu.\\
Je ne suis pas d'accord. La lettre manuscrite est plus sincère, on voit l'écriture et l'attitude.

\textbf{B:} 但是现在很多公司环保，不收纸质信。\\
Dànshì xiànzài hěnduō gōngsī huánbǎo, bù shōu zhǐzhì xìn.\\
Mais maintenant beaucoup d'entreprises sont écolo, ne prennent pas les lettres papier.

\textbf{A:} 对了，而且邮件可以附件发成绩单。\\
Duì le, érqiě yóujiàn kěyǐ fùjiàn fā chéngjīdān.\\
Exact, et en plus l'email peut envoyer les relevés en pièce jointe.

\textbf{B:} 也是。总之，看公司要求，区分场合最重要。\\
Yě shì. Zǒngzhī, kàn gōngsī yāoqiú, qūfēn chǎnghé zuì zhòngyào.\\
C'est vrai aussi. Bref, selon les exigences de l'entreprise, distinguer le contexte est le plus important.
\end{corrige}

\begin{attention}[Bilan des erreurs à ne plus commettre à l'oral (Check-list examen)]
\begin{itemize}
    \item \textbf{Ton \cle{不} (bù) :} \textbf{Toujours} \cle{bú} (ton 2) devant un ton 4 (\cle{不同意 bú tóngyì}, \cle{不是 bú shì}, \cle{不对 bú duì}, \cle{不要 bú yào}). \textbf{Point crucial pour l'examen oral.}
    \item \textbf{\cle{和} vs \cle{跟} vs \cle{与} :} Oral courant : \cle{和} (hé) ou \cle{跟} (gēn). \cle{与} (yǔ) = écrit/formel/littéraire. Ne pas mélanger dans la même phrase (ex: *\cle{我和你与他}).
    \item \textbf{« Je pense que » :} \cle{我认为} (rènwéi : formel, exposé/débat, opinion raisonnée) > \cle{我觉得} (juéde : subjectif, sentiment, impression) > \cle{我想} (xiǎng : intention/volonté/plan). Choisir selon le contexte.
    \item \textbf{Vitesse vs Clarté :} Mieux vaut parler \textbf{moins vite} avec des tons justes et des pauses aux virgules (，) qu'aller vite et « avaler » les tons. L'intelligibilité prime sur la fluidité apparente.
    \item \textbf{Classificateur \cle{封} (fēng) :} Une lettre = \cle{一封信}. Ne pas dire *\cle{一张信} (zhāng = feuille plate) ni *\cle{一个信}.
    \item \textbf{\cle{此致敬礼} :} Ne pas oublier le retour à la ligne et le retrait pour \cle{敬礼}. C'est un critère de correction à l'écrit comme à l'oral (description de la mise en page).
\end{itemize}
\end{attention}

\coursec{Les tons à travailler}

Dans la section précédente, nous avons lu et répété les dialogues modèles de la lettre formelle : demande de stage, remerciement, relance. Vous avez remarqué que ces dialogues contiennent des formules très solennelles : \zh{尊敬的}, \zh{此致敬礼}, \zh{贵公司}. Ces mots ne pardonnent aucune faute de ton : mal prononcés, ils perdent leur politesse, voire changent de sens. Cette section est donc consacrée au \cle{travail phonétique} : nous allons décortiquer les tons des formules clés, découvrir la règle de transformation du troisième ton, et apprendre à poser une question avec la bonne mélodie.

\courssub{Observer d'abord : les tons des formules clés}

Écoutons (ou lisons à voix haute) les quatre formules essentielles de la lettre formelle, en suivant la mélodie de chaque syllabe :

\begin{exemplebox}[Les quatre formules à mémoriser avec leurs tons]
\zh{尊敬} \emph{zūnjìng} (1\ier{} ton -- 4\ieme{} ton) : « honorable, respecté ». La première syllabe reste haute et plate, la deuxième tombe brusquement.\\
\zh{此致敬礼} \emph{cǐzhì jìnglǐ} (3\ieme{} -- 4\ieme{} -- 4\ieme{} -- 3\ieme{} ton) : « avec mes salutations respectueuses » (formule de politesse finale).\\
\zh{申请} \emph{shēnqǐng} (1\ier{} -- 3\ieme{} ton) : « demander, solliciter ».\\
\zh{同意} \emph{tóngyì} (2\ieme{} -- 4\ieme{} ton) : « être d'accord, approuver ».
\end{exemplebox}

Observez : aucune de ces formules n'est « plate ». Le chinois est une \cle{langue à tons} : chaque syllabe porte une mélodie qui fait partie du mot, exactement comme une voyelle ou une consonne. En français, la mélodie sert à exprimer l'étonnement ou la question ; en chinois, elle sert à \emph{distinguer les mots}. Dire \zh{tóngyì} avec deux tons plats, c'est comme écrire « tongui » en français : le mot n'existe plus.

\begin{center}
\begin{tabularx}{\linewidth}{|X|X|X|X|}
\hline
\rowcolor{popblueL} Caractères & Pinyin & Nature / Tons & Traduction \\
\hline
尊敬 & zūnjìng & Adjectif (1--4) & honorable, respecté \\
\hline
此致敬礼 & cǐzhì jìnglǐ & Locution (3--4--4--3) & salutations respectueuses \\
\hline
申请 & shēnqǐng & Verbe (1--3) & demander, solliciter \\
\hline
同意 & tóngyì & Verbe (2--4) & approuver, accepter \\
\hline
贵公司 & guì gōngsī & Nom (4--1) & votre honorable société \\
\hline
实习 & shíxí & Verbe / Nom (2--2) & stage ; faire un stage \\
\hline
\end{tabularx}
\end{center}

\begin{exemplebox}[Deux exemples en contexte]
\zh{我想在贵公司实习。} \emph{Wǒ xiǎng zài guì gōngsī shíxí.} « Je voudrais faire un stage dans votre honorable société. »\\
\zh{老师同意了我的申请。} \emph{Lǎoshī tóngyì le wǒ de shēnqǐng.} « Le professeur a accepté ma demande. »
\end{exemplebox}

\courssub{Les quatre tons en images}

Avant d'aller plus loin, visualisons les quatre tons (plus le ton neutre) sur quatre syllabes de notre leçon : \zh{尊} \emph{zūn}, \zh{同} \emph{tóng}, \zh{礼} \emph{lǐ} et la particule \zh{吗} \emph{ma}.

\begin{popfigure}
\begin{tikzpicture}[scale=1]
% Ton 1
\begin{scope}[shift={(0,0)}]
\draw[popline] (0,0) rectangle (2.4,1.8);
\draw[popline, dashed] (0,1.35) -- (2.4,1.35);
\draw[popline, dashed] (0,0.45) -- (2.4,0.45);
\draw[very thick, popblue] (0.25,1.35) -- (2.15,1.35);
\node[align=center] at (1.2,-0.5) {\zh{尊} zūn\\ \footnotesize 1\ier{} ton : haut et plat};
\end{scope}
% Ton 2
\begin{scope}[shift={(3.2,0)}]
\draw[popline] (0,0) rectangle (2.4,1.8);
\draw[popline, dashed] (0,1.35) -- (2.4,1.35);
\draw[popline, dashed] (0,0.45) -- (2.4,0.45);
\draw[very thick, popgreen] (0.25,0.55) -- (2.15,1.35);
\node[align=center] at (1.2,-0.5) {\zh{同} tóng\\ \footnotesize 2\ieme{} ton : montant};
\end{scope}
% Ton 3
\begin{scope}[shift={(6.4,0)}]
\draw[popline] (0,0) rectangle (2.4,1.8);
\draw[popline, dashed] (0,1.35) -- (2.4,1.35);
\draw[popline, dashed] (0,0.45) -- (2.4,0.45);
\draw[very thick, poporange] (0.25,0.95) .. controls (0.9,0.15) and (1.5,0.15) .. (2.15,0.95);
\node[align=center] at (1.2,-0.5) {\zh{礼} lǐ\\ \footnotesize 3\ieme{} ton : descend puis monte};
\end{scope}
% Ton 4
\begin{scope}[shift={(9.6,0)}]
\draw[popline] (0,0) rectangle (2.4,1.8);
\draw[popline, dashed] (0,1.35) -- (2.4,1.35);
\draw[popline, dashed] (0,0.45) -- (2.4,0.45);
\draw[very thick, poppink] (0.25,1.35) -- (2.15,0.35);
\node[align=center] at (1.2,-0.5) {\zh{敬} jìng\\ \footnotesize 4\ieme{} ton : tombant};
\end{scope}
% Ton neutre
\begin{scope}[shift={(12.8,0)}]
\draw[popline] (0,0) rectangle (2.4,1.8);
\draw[popline, dashed] (0,1.35) -- (2.4,1.35);
\draw[popline, dashed] (0,0.45) -- (2.4,0.45);
\fill[popmuted] (1.2,0.9) circle (0.07);
\node[align=center] at (1.2,-0.5) {\zh{吗} ma\\ \footnotesize ton neutre : court et léger};
\end{scope}
\end{tikzpicture}
\legende{Les quatre tons et le ton neutre : haut-plat (zūn), montant (tóng), descendant-montant (lǐ), tombant (jìng), neutre (ma).}
\end{popfigure}

\begin{aretenir}[Les cinq mélodies du chinois]
\begin{itemize}
\item \cle{1\ier{} ton} (ā) : voix haute et stable, comme quand le médecin vous demande de dire « aaaa ».
\item \cle{2\ieme{} ton} (á) : voix qui monte, comme le « hein ? » français de la surprise.
\item \cle{3\ieme{} ton} (ǎ) : voix qui descend puis remonte ; dans la conversation rapide, on n'entend souvent que la partie basse.
\item \cle{4\ieme{} ton} (à) : voix qui tombe sec, comme un « non ! » décidé.
\item \cle{Ton neutre} (a) : syllabe courte, légère, sans mélodie propre.
\end{itemize}
\end{aretenir}

\courssub{La règle d'or : deux troisièmes tons qui se suivent}

Observons deux mots que vous connaissez déjà : \zh{你好} et \zh{敬礼}. Sur le papier, le pinyin s'écrit \emph{nǐ hǎo} et \emph{jìng lǐ} : deux troisièmes tons à chaque fois. Mais à l'oral, personne ne prononce deux creux successifs : ce serait lent et fatigant. Écoutez un locuteur natif dire « bonjour » : vous entendez \emph{ní hǎo}, avec une première syllabe qui monte.

\begin{propriete}[Sandhi du troisième ton]
Quand \textbf{deux syllabes au troisième ton se suivent}, la \textbf{première se prononce comme un deuxième ton} (montant). La deuxième garde son troisième ton.
\begin{itemize}
\item \zh{你好} : écrit \emph{nǐ hǎo}, prononcé \emph{ní hǎo} « bonjour ».
\item \zh{敬礼} : écrit \emph{jìng lǐ}, prononcé \emph{jíng lǐ} « salutation respectueuse ».
\item \zh{很好} : écrit \emph{hěn hǎo}, prononcé \emph{hén hǎo} « très bien ».
\item \zh{可以} : écrit \emph{kěyǐ}, prononcé \emph{kéyǐ} « pouvoir, d'accord ».
\end{itemize}
\textbf{Attention :} la transformation est uniquement \cle{orale}. On continue d'écrire le pinyin avec les tons d'origine : \emph{nǐ hǎo}, jamais \emph{ní hǎo} à l'écrit. Les caractères, eux, ne changent jamais.
\end{propriete}

Cas particulier de notre formule finale : dans \zh{此致敬礼} \emph{cǐzhì jìnglǐ}, seul le dernier mot \zh{敬礼} subit le sandhi (\emph{jìng lǐ} devient \emph{jíng lǐ}), car \zh{此} \emph{cǐ} est suivi de \zh{致} \emph{zhì}, un quatrième ton : pas de transformation. La mélodie réelle est donc : 3 -- 4 -- 2 -- 3.

\begin{attention}[Ce que le francophone a tendance à oublier]
Deux erreurs symétriques guettent l'apprenant camerounais :
\begin{enumerate}
\item Prononcer \emph{nǐ hǎo} avec deux vrais troisièmes tons (deux creux) : cela sonne robotique et étrange.
\item Inversement, \textbf{écrire} \emph{ní hǎo} dans une copie : à l'écrit, le pinyin garde toujours les tons d'origine. Le sandhi est une règle de prononciation, pas d'orthographe.
\end{enumerate}
Retenez : \textbf{on écrit l'original, on dit le transformé}.
\end{attention}

\courssub{L'intonation des questions avec 吗 et 呢}

Pour relancer un correspondant ou demander une réponse, nos dialogues utilisent deux particules interrogatives :

\begin{exemplebox}[Deux types de questions]
\zh{您同意我的申请吗？} \emph{Nín tóngyì wǒ de shēnqǐng ma?} « Acceptez-vous ma demande ? » (question fermée, réponse oui/non).\\
\zh{我的信呢？} \emph{Wǒ de xìn ne?} « Et ma lettre ? » (question de relance : « où en est... ? »).
\end{exemplebox}

\begin{propriete}[Mélodie des questions à particule]
\begin{itemize}
\item \zh{吗} \emph{ma} se place à la fin d'une phrase affirmative et la transforme en question oui/non. \zh{吗} est toujours au \textbf{ton neutre} : court, léger, sans mélodie. La voix monte légèrement en fin de phrase, mais c'est l'intonation générale, pas un ton sur la particule.
\item \zh{呢} \emph{ne} sert à relancer (« et... ? », « où en est... ? »). Elle est aussi au ton neutre, avec une légère montée de la voix.
\item Structure : Phrase affirmative + \zh{吗} ? / Nom ou pronom + \zh{呢} ?
\end{itemize}
Comparaison avec le français : le français forme la question oui/non surtout par l'intonation montante (« Vous acceptez ? ») ou l'inversion. Le chinois, lui, garde l'ordre des mots de la phrase affirmative et ajoute simplement \zh{吗} : \textbf{Phrase + 吗 ?}. C'est plus simple : pas d'inversion, pas de « est-ce que ».
\end{propriete}

\begin{attention}[Le piège du ton plein sur 吗]
L'erreur la plus fréquente des francophones est de prononcer \zh{吗} avec un ton plein, par exemple un premier ton \emph{mā}, parce que la syllabe est écrite « ma ». Résultat : la fin de la phrase devient lourde et le mot ressemble à \zh{妈} \emph{mā} « maman » ! Le ton neutre doit rester \textbf{court, léger, presque soufflé}. Imaginez que \zh{吗} est une petite bulle accrochée à la fin de la phrase : elle flotte, elle ne pèse pas. Même vigilance pour \zh{呢} \emph{ne} et \zh{的} \emph{de}.
\end{attention}

\courssub{Le quatrième ton énergique : 贵 dans 贵公司}

La lettre formelle exige de désigner poliment l'entreprise du destinataire : on dit \zh{贵公司} \emph{guì gōngsī} (4\ieme{} -- 1\ier{} ton), littéralement « honorable société », c'est-à-dire « votre société ». Le mot \zh{贵} \emph{guì} porte un quatrième ton qui doit \textbf{tomber net et fort}, comme un coup de tampon officiel sur un document. Un quatrième ton mou ou hésitant donne une impression de manque de sérieux, exactement l'inverse de l'effet recherché dans une lettre formelle.

\begin{exemplebox}[贵 en contexte]
\zh{贵公司很有名。} \emph{Guì gōngsī hěn yǒumíng.} « Votre honorable société est très connue. »\\
\zh{请问，贵公司在哪里？} \emph{Qǐngwèn, guì gōngsī zài nǎli?} « Excusez-moi, où se trouve votre société ? »
\end{exemplebox}

\begin{methode}[Travailler les tons avec la main]
Pour ancrer les tons dans la mémoire du corps, suivez ces étapes à chaque répétition :
\begin{enumerate}
\item Levez la main à hauteur d'épaule : c'est le niveau haut de votre voix.
\item Premier ton : la main avance à l'horizontale, sans bouger.
\item Deuxième ton : la main monte en diagonale, comme pour poser une question.
\item Troisième ton : la main descend puis remonte en creux.
\item Quatrième ton : la main tombe d'un coup sec, comme pour couper.
\item Répétez chaque formule (\zh{尊敬}, \zh{申请}, \zh{同意}, \zh{贵公司}) trois fois en frappant le rythme, puis sans la main, à voix haute.
\end{enumerate}
En classe, travaillez en chœur : le professeur montre la courbe de la main, la classe prononce ; puis l'inverse, le professeur prononce et les élèves dessinent la courbe dans les airs.
\end{methode}

\courssub{Paires minimales : entraînement sur mā / má / mǎ / ma}

Pour vérifier que vos tons sont bien distincts, rien de tel que les \cle{paires minimales} : des syllabes identiques qui ne diffèrent que par le ton.

\begin{center}
\begin{tabularx}{\linewidth}{|X|X|X|X|}
\hline
\rowcolor{popblueL} Caractère & Pinyin & Nature / Ton & Traduction \\
\hline
妈 & mā & Particule / 1\ier{} ton & maman \\
\hline
麻 & má & Nom / 2\ieme{} ton & chanvre ; engourdi \\
\hline
马 & mǎ & Nom / 3\ieme{} ton & cheval \\
\hline
骂 & mà & Verbe / 4\ieme{} ton & insulter, gronder \\
\hline
吗 & ma & Particule / ton neutre & particule de question \\
\hline
\end{tabularx}
\end{center}

Exercice de discrimination : le professeur prononce l'une de ces cinq syllabes ; les élèves lèvent le nombre de doigts correspondant au ton (1, 2, 3, 4, ou un poing fermé pour le neutre). Appliquez ensuite le même exercice aux syllabes de la leçon : \emph{zūn / zún / zǔn / zùn}, puis \emph{jīng / jíng / jǐng / jìng}.

\begin{exoresolu}[Identifier et corriger les tons]
\textbf{Énoncé.} Voici la fin d'une lettre formelle prononcée par un élève : « cǐzhì jìng lǐ » avec deux troisièmes tons pleins, puis « Nín tóngyì mā ? » avec un premier ton sur la particule. (1) Indiquez les deux erreurs de prononciation. (2) Donnez la prononciation correcte. (3) Écrivez le pinyin correct tel qu'il doit apparaître à l'écrit.
\tcblower
\textbf{Solution détaillée.}
(1) Première erreur : dans \zh{敬礼}, deux troisièmes tons se suivent ; l'élève n'a pas appliqué le sandhi. Deuxième erreur : la particule \zh{吗} a été prononcée avec un premier ton plein \emph{mā} au lieu du ton neutre.
(2) Prononciation correcte : \emph{cǐzhì jíng lǐ} (le premier 3\ieme{} ton de \zh{敬礼} devient un 2\ieme{} ton) ; \emph{Nín tóngyì ma ?} avec un \emph{ma} court et léger, et une légère montée d'intonation en fin de phrase.
(3) À l'écrit, le pinyin garde les tons d'origine : \emph{cǐzhì jìnglǐ} et \emph{Nín tóngyì ma?} — on n'écrit jamais \emph{jíng lǐ}, car le sandhi est une règle orale uniquement.
\end{exoresolu}

\begin{exercice}[À vous de chanter les tons]
\begin{enumerate}
\item Classez ces mots de la leçon selon leur schéma de tons : \zh{尊敬}, \zh{申请}, \zh{同意}, \zh{实习}, \zh{贵公司}.
\item Prononcez à voix haute en appliquant le sandhi : \zh{你好}, \zh{敬礼}, \zh{很好}, \zh{可以}. Écrivez ensuite leur pinyin correct (tons d'origine).
\item Transformez en questions avec \zh{吗} : \zh{您同意。} / \zh{他申请了。} Puis relancez avec \zh{呢} : \zh{我的信} / \zh{我的申请}.
\item En binôme : l'un prononce \emph{mā, má, mǎ, mà} ou \emph{ma} au hasard, l'autre identifie le ton et donne la traduction.
\end{enumerate}
\end{exercice}

\begin{corrige}[Exercice : À vous de chanter les tons]
\begin{enumerate}
\item 1--4 : \zh{尊敬} ; 1--3 : \zh{申请} ; 2--4 : \zh{同意} ; 2--2 : \zh{实习} ; 4--1 : \zh{贵公司}.
\item Prononciation : \emph{ní hǎo}, \emph{jíng lǐ}, \emph{hén hǎo}, \emph{kéyǐ}. Écriture correcte : \emph{nǐ hǎo}, \emph{jìnglǐ}, \emph{hěn hǎo}, \emph{kěyǐ}.
\item \zh{您同意吗？} \emph{Nín tóngyì ma?} ; \zh{他申请了吗？} \emph{Tā shēnqǐng le ma?} ; \zh{我的信呢？} \emph{Wǒ de xìn ne?} ; \zh{我的申请呢？} \emph{Wǒ de shēnqǐng ne?}
\item mā = maman ; má = chanvre/engourdi ; mǎ = cheval ; mà = gronder ; ma = particule de question.
\end{enumerate}
\end{corrige}

\begin{aretenir}[L'essentiel phonétique de la leçon]
\begin{itemize}
\item Les formules de la lettre formelle ont des tons fixes à respecter : \zh{尊敬} (1--4), \zh{申请} (1--3), \zh{同意} (2--4), \zh{贵公司} (4--1), \zh{此致敬礼} (3--4--4--3).
\item Deux troisièmes tons consécutifs : le premier devient un deuxième ton \textbf{à l'oral seulement} ; l'écriture du pinyin ne change pas.
\item \zh{吗} et \zh{呢} restent au ton neutre : courts et légers ; la question se marque par la particule, pas par l'inversion.
\item Le quatrième ton de \zh{贵} doit tomber net : c'est la marque sonore du respect dans la correspondance formelle.
\end{itemize}
Vous êtes maintenant prêts pour la dernière étape : les jeux de rôle et le débat guidé, où vous mettrez ces tons en musique dans de vraies conversations.
\end{aretenir}

\coursec{Jeux de rôle et débat guidé}

Dans la section précédente, nous avons travaillé les tons à l'oral : le premier ton haut et plat de \zh{公司} (gōngsī), le quatrième ton descendant de \zh{意见} (yìjiàn), le ton léger de \zh{吗} (ma). Il est maintenant temps de mettre tout cela en action. Une lettre formelle ne vit pas seulement sur le papier : elle se présente, se défend et se discute à voix haute. C'est exactement ce que nous allons faire dans cette dernière étape : jouer, débattre et parler chinois comme dans la vraie vie professionnelle.

\courssub{Jeu de rôle 1 : la demande de stage face au directeur}

\textbf{Mise en situation.}\par

La classe se divise en binômes. L'élève A joue le rôle du \cle{directeur d'entreprise} (par exemple le directeur d'une société de télécommunications à Douala). L'élève B joue le rôle de l'élève qui a rédigé sa lettre formelle et qui vient \cle{présenter oralement sa demande de stage}.

\begin{experience}[Jeu de rôle 1 — 面试实习生 Miànshì shíxíshēng « Entretien avec le stagiaire »]
Déroulement conseillé :
\begin{enumerate}
\item L'élève B frappe à la porte, salue poliment : \zh{您好！我是……} (Nín hǎo! Wǒ shì…) — « Bonjour ! Je suis… »
\item Il présente sa demande en trois phrases : qui il est, ce qu'il demande, pourquoi.
\item Le directeur (élève A) pose au moins deux questions : \zh{你为什么想来我们公司？} (Nǐ wèishénme xiǎng lái wǒmen gōngsī?) — « Pourquoi voulez-vous venir dans notre entreprise ? »
\item L'élève B remercie et prend congé : \zh{谢谢您的时间，再见！} (Xièxie nín de shíjiān, zàijiàn!) — « Merci pour votre temps, au revoir ! »
\item On échange les rôles.
\end{enumerate}
\end{experience}

\begin{exemplebox}[Phrases utiles pour l'entretien et le débat]
\zh{我叫王华，是第一中学的学生。} (Wǒ jiào Wáng Huá, shì dì-yī zhōngxué de xuésheng.) — « Je m'appelle Wang Hua, je suis élève du lycée n\textsuperscript{o} 1. »\\[2pt]
\zh{我想在贵公司实习。} (Wǒ xiǎng zài guì gōngsī shíxí.) — « Je voudrais faire un stage dans votre honorable entreprise. »\\[2pt]
\zh{我准备好了简历。} (Wǒ zhǔnbèi hǎo le jiǎnlì.) — « J'ai préparé mon CV. »\\[2pt]
\zh{我同意李明的看法。} (Wǒ tóngyì Lǐ Míng de kànfǎ.) — « Je suis d'accord avec l'avis de Li Ming. »\\[2pt]
\zh{还有别的意见吗？} (Hái yǒu bié de yìjiàn ma?) — « Y a-t-il d'autres avis ? »
\end{exemplebox}

\begin{attention}[Le mot de politesse 贵 guì]
Dans une lettre ou un entretien formel, on ne dit jamais \zh{你的公司} (nǐ de gōngsī, « ton entreprise ») au directeur : on dit \zh{贵公司} (guì gōngsī), littéralement « l'honorable entreprise ». C'est l'équivalent chinois de « votre honorable établissement » en français administratif. Oublier \zh{贵} n'est pas une faute de grammaire, mais c'est une faute de politesse !
\end{attention}

\courssub{Jeu de rôle 2 : l'interview du journaliste}

\textbf{Mise en situation.}\par

Un élève joue le \cle{journaliste} d'un journal lycéen (par exemple le journal du lycée de Yaoundé), l'autre joue un \cle{expert} qui vient d'écrire une lettre formelle. Le journaliste pose des questions sur la lettre : à qui, pourquoi, comment.

\begin{experience}[Jeu de rôle 2 — 记者采访 Jìzhě cǎifǎng « Interview de presse »]
Questions modèles du journaliste :
\begin{itemize}
\item \zh{你的信是写给谁的？} (Nǐ de xìn shì xiě gěi shéi de?) — « À qui ta lettre est-elle adressée ? »
\item \zh{你为什么写这封信？} (Nǐ wèishénme xiě zhè fēng xìn?) — « Pourquoi as-tu écrit cette lettre ? »
\item \zh{写信难不难？} (Xiě xìn nán bu nán?) — « Écrire une lettre, est-ce difficile ? »
\end{itemize}
Réponses modèles de l'expert :
\begin{itemize}
\item \zh{我的信是写给公司经理的。} (Wǒ de xìn shì xiě gěi gōngsī jīnglǐ de.) — « Ma lettre est adressée au directeur de l'entreprise. »
\item \zh{因为我想找一份实习工作。} (Yīnwèi wǒ xiǎng zhǎo yí fèn shíxí gōngzuò.) — « Parce que je voudrais trouver un stage. »
\item \zh{不难，只要掌握格式。} (Bù nán, zhǐyào zhǎngwò géshì.) — « Pas difficile, tant qu'on maîtrise le format. »
\end{itemize}
\end{experience}

\courssub{Débat guidé : la lettre formelle est-elle encore utile ?}

\textbf{Question du débat.}\par

\begin{center}
\zh{在WhatsApp和电子邮件的时代，正式信件还有用吗？}\\
(Zài WhatsApp hé diànzǐ yóujiàn de shídài, zhèngshì xìnjiàn hái yǒu yòng ma?)\\
« À l'ère de WhatsApp et des e-mails, la lettre formelle est-elle encore utile ? »
\end{center}

\textbf{Répartition des rôles.}\par

\begin{itemize}
\item \cle{L'animateur} (\zh{主持人} zhǔchírén) : il pose la question, donne la parole et relance avec \zh{你觉得呢？} (Nǐ juéde ne?) — « Et toi, qu'en penses-tu ? »
\item \cle{Le camp « pour »} (\zh{同意方} tóngyì fāng) : deux élèves, chacun avec un argument.
\item \cle{Le camp « contre »} (\zh{不同意方} bù tóngyì fāng) : deux élèves, chacun avec un argument.
\item \cle{Le secrétaire de séance} (\zh{记录员} jìlùyuán) : il note au tableau les formules utilisées.
\end{itemize}

\begin{methode}[Animer un débat en cinq étapes]
\begin{enumerate}
\item \textbf{Saluer} : \zh{大家好！} (Dàjiā hǎo!) — « Bonjour à tous ! »
\item \textbf{Poser la question} : \zh{今天的题目是……} (Jīntiān de tímù shì…) — « Le sujet d'aujourd'hui est… »
\item \textbf{Donner la parole} : \zh{请你先说。} (Qǐng nǐ xiān shuō.) — « Je te prie de parler d'abord. »
\item \textbf{Relancer} : \zh{你觉得呢？还有别的意见吗？} (Nǐ juéde ne? Hái yǒu bié de yìjiàn ma?) — « Qu'en penses-tu ? Y a-t-il d'autres avis ? »
\item \textbf{Conclure et remercier} : \zh{谢谢大家！} (Xièxie dàjiā!) — « Merci à tous ! »
\end{enumerate}
\end{methode}

\begin{popfigure}
\begin{tikzpicture}[
  node distance=1.0cm,
  box/.style={draw=popblue, fill=popblueL, rounded corners=3pt, align=center, minimum width=3.5cm, minimum height=0.9cm, font=\small},
  pour/.style={draw=popgreen, fill=popgreenL, rounded corners=3pt, align=center, minimum width=3.5cm, font=\small},
  contre/.style={draw=poporange, fill=poporangeL, rounded corners=3pt, align=center, minimum width=3.5cm, font=\small},
  vote/.style={draw=poppurple, fill=poppurpleL, rounded corners=3pt, align=center, minimum width=3.8cm, font=\small},
  fl/.style={-{Stealth[length=2.5mm]}, thick, popdark}
]
\node[box, align=center] (anim){Animateur 主持人\\pose la question};
\node[pour, below left=1.2cm and 0.5cm of anim, align=center] (pour){Pour 同意方\\Arg. 1 : demandes officielles\\Arg. 2 : examens et emplois};
\node[contre, below right=1.2cm and 0.5cm of anim, align=center] (contre){Contre 不同意方\\Arg. 1 : WhatsApp plus rapide\\Arg. 2 : e-mail moins cher};
\node[vote, below=3.8cm of anim] (vote) {Vote de la classe 全班投票};
\draw[fl] (anim) -- (pour);
\draw[fl] (anim) -- (contre);
\draw[fl] (pour) -- (vote);
\draw[fl] (contre) -- (vote);
\end{tikzpicture}
\legende{Organigramme du débat : l'animateur lance, les deux camps argumentent, la classe vote.}
\end{popfigure}

\begin{attention}[Contredire poliment, jamais brutalement]
Dans un débat formel chinois, on contredit \textbf{poliment} : \zh{我不同意，因为……} (Wǒ bù tóngyì, yīnwèi…) — « Je ne suis pas d'accord, parce que… ». On ne dit \textbf{jamais} \zh{你错了！} (Nǐ cuò le!) — « Tu as tort ! », beaucoup trop direct et blessant. En français aussi, on préfère « je ne partage pas ton avis » à « tu as tort » : la culture chinoise, qui accorde une grande importance à la « face » (\zh{面子} miànzi), est encore plus sensible à cette politesse.
\end{attention}

\courssub{Grille d'écoute et bilan collectif}

Pendant le débat, chaque élève qui ne parle pas remplit une \cle{grille d'écoute} : il note \textbf{une formule entendue} et \textbf{une question à poser} à la fin.

\begin{center}
\begin{tabularx}{\linewidth}{|X|X|X|X|}
\hline
\rowcolor{popblueL} Caractères & Pinyin & Nature & Traduction \\
\hline
主持人 & zhǔchírén & nom & animateur (du débat) \\
\hline
记录员 & jìlùyuán & nom & secrétaire de séance \\
\hline
意见 & yìjiàn & nom & avis, opinion \\
\hline
看法 & kànfǎ & nom & point de vue \\
\hline
同意 & tóngyì & verbe & être d'accord \\
\hline
有用 & yǒuyòng & adjectif & utile \\
\hline
投票 & tóupiào & verbe & voter \\
\hline
时代 & shídài & nom & ère, époque \\
\hline
\end{tabularx}
\end{center}

\begin{exoresolu}[Grille d'écoute — exemple rempli]
Énoncé : pendant le débat, note une formule entendue et rédige une question à poser au camp « contre ».
\tcblower
Solution détaillée :
\begin{itemize}
\item Formule entendue : \zh{我不同意，因为电子邮件更快。} (Wǒ bù tóngyì, yīnwèi diànzǐ yóujiàn gèng kuài.) — « Je ne suis pas d'accord, parce que l'e-mail est plus rapide. » (Structure : \zh{我} + \zh{不同意} + \zh{因为} + raison.)
\item Question à poser : \zh{如果没有网络，怎么办？} (Rúguǒ méiyǒu wǎngluò, zěnme bàn?) — « Et s'il n'y a pas d'Internet, que fait-on ? » (Structure : \zh{如果} + hypothèse + ，\zh{怎么办}?)
\end{itemize}
Méthode : pour former la question, on utilise \zh{如果} (si) + situation, puis \zh{怎么办} (que faire ?). Attention à ne pas traduire mot à mot le français « que se passe-t-il si… » : le chinois place l'hypothèse d'abord, la question ensuite.
\end{exoresolu}

\textbf{Bilan collectif.}\par

À la fin du débat, la classe procède au \cle{vote} : \zh{我们投票吧！} (Wǒmen tóupiào ba!) — « Votons ! ». Puis le secrétaire de séance lit sa liste, et le professeur écrit au tableau \textbf{les meilleures formules utilisées} : celles qui étaient correctes, bien prononcées (tons justes !) et polies. Chacun les recopie dans son cahier : elles deviennent le répertoire personnel de la classe pour les futurs exposés.

\begin{savaistu}[La lettre formelle : un atout réel]
Au Cameroun comme en Chine, les demandes officielles — bourse, emploi, stage, inscription universitaire — exigent toujours un courrier formel, même à l'ère de WhatsApp. En Chine, la lettre de motivation (\zh{求职信} qiúzhíxìn) suit des règles de politesse très strictes ; au Cameroun, aucune demande de stage n'aboutit sans lettre manuscrite ou tapée. Maîtriser ce genre en chinois, c'est donc doubler vos chances : réussir vos examens de LV2 \emph{et} vous démarquer dans la vie professionnelle, où la coopération sino-camerounaise multiplie les occasions d'écrire en chinois.
\end{savaistu}

\begin{exercice}[À vous de jouer]
\begin{enumerate}
\item En binôme, rejouez l'entretien de stage en changeant d'entreprise (banque, hôpital, station de radio).
\item Rédigez deux arguments « pour » et deux arguments « contre » avec la structure \zh{我认为……因为……} (Wǒ rènwéi… yīnwèi…) — « Je pense que… parce que… ».
\item Animez à tour de rôle un mini-débat de trois minutes en utilisant au moins trois formules du tableau de vocabulaire.
\end{enumerate}
\end{exercice}

\newpage

\coursec{Activité d'intégration}

\coursec{Leçon 41 — Expression orale : rédaction d’une lettre formelle}

\courssub{Découverte : la lettre formelle chinoise}

\begin{definition}[La lettre formelle (正式信 \zh{zhèngshì xìn})]
Dans le monde sinophone (Chine, Taïwan, Singapour) comme dans les relations professionnelles sino-camerounaises, la lettre formelle conserve une place irremplaçable pour les démarches officielles : candidature (stage, emploi, bourse), réclamation, correspondance administrative ou commerciale. Contrairement au message instantané (WeChat, WhatsApp), elle engage la responsabilité de l'expéditeur, témoigne du respect de la hiérarchie (\cle{guānxì 关系}) et suit un \cle{code rituel strict} (appel honorifique, formules figées, présentation soignée). Au Cameroun, les grandes entreprises (CamTel, MTN, Orange, SONARA) et les ambassades exigent souvent une version chinoise pour les dossiers impliquant des partenaires de Chine.
\end{definition}

\begin{propriete}[Structure type d'une lettre formelle]
\begin{center}
\begin{tikzpicture}[
    node distance=0.8cm and 1.2cm,
    box/.style={draw=popblue, thick, rounded corners, fill=popblueL, text width=2.8cm, align=center, font=\small},
    arrow/.style={-Stealth, popblue, thick}
]
\node[box, align=center] (appel){\textbf{1. Appel}\\\zh{尊敬的 + Titre + :}\\(ex. \zh{尊敬的王经理：})};
\node[box, below=of appel, align=center] (presentation){\textbf{2. Présentation}\\\zh{姓名 + Identité + 单位}\\(ex. \zh{我叫李雅文，是…学生})};
\node[box, below=of presentation, align=center] (motif){\textbf{3. Motif / Demande}\\\zh{申请 / 感兴趣 + 理由}\\(ex. \zh{想申请贵公司暑期实习})};
\node[box, below=of motif, align=center] (qualites){\textbf{4. Atouts / Qualités}\\\zh{技能 + 经验 + 语言}\\(ex. \zh{学习中文两年了})};
\node[box, below=of qualites, align=center] (politesse){\textbf{5. Formules finales}\\\zh{谢谢您的时间 / 期待回复}\\\zh{此致，敬礼！}};
\node[box, below=of politesse, align=center] (signature){\textbf{6. Signature / Date}\\\zh{学生：姓名}\\\zh{二〇二五年五月十日}};
\draw[arrow] (appel) -- (presentation);
\draw[arrow] (presentation) -- (motif);
\draw[arrow] (motif) -- (qualites);
\draw[arrow] (qualites) -- (politesse);
\draw[arrow] (politesse) -- (signature);
\end{tikzpicture}
\end{center}
\emph{Ordonnez toujours les informations du plus formel (destinataire) au plus personnel (expéditeur).}
\end{propriete}

\begin{savaistu}[Cameroun / Chine : codes épistolaires]
\begin{itemize}
\item \textbf{Chine :} L'en-tête place le destinataire en haut à gauche (\zh{尊敬的……：}). La formule \zh{此致，敬礼！} s'écoule sur deux lignes : \zh{此致} aligné à gauche, \zh{敬礼！} indenté de deux caractères sur la ligne suivante. L'enveloppe rouge est de rigueur pour les occasions solennelles.
\item \textbf{Cameroun :} L'usage du français impose « Monsieur le Directeur, » ou « Madame, ». La signature précède souvent la date à droite.
\item \textbf{Point commun :} L'usage du \textbf{vous honorifique} (\zh{您} / \emph{Vous}) et du vocabulaire valorisant le destinataire (\zh{贵公司} « votre honorable entreprise », \zh{贵校} « votre honorable établissement ») est obligatoire dans les deux cultures.
\end{itemize}
\end{savaistu}

\courssub{Objectifs de l'activité}

À la fin de cette activité d'intégration, tu seras capable de :
\begin{itemize}
\item présenter oralement, en chinois, une lettre formelle simple (demande de stage, de bourse ou d'emploi) en 1 à 2 minutes ;
\item structurer ton exposé avec les connecteurs \zh{首先} (d'abord), \zh{然后} (ensuite), \zh{最后} (enfin) ;
\item utiliser au moins deux formules de politesse adaptées à la lettre formelle (\zh{尊敬的……}, \zh{谢谢您的时间}, \zh{期待您的回复}) ;
\item poser des questions avec \zh{请问……} et donner ton avis avec \zh{我认为……} ou \zh{我不同意，因为……} ;
\item participer à un débat guidé en respectant la prononciation des tons et le tour de parole.
\end{itemize}

\courssub{Situation}

Ton lycée de Yaoundé organise chaque année le « Forum des jeunes reporters ». Cette année, le thème choisi est : \emph{la lettre formelle à l'ère des réseaux sociaux}. La grande entreprise de télécommunications camerounaise « CamTel » vient d'annoncer des stages d'été pour les lycéens intéressés par les métiers des médias et de la communication. Pour postuler, il faut rédiger une lettre formelle en chinois — l'entreprise travaille en effet avec des partenaires chinois (Huawei, ZTE, China Telecom).

Ton professeur de chinois te propose de préparer ce forum : chaque groupe de quatre élèves rédige une lettre formelle, la présente oralement devant la classe, puis toute la classe débat de la question : \emph{Faut-il encore apprendre à écrire des lettres formelles quand tout le monde utilise WhatsApp et les réseaux sociaux ?}

Voici le modèle de lettre distribué par le professeur :

\begin{textebox}[Modèle de lettre formelle : demande de stage]
尊敬的王经理：您好！\\
Zūnjìng de Wáng jīnglǐ : Nín hǎo !\\
« Monsieur le Directeur Wang, bonjour ! »\\[4pt]
我叫李雅文，是雅温得第二中学高一年级的学生。\\
Wǒ jiào Lǐ Yǎwén, shì Yǎwēndé Dì-èr Zhōngxué gāo yī niánjí de xuésheng.\\
« Je m'appelle Li Yawen, je suis élève de Première (Seconde année secondaire) au Lycée n°2 de Yaoundé. »\\[4pt]
我对通信技术很感兴趣，所以想申请贵公司的暑期实习。\\
Wǒ duì tōngxìn jìshù hěn gǎn xìngqù, suǒyǐ xiǎng shēnqǐng guì gōngsī de shǔqī shíxí.\\
« Je m'intéresse beaucoup aux technologies de communication, c'est pourquoi je souhaite demander un stage d'été dans votre entreprise. »\\[4pt]
我学习中文已经两年了，英语和法语说得也不错。\\
Wǒ xuéxí Zhōngwén yǐjīng liǎng nián le, Yīngyǔ hé Fǎyǔ shuō de yě búcuò.\\
« J'apprends le chinois depuis deux ans, et je parle aussi assez bien anglais et français. »\\[4pt]
谢谢您的时间，期待您的回复。\\
Xièxie nín de shíjiān, qīdài nín de huífù.\\
« Merci pour votre temps, j'attends votre réponse avec impatience. »\\[4pt]
此致，\\
敬礼！\\
Cǐzhì,\\
Jìnglǐ !\\
« Veuillez agréer mes salutations distinguées ! »\\[4pt]
学生：李雅文　二〇二五年五月十日\\
Xuésheng : Lǐ Yǎwén　Èr líng èr wǔ nián wǔ yuè shí rì\\
« L'élève : Li Yawen, le 10 mai 2025 »
\end{textebox}

\begin{attention}[Pièges de la lettre formelle — Erreurs fréquentes des francophones]
\begin{itemize}
\item \textbf{\zh{您} vs \zh{你}} : On écrit \zh{您} (vous, poli, avec le radical cœur \zh{心}) et non \zh{你} dans une lettre formelle. Oublier le cœur change le registre en familier.
\item \textbf{\zh{贵公司} (guì gōngsī)} : Le préfixe honorifique \zh{贵} est indispensable devant le nom de l'entreprise du destinataire ; on ne dit jamais \zh{你的公司} (votre entreprise, neutre/impoli) ni \zh{我们公司} (notre entreprise) pour désigner celle du destinataire.
\item \textbf{Formule finale \zh{此致，敬礼！}} : Elle est \textbf{figée}. \zh{此致} = « Ici s'arrête (la lettre) », \zh{敬礼} = « Salutations respectueuses ». Ne la traduis pas mot à mot, apprends-la par cœur. Mise en page : \zh{此致} en début de ligne, \zh{敬礼！} sur la ligne suivante, indenté de deux caractères.
\item \textbf{Classificateur des lettres} : Une lettre = \zh{一封信} (yī fēng xìn). Jamais \zh{一个信} ni \zh{一张信}.
\item \textbf{Attention aux tons (sandhi et homophones)} :
      \begin{itemize}
      \item \zh{shēnqǐng} 申请 (demander, postuler, 1\textsuperscript{er} + 3\textsuperscript{e} ton) $\neq$ \zh{shēngqì} 生气 (se fâcher) ni \zh{shēngqǐng} 请 (inviter, autre caractère).
      \item \zh{jīnglǐ} 经理 (directeur, gérant) : 1\textsuperscript{er} ton + 3\textsuperscript{e} ton. \textbf{Sandhi} : deux 3\textsuperscript{e} tons consécutifs $\rightarrow$ le premier devient 2\textsuperscript{e} ton : se prononce \emph{jínglǐ}.
      \item \zh{tóngyì} 同意 (être d'accord, 2\textsuperscript{e} + 4\textsuperscript{e} ton) $\neq$ \zh{tǒngyì} 统一 (unifier).
      \item \zh{yìjiàn} 意见 (opinion, 4\textsuperscript{e} + 4\textsuperscript{e} ton).
      \end{itemize}
\end{itemize}
\end{attention}

\courssub{Formules fonctionnelles : présenter, donner son avis, relancer}

\begin{propriete}[Boîte à outils orale pour le forum]
\begin{center}
\begin{tabularx}{\linewidth}{|X|X|X|X|}
\hline
\rowcolor{popblueL} \textbf{Fonction} & \textbf{Structure chinoise} & \textbf{Pinyin} & \textbf{Équivalent français} \\
\hline
\bfseries Ouvrir / Présenter & \zh{首先，我们……} & Shǒuxiān, wǒmen… & D'abord, nous… \\
\hline
\bfseries Enchaîner & \zh{然后，在信里……} & Ránhòu, zài xìn lǐ… & Ensuite, dans la lettre… \\
\hline
\bfseries Conclure l'exposé & \zh{最后，我们用了……} & Zuìhòu, wǒmen yòng le… & Enfin, nous avons utilisé… \\
\hline
\bfseries Donner son avis & \zh{我认为 + [phrase]} & Wǒ rènwéi… & Je pense que… / Selon moi… \\
\hline
\bfseries Exprimer désaccord & \zh{我不同意，因为 + [raison]} & Wǒ bù tóngyì, yīnwèi… & Je ne suis pas d'accord, car… \\
\hline
\bfseries Poser question polie & \zh{请问，+ [question]?} & Qǐngwèn, …? & S'il vous plaît, … ? / Permettez-moi de demander… \\
\hline
\bfseries Relancer le débat (animateur) & \zh{谁有不同的意见？} & Shéi yǒu bùtóng de yìjiàn? & Qui a un avis différent ? \\
\hline
\bfseries Relancer (animateur) & \zh{你同意他的看法吗？} & Nǐ tóngyì tā de kànfǎ ma? & Es-tu d'accord avec son point de vue ? \\
\hline
\bfseries Relancer (animateur) & \zh{还有谁想发言？} & Hái yǒu shéi xiǎng fāyán? & Qui d'autre veut prendre la parole ? \\
\hline
\end{tabularx}
\end{center}
\end{propriete}

\courssub{Dialogues modèles avec pinyin (Entraînement aux tons)}

\begin{textebox}[Dialogue 1 : Présentation de la lettre (Groupe $\rightarrow$ Classe)]
\textbf{A (Groupe) :} \zh{老师好，同学们好！我们是第一组。首先，我们给CamTel公司的经理写了一封正式的信。}\\
\emph{Lǎoshī hǎo, tóngxuémen hǎo! Wǒmen shì dì-yī zǔ. Shǒuxiān, wǒmen gěi CamTel gōngsī de jīnglǐ xiě le yī fēng zhèngshì de xìn.}\\[4pt]
\textbf{B (Groupe) :} \zh{然后，我们介绍了我们的身份：我们是雅温得的高一学生，学习中文两年了。}\\
\emph{Ránhòu, wǒmen jièshào le wǒmen de shēnfèn: wǒmen shì Yǎwēndé de gāo yī xuésheng, xuéxí Zhōngwén liǎng nián le.}\\[4pt]
\textbf{C (Groupe) :} \zh{我们申请暑期实习，因为我们对5G技术很感兴趣。最后，我们写了：谢谢您的时间，期待您的回复。}\\
\emph{Wǒmen shēnqǐng shǔqī shíxí, yīnwèi wǒmen duì Wǔ-Jì jìshù hěn gǎn xìngqù. Zuìhòu, wǒmen xiě le: Xièxie nín de shíjiān, qīdài nín de huífù.}\\[4pt]
\textbf{D (Groupe) :} \zh{我们的介绍完了，谢谢大家！}\\
\emph{Wǒmen de jièshào wán le, xièxie dàjiā!}
\end{textebox}

\begin{textebox}[Dialogue 2 : Échanges après l'exposé (Classe $\rightarrow$ Groupe)]
\textbf{Élève 1 :} \zh{请问，你们为什么选择CamTel公司？}\\
\emph{Qǐngwèn, nǐmen wèishénme xuǎnzé CamTel gōngsī?}\\[4pt]
\textbf{Groupe (A) :} \zh{因为CamTel是喀麦隆最大的电信公司，有中国合作伙伴。}\\
\emph{Yīnwèi CamTel shì Kāmàilóng zuìdà de diànxìn gōngsī, yǒu Zhōngguó hézuǒ huǒbàn.}\\[4pt]
\textbf{Élève 2 :} \zh{我认为写中文信很有用，因为中国企业重视礼貌。}\\
\emph{Wǒ rènwéi xiě Zhōngwén xìn hěn yǒuyòng, yīnwèi Zhōngguó qǐyè zhòngshì lǐmào.}\\[4pt]
\textbf{Élève 3 :} \zh{我不同意，因为现在微信、邮件更快，写信太慢了。}\\
\emph{Wǒ bù tóngyì, yīnwèi xiànzài Wēixìn, yóujiàn gèng kuài, xiě xìn tài màn le.}
\end{textebox}

\courssub{Les tons à travailler : focus phonétique}

\begin{aretenir}[Points de vigilance tonale pour cette leçon]
\begin{center}
\begin{tabularx}{\linewidth}{|X|X|X|X|}
\hline
\rowcolor{poporangeL} \textbf{Mot / Groupe} & \textbf{Pinyin citation} & \textbf{Pinyin en contexte (sandhi)} & \textbf{Exercice ciblé} \\
\hline
\zh{申请} & shēn-qǐng (1-3) & shēn-qǐng & Répéter 5x : \zh{我想申请} (Wǒ xiǎng shēnqǐng) \\
\hline
\zh{经理} & jīng-lǐ (1-3) & \textbf{jíng-lǐ} (2-3) & \zh{王经理} (Wáng \textbf{jíng}lǐ) — \textbf{Attention sandhi 3+3 !} \\
\hline
\zh{贵公司} & guì-gōngsī (4-1) & guì-gōngsī & \zh{贵公司的实习} (guì gōngsī de shíxí) \\
\hline
\zh{感兴趣} & gǎn-xìngqù (3-4-4) & \textbf{gán}-xìngqù (2-4-4) & \zh{很感兴趣} (hěn \textbf{gán} xìngqù) — sandhi 3+4 \\
\hline
\zh{此致，敬礼} & cǐ-zhì, jìng-lǐ (3-4, 4-3) & cǐ-zhì, jìng-lǐ & Lire la formule finale 3x sans hésiter \\
\hline
\zh{我认为 / 我不同意} & wǒ rèn-wéi / wǒ bù tóng-yì & wǒ rèn-wéi / wǒ bù tóng-yì & Alterner : \zh{我认为…} / \zh{我不同意…} \\
\hline
\end{tabularx}
\end{center}
\textbf{Consigne prof :} Faire un « chorus » collectif sur chaque ligne du tableau avant les jeux de rôle. Insister sur le \textbf{2\textsuperscript{e} ton montant} pour les sandhi (jínglǐ, gán xìngqù).
\end{aretenir}

\courssub{Consignes}

\textbf{Tâche 1 — Comprendre le modèle (10 min).} Lis la lettre modèle à voix haute avec ton groupe. Souligne : (a) la formule d'appel, (b) la présentation de l'expéditeur, (c) la raison de la demande, (d) les qualités du candidat, (e) la formule de politesse finale. Vérifie le sens des mots nouveaux dans le tableau de vocabulaire.

\textbf{Tâche 2 — Rédiger votre lettre (20 min).} Par groupes de quatre, choisissez un type de lettre : demande de \cle{stage}, demande de \cle{bourse} ou demande d'\cle{emploi} dans une entreprise de télécommunications (CamTel, MTN, Orange ou une entreprise chinoise partenaire). Rédigez une lettre de 6 à 8 phrases en suivant le plan du modèle : appel $\rightarrow$ présentation $\rightarrow$ raison $\rightarrow$ qualités $\rightarrow$ remerciement $\rightarrow$ formule finale. Utilisez au moins \textbf{deux formules de politesse} (\zh{您}, \zh{贵公司}, \zh{谢谢您的时间}, \zh{期待您的回复}).

\textbf{Tâche 3 — Préparer le mini-exposé (15 min).} Préparez un exposé oral de 1 à 2 minutes qui présente votre lettre. Structure obligatoire : \zh{首先……} (d'abord, nous présentons…), \zh{然后……} (ensuite…), \zh{最后……} (enfin…). Répartissez la parole entre les quatre membres du groupe. Répétez en faisant attention aux tons (voir tableau « Les tons à travailler »).

\textbf{Tâche 4 — Présentation devant la classe (20 min).} Chaque groupe présente son exposé. Les auditeurs doivent :
\begin{itemize}
\item poser au moins une question commençant par \zh{请问……} (ex. : \zh{请问，你为什么想在这个公司实习？}) ;
\item donner leur avis avec \zh{我认为……} ou \zh{我不同意，因为……}.
\end{itemize}

\textbf{Tâche 5 — Débat de classe (20 min).} Un élève animateur lance le débat : \emph{« Faut-il encore apprendre à écrire des lettres formelles à l'ère des réseaux sociaux ? »} L'animateur relance la parole avec les formules : \zh{谁有不同的意见？}, \zh{你同意他的看法吗？}, \zh{还有谁想发言？}. Chaque élève doit parler au moins une fois.

\courssub{Ressources à mobiliser}

\begin{center}
\begin{tabularx}{\linewidth}{|X|X|X|X|}
\hline
\rowcolor{popblueL} Caractères & Pinyin & Nature & Traduction \\
\hline
尊敬的 & zūnjìng de & adj. & honorable, respecté (appel) \\
\hline
申请 & shēnqǐng & v. & demander, postuler (stage, emploi, bourse) \\
\hline
实习 & shíxí & v./n. & faire un stage ; stage \\
\hline
奖学金 & jiǎngxuéjīn & n. & bourse d'études \\
\hline
贵公司 & guì gōngsī & n. & votre honorable entreprise (honorifique) \\
\hline
感兴趣 & gǎn xìngqù & loc. v. & s'intéresser à (\zh{对……感兴趣}) \\
\hline
回复 & huífù & v./n. & répondre ; réponse \\
\hline
首先 & shǒuxiān & adv. & d'abord, premièrement \\
\hline
然后 & ránhòu & conj. & ensuite, puis \\
\hline
最后 & zuìhòu & adv. & enfin, finalement \\
\hline
意见 & yìjiàn & n. & opinion, avis \\
\hline
同意 & tóngyì & v. & être d'accord \\
\hline
看法 & kànfǎ & n. & point de vue, opinion \\
\hline
发言 & fāyán & v. & prendre la parole, s'exprimer \\
\hline
电信 & diànxìn & n. & télécommunications \\
\hline
合作伙伴 & hézuǒ huǒbàn & n. & partenaire de coopération \\
\hline
\end{tabularx}
\end{center}

\begin{aretenir}[Structures à réutiliser — Fiche mémo élève]
\begin{itemize}
\item \textbf{Donner son avis} : \zh{我认为 + phrase.}
      \zh{我认为写信还很重要。} (Wǒ rènwéi xiě xìn hái hěn zhòngyào.) — « Je pense qu'écrire des lettres est encore très important. »
\item \textbf{Ne pas être d'accord} : \zh{我不同意，因为 + phrase.}
      \zh{我不同意，因为电子邮件更快。} (Wǒ bù tóngyì, yīnwèi diànzǐ yóujiàn gèng kuài.) — « Je ne suis pas d'accord, car le courriel est plus rapide. »
\item \textbf{Poser une question polie} : \zh{请问，……？}
      \zh{请问，你们为什么选这个公司？} (Qǐngwèn, nǐmen wèishénme xuǎn zhège gōngsī?) — « S'il vous plaît, pourquoi avez-vous choisi cette entreprise ? »
\item \textbf{Structurer un exposé} : \zh{首先……然后……最后……}
\item \textbf{Formule d'appel} : \zh{尊敬的 + Titre/Nom + ：} (ex. \zh{尊敬的李经理：})
\item \textbf{Formule de clôture} : \zh{此致，\newline 敬礼！} (sur deux lignes, \zh{敬礼} indenté)
\end{itemize}
\end{aretenir}

\courssub{Proposition de production et corrigé commenté}

\begin{exoresolu}[Exposé modèle du groupe « Les jeunes reporters »]
Présente oralement la lettre de ton groupe en 1 à 2 minutes, avec \zh{首先/然后/最后} et deux formules de politesse.
\tcblower
\textbf{Production modèle (exposé oral) :}\\[4pt]
\zh{老师好，同学们好！我们是第三组。}\\
\emph{Lǎoshī hǎo, tóngxuémen hǎo! Wǒmen shì dì-sān zǔ.}\\
« Bonjour Madame/Monsieur, bonjour camarades ! Nous sommes le groupe 3. »\\[4pt]
\zh{首先，我们给Orange公司的经理写了一封正式的信。}\\
\emph{Shǒuxiān, wǒmen gěi Orange gōngsī de jīnglǐ xiě le yī fēng zhèngshì de xìn.}\\
« D'abord, nous avons écrit une lettre formelle au directeur de l'entreprise Orange. »\\[4pt]
\zh{然后，在信里，我们先介绍自己：我们是雅温得的学生，学习中文两年了。}\\
\emph{Ránhòu, zài xìn lǐ, wǒmen xiān jièshào zìjǐ: wǒmen shì Yǎwēndé de xuésheng, xuéxí Zhōngwén liǎng nián le.}\\
« Ensuite, dans la lettre, nous nous présentons d'abord : nous sommes des élèves de Yaoundé et nous apprenons le chinois depuis deux ans. »\\[4pt]
\zh{我们说了我们对通信技术很感兴趣，所以想申请暑期实习。}\\
\emph{Wǒmen shuō le wǒmen duì tōngxìn jìshù hěn gǎn xìngqù, suǒyǐ xiǎng shēnqǐng shǔqī shíxí.}\\
« Nous avons dit que nous nous intéressons beaucoup aux technologies de communication et que nous voulons donc demander un stage d'été. »\\[4pt]
\zh{最后，我们用了两个礼貌的句子：谢谢您的时间，期待您的回复。}\\
\emph{Zuìhòu, wǒmen yòng le liǎng ge lǐmào de jùzi: Xièxie nín de shíjiān, qīdài nín de huífù.}\\
« Enfin, nous avons utilisé deux phrases polies : merci pour votre temps, nous attendons votre réponse. »\\[4pt]
\zh{我的介绍完了，谢谢大家！}\\
\emph{Wǒ de jièshào wán le, xièxie dàjiā!}\\
« Ma présentation est terminée, merci à tous ! »\\[6pt]
\textbf{Commentaire des choix de langue :}
\begin{itemize}
\item Le plan \zh{首先 $\rightarrow$ 然后 $\rightarrow$ 最后} donne une structure claire et facile à suivre pour les auditeurs : c'est la marque d'un exposé organisé (compétence « Expression orale continue »).
\item Les formules de politesse \zh{您} (dans \zh{谢谢您的时间}) et \zh{贵公司} (implicite dans le contexte « directeur d'Orange ») répondent à la consigne et montrent la maîtrise du registre formel.
\item \zh{一封正式的信} : le classificateur des lettres est \zh{封} (fēng) ; ne dis jamais \zh{一个信} ni \zh{一张信}.
\item \zh{介绍完了} : \zh{完} (wán) est un complément de résultat signifiant « finir de, achever » ; \zh{了} (le) marque ici l'accomplissement de l'action (aspect perfectif).
\item \textbf{Prononciation} : soigne le 3\textsuperscript{e} ton de \zh{shǒuxiān} (souvent prononcé 2\textsuperscript{e} ton par sandhi si suivi d'un autre 3\textsuperscript{e} ton, mais ici isolé $\rightarrow$ 3\textsuperscript{e} ton bas) et la succession de tons de \zh{qīdài nín de huífù} (1-4 / 2 / neutre / 2-4).
\end{itemize}
\end{exoresolu}

\courssub{Exemple de lettre rédigée par un groupe (Demande de bourse)}

\begin{textebox}[Production élève : Lettre de demande de bourse — Groupe 2]
尊敬的陈主任：您好！\\
Zūnjìng de Chén zhǔrèn: Nín hǎo!\\
« Monsieur le Directeur Chen, bonjour ! »\\[4pt]
我叫张伟，是布埃阿第一中学高一学生。\\
Wǒ jiào Zhāng Wěi, shì Bù'āi'ā Dì-yī Zhōngxué gāo yī xuésheng.\\
« Je m'appelle Zhang Wei, élève de Première au Lycée n°1 de Buea. »\\[4pt]
我成绩优秀，特别喜欢中文和计算机。\\
Wǒ chéngjì yōuxiù, tèbié xǐhuan Zhōngwén hé jìsuànjī.\\
« Mes résultats sont excellents, j'aime particulièrement le chinois et l'informatique. »\\[4pt]
我想申请贵单位的“中喀友谊奖学金”。\\
Wǒ xiǎng shēnqǐng guì dānwèi de "Zhōng-Kā Yǒuyì Jiǎngxuéjīn".\\
« Je souhaite postuler à la "Bourse d'Amitié Sino-Camerounaise" de votre organisme. »\\[4pt]
谢谢您的时间，期待您的回复。\\
Xièxie nín de shíjiān, qīdài nín de huífù.\\
« Merci pour votre temps, dans l'attente de votre réponse. »\\[4pt]
此致，\\
敬礼！\\
Cǐzhì,\\
Jìnglǐ!\\[4pt]
学生：张伟　二〇二五年五月十五日\\
Xuésheng: Zhāng Wěi　Èr líng èr wǔ nián wǔ yuè shíwǔ rì
\end{textebox}

\courssub{Bilan de l'activité}

À l'issue du forum, vérifie que tu as atteint les objectifs :
\begin{itemize}
\item je sais rédiger le squelette d'une lettre formelle chinoise : appel (\zh{尊敬的……}), présentation, demande, remerciement, formule finale (\zh{此致，敬礼！}) ;
\item je sais présenter un document à l'oral avec \zh{首先}, \zh{然后}, \zh{最后} ;
\item je sais réagir dans un débat : questionner avec \zh{请问……}, donner mon avis avec \zh{我认为……}, exprimer mon désaccord avec \zh{我不同意，因为……} ;
\item je fais attention aux tons, en particulier dans \zh{shēnqǐng}, \zh{jīnglǐ} (sandhi !), \zh{tóngyì} et \zh{yìjiàn}.
\end{itemize}

\begin{attention}[Critères d'évaluation de l'enseignant — Grille simplifiée]
Ton professeur évalue trois points :
\begin{enumerate}
\item \textbf{Prononciation des tons} (20 pts) : justesse des tons lexicaux + sandhi (3+3 $\rightarrow$ 2+3 ; 3+4 $\rightarrow$ 2+4) pendant l'exposé et le débat.
\item \textbf{Usage des formules fonctionnelles} (20 pts) : politesse (\zh{您}, \zh{贵公司}, \zh{此致敬礼}), connecteurs d'exposé (\zh{首先/然后/最后}), formules d'interaction (\zh{请问}, \zh{我认为}, \zh{我不同意}, \zh{谁有不同意见}).
\item \textbf{Participation au débat} (10 pts) : prise de parole au moins une fois, écoute active, réaction pertinente aux camarades (relance, approbation, objection).
\end{enumerate}
\emph{Astuce :} Prépare à l'avance deux ou trois phrases d'opinion types (ex. \zh{我认为手写信更有温度}, \zh{我不同意，电子邮件环保}) pour ne pas rester silencieux pendant le débat.
\end{attention}

\courssub{Exercices d'entraînement (Devoirs / Remédiation)}

\begin{exercice}[Entraînement écrit : Compléter la lettre]
Complète la lettre ci-dessous avec les éléments manquants (choisis dans la boîte). Écris les caractères, le pinyin et la traduction.
\begin{center}
\begin{tabularx}{\linewidth}{|X|X|}
\hline
\zh{尊敬的} & \zh{申请} \\
\hline
\zh{贵公司} & \zh{此致，敬礼} \\
\hline
\zh{感兴趣} & \zh{期待您的回复} \\
\hline
\end{tabularx}
\end{center}

\zh{\_\_\_\_\_\_ 刘经理：您好！}\\
\zh{我叫王明，是杜阿拉第三中学高一学生。}\\
\zh{我对中非贸易很 \_\_\_\_\_\_，所以想 \_\_\_\_\_\_ \_\_\_\_\_\_ 的暑期实习。}\\
\zh{谢谢您的时间，\_\_\_\_\_\_。}\\
\zh{\_\_\_\_\_\_}\\
\zh{学生：王明　二〇二五年六月一日}
\end{exercice}

\begin{exercice}[Entraînement oral : Transformer en style formel]
Transforme les phrases familières en style lettre formelle (utilise \zh{您}, \zh{贵公司}, \zh{申请}).
\begin{enumerate}
\item \zh{我想在你公司打工。} $\rightarrow$ \zh{\_\_\_\_\_\_\_\_\_\_\_\_}
\item \zh{请回信给我。} $\rightarrow$ \zh{\_\_\_\_\_\_\_\_\_\_\_\_}
\item \zh{我是雅温得的学生。} $\rightarrow$ \zh{\_\_\_\_\_\_\_\_\_\_\_\_} (ajoute \zh{尊敬的} au début de la phrase si c'est l'ouverture)
\end{enumerate}
\end{exercice}

\begin{corrige}[Exercices d'entraînement]
\textbf{Exercice 1 (Lettre complète) :}
\zh{尊敬的 刘经理：您好！}\\
\zh{我叫王明，是杜阿拉第三中学高一学生。}\\
\zh{我对中非贸易很 \textbf{感兴趣}，所以想 \textbf{申请} \textbf{贵公司} 的暑期实习。}\\
\zh{谢谢您的时间，\textbf{期待您的回复}。}\\
\zh{\textbf{此致，}}\\
\zh{\textbf{敬礼！}}\\
\zh{学生：王明　二〇二五年六月一日}

\textbf{Exercice 2 (Transformation) :}
\begin{enumerate}
\item \zh{我想申请贵公司的实习。} (ou \zh{我想申请在贵公司实习。}) — \emph{Remplace « 打工 » (familier) par « 实习 » + « 贵公司 ».}
\item \zh{期待您的回复。} (ou \zh{恳请您回复。}) — \emph{« 请回信 » est trop direct ; « 期待/恳请 » + « 您 » est la norme.}
\item \zh{尊敬的领导：我是雅温得的学生。} — \emph{L'appel \zh{尊敬的……} ouvre la lettre.}
\end{enumerate}
\end{corrige}

\newpage

\coursec{Méthodes et exercices résolus}

\courssub{Savoir-faire 1 : Présenter oralement un exposé simple en chinois}

\begin{methode}[Présenter un exposé simple]
Pour présenter oralement un sujet simple (par exemple : la lettre formelle), suis toujours le même plan en trois temps.
\begin{enumerate}
\item \textbf{Saluer et annoncer le sujet.} Commence par une salutation : \zh{大家好！} (\emph{Dàjiā hǎo!} « Bonjour à tous ! »), puis annonce : \zh{我今天要讲的是……} (\emph{Wǒ jīntiān yào jiǎng de shì...} « Ce dont je vais parler aujourd'hui, c'est... »).
\item \textbf{Développer en 2 ou 3 idées reliées.} Utilise les connecteurs \zh{首先} (\emph{shǒuxiān} « d'abord »), \zh{然后} (\emph{ránhòu} « ensuite »), \zh{最后} (\emph{zuìhòu} « enfin »). Une idée = une phrase courte : Sujet + Verbe + Complément.
\item \textbf{Conclure et remercier.} Termine par : \zh{我的介绍完了，谢谢大家！} (\emph{Wǒ de jièshào wán le, xièxie dàjiā!} « Ma présentation est terminée, merci à tous ! »).
\end{enumerate}
Réflexes : phrases courtes, aucun mot que tu ne maîtrises pas, regard vers le public, débit lent pour poser les tons.
\end{methode}

\begin{center}
\begin{tabularx}{\linewidth}{|X|X|X|X|}
\hline \rowcolor{popblueL} \textbf{Caractères} & \textbf{Pinyin} & \textbf{Nature} & \textbf{Traduction} \\ \hline
大家好 & dàjiā hǎo & loc. & Bonjour à tous \\ \hline
讲 & jiǎng & v. & parler, expliquer \\ \hline
正式 & zhèngshì & adj. & formel, officiel \\ \hline
信 & xìn & n. & lettre \\ \hline
校长 & xiàozhǎng & n. & proviseur, directeur d'école \\ \hline
称呼 & chēnghu & v. & s'adresser à, appeler \\ \hline
尊敬的 & zūnjìng de & adj. & respecté (formule de politesse) \\ \hline
结尾 & jiéwěi & n. & fin, conclusion \\ \hline
此致敬礼 & cǐ zhì jìnglǐ & loc. & Veuillez agréer mes salutations respectueuses \\ \hline
介绍 & jièshào & v./n. & présenter / présentation \\ \hline
\end{tabularx}
\end{center}

\begin{exoresolu}[Exposé : la lettre formelle]
\textbf{Énoncé.} Prépare un mini-exposé oral de 5 phrases sur le thème « la lettre formelle en chinois », en utilisant \zh{首先}, \zh{然后}, \zh{最后} et en terminant par un remerciement.
\tcblower
\textbf{Raisonnement.} On applique le plan en trois temps : salutation + annonce du sujet (1 phrase), développement avec les trois connecteurs (3 idées : définition, destinataires, structure), conclusion + remerciement.

\textbf{Réponse rédigée.}\\
大家好！我今天要讲的是正式的信。\\
\emph{Dàjiā hǎo! Wǒ jīntiān yào jiǎng de shì zhèngshì de xìn.}\\
« Bonjour à tous ! Aujourd'hui, je vais parler de la lettre formelle. »\\
首先，正式的信是写给老师、校长或者领导的。\\
\emph{Shǒuxiān, zhèngshì de xìn shì xiě gěi lǎoshī, xiàozhǎng huòzhě lǐngdǎo de.}\\
« D'abord, une lettre formelle s'écrit à un professeur, à un proviseur ou à un supérieur. »\\
然后，信的开头要称呼对方，比如说“尊敬的老师”。\\
\emph{Ránhòu, xìn de kāitóu yào chēnghu duìfāng, bǐrú shuō “zūnjìng de lǎoshī”.}\\
« Ensuite, au début de la lettre, il faut s'adresser au destinataire, par exemple "Monsieur le Professeur". »\\
最后，信的结尾要写“此致敬礼”和自己的名字。\\
\emph{Zuìhòu, xìn de jiéwěi yào xiě “cǐ zhì jìnglǐ” hé zìjǐ de míngzi.}\\
« Enfin, à la fin de la lettre, on écrit la formule de politesse et son nom. »\\
我的介绍完了，谢谢大家！\\
\emph{Wǒ de jièshào wán le, xièxie dàjiā!} « Ma présentation est terminée, merci à tous ! »

\textbf{Justifications grammaticales.} \zh{是写给……的} : structure emphatique \zh{是……的} pour insister sur le destinataire. \zh{比如说} : « par exemple », suivi d'une citation entre guillemets chinois “ ”. \zh{要} exprime l'obligation (« il faut »). \zh{领导} (\emph{lǐngdǎo}) est plus précis que \zh{公司} pour désigner un destinataire humain. \textbf{Conclusion.} Un exposé réussi = plan clair + connecteurs + phrases courtes + tons posés.
\end{exoresolu}

\courssub{Savoir-faire 2 : Échanger avec ses camarades dans un petit débat}

\begin{methode}[Donner son avis, réagir et relancer]
Un débat guidé repose sur trois gestes de communication.
\begin{enumerate}
\item \textbf{Donner son avis.} \zh{我觉得……} (\emph{Wǒ juéde...} « je pense que... », courant), \zh{我认为……} (\emph{Wǒ rènwéi...} « je considère que... », plus formel), \zh{在我看来，……} (\emph{zài wǒ kànlái} « à mon avis »).
\item \textbf{Être d'accord ou pas d'accord.} \zh{我同意你的看法。} (\emph{Wǒ tóngyì nǐ de kànfǎ.} « Je suis d'accord avec ton point de vue. ») ; \zh{我不同意，因为……} (\emph{Wǒ bù tóngyì, yīnwèi...} « Je ne suis pas d'accord, parce que... ») ; \zh{你说得对，可是……} (\emph{Nǐ shuō de duì, kěshì...} « Tu as raison, mais... »).
\item \textbf{Relancer la parole.} \zh{你呢？你怎么看？} (\emph{Nǐ ne? Nǐ zěnme kàn?} « Et toi ? Qu'en penses-tu ? »), \zh{谁还有别的意见？} (\emph{Shéi hái yǒu bié de yìjiàn?} « Qui a encore un autre avis ? »).
\end{enumerate}
Réflexe d'or : toujours justifier avec \zh{因为……所以……} (\emph{yīnwèi... suǒyǐ...} « parce que... donc... »). En chinois, on garde souvent les deux connecteurs dans la même phrase, contrairement au français où « donc » est souvent omis.
\end{methode}

\begin{center}
\begin{tabularx}{\linewidth}{|X|X|X|X|}
\hline \rowcolor{popgreenL} \textbf{Caractères} & \textbf{Pinyin} & \textbf{Nature} & \textbf{Traduction} \\ \hline
觉得 & juéde & v. & penser, trouver (avis subjectif) \\ \hline
认为 & rènwéi & v. & considérer que (avis réfléchi) \\ \hline
看法 & kànfǎ & n. & point de vue, opinion \\ \hline
同意 & tóngyì & v. & être d'accord \\ \hline
因为 & yīnwèi & conj. & parce que \\ \hline
所以 & suǒyǐ & conj. & donc, c'est pourquoi \\ \hline
可是 & kěshì & conj. & mais, cependant \\ \hline
怎么看 & zěnme kàn & loc. & qu'en penses-tu ? \\ \hline
意见 & yìjiàn & n. & avis, opinion \\ \hline
电子邮件 & diànzǐ yóujiàn & n. & courriel, e-mail \\ \hline
\end{tabularx}
\end{center}

\begin{exoresolu}[Débat : lettre ou e-mail ?]
\textbf{Énoncé.} Complète ce débat entre deux élèves de Première A avec les formules fonctionnelles appropriées, puis traduis les répliques de B (1) et de A (2).\\
A : 我觉得写信比发电子邮件好。\\
B : (1) ————，因为电子邮件很快。\\
A : 你说得对，(2) ———— 手写的信更礼貌。\\
B : (3) ————？（relancer un troisième camarade）
\tcblower
\textbf{Raisonnement.} (1) B s'oppose à A et justifie avec \zh{因为} : il faut une formule de désaccord. (2) A concède partiellement (\zh{你说得对}) puis nuance : il faut \zh{可是} « mais ». (3) Il faut une formule de relance interrogative.

\textbf{Réponses.}\\
(1) \zh{我不同意} (\emph{Wǒ bù tóngyì} « Je ne suis pas d'accord ») — la formule de désaccord direct, suivie de la justification \zh{因为电子邮件很快} « parce que l'e-mail est rapide ».\\
(2) \zh{可是} (\emph{kěshì} « mais ») — la concession \zh{你说得对} « tu as raison » appelle obligatoirement une opposition : \zh{可是手写的信更礼貌} « mais une lettre écrite à la main est plus polie ». Remarque : \zh{更} (\emph{gèng}) marque le comparatif de supériorité « plus... ».\\
(3) \zh{你呢？你怎么看？} (\emph{Nǐ ne? Nǐ zěnme kàn?} « Et toi ? Qu'en penses-tu ? ») — double relance : \zh{呢} renvoie la question, \zh{你怎么看} demande l'opinion.

\textbf{Traductions.}\\
Réplique de B (1) : « Je ne suis pas d'accord, parce que l'e-mail est rapide. »\\
Réplique de A (2) : « Tu as raison, mais une lettre écrite à la main est plus polie. »

\textbf{Conclusion.} Dans un débat, chaque prise de parole = formule d'opinion + justification avec \zh{因为}, et on n'oublie jamais de relancer un camarade.
\end{exoresolu}

\courssub{Savoir-faire 3 : Prononcer correctement les tons}

\begin{methode}[Travailler les tons à l'oral]
\begin{enumerate}
\item \textbf{Identifier le ton de chaque syllabe} avant de parler : 1\textsuperscript{er} ton haut et plat (mā), 2\textsuperscript{e} ton montant (má), 3\textsuperscript{e} ton descendant-montant (mǎ), 4\textsuperscript{e} ton descendant sec (mà), ton neutre léger et court (ma).
\item \textbf{Appliquer la règle des deux 3\textsuperscript{es} tons (sandhi).} Deux 3\textsuperscript{es} tons qui se suivent : le premier devient un 2\textsuperscript{e} ton. \zh{你好} se prononce \emph{ní hǎo} (et non nǐ hǎo) ; \zh{很好} se prononce \emph{hén hǎo}.
\item \textbf{Poser les tons dans les formules de politesse.} \zh{谢谢} \emph{xièxie} (4\textsuperscript{e} + neutre), \zh{老师} \emph{lǎoshī} (3\textsuperscript{e} + 1\textsuperscript{er}), \zh{尊敬} \emph{zūnjìng} (1\textsuperscript{er} + 4\textsuperscript{e}), \zh{敬礼} \emph{jìnglǐ} (4\textsuperscript{e} + 3\textsuperscript{e}).
\item \textbf{S'entraîner en paires minimales.} Répéter des mots proches : \zh{买} \emph{mǎi} « acheter » / \zh{卖} \emph{mài} « vendre » ; \zh{问} \emph{wèn} « demander » / \zh{闻} \emph{wén} « sentir » ; \zh{写} \emph{xiě} « écrire » / \zh{血} \emph{xuè} « sang » (attention voyelle).
\end{enumerate}
Réflexe : parler lentement vaut mieux que parler vite avec des tons faux ; un ton faux change le sens du mot.
\end{methode}

\begin{center}
\begin{tabularx}{\linewidth}{|X|X|X|X|}
\hline \rowcolor{poporangeL} \textbf{Caractères} & \textbf{Pinyin} & \textbf{Nature} & \textbf{Traduction} \\ \hline
买 & mǎi & v. & acheter \\ \hline
卖 & mài & v. & vendre \\ \hline
问 & wèn & v. & demander \\ \hline
闻 & wén & v. & sentir (odeur) \\ \hline
写 & xiě & v. & écrire \\ \hline
血 & xuè & n. & sang \\ \hline
封 & fēng & m. & classificateur (lettres, enveloppes) \\ \hline
辅导 & fǔdǎo & v. & tutorer, donner des cours de soutien \\ \hline
\end{tabularx}
\end{center}

\begin{exoresolu}[Dictée de tons et correction]
\textbf{Énoncé.} 1) Indique les tons des mots suivants et dis lesquels subissent la règle de sandhi du 3\textsuperscript{e} ton : 你好、老师、敬礼、看法、同意。 2) Un élève *veut dire* « J'achète une lettre » (\emph{wǒ mǎi yì fēng xìn}) mais le professeur comprend « Je vends une lettre ». Identifie l'erreur de ton et corrige.
\tcblower
\textbf{1) Tons et sandhi.}
\begin{itemize}
\item \zh{你好} : \emph{nǐ hǎo} (3\textsuperscript{e} + 3\textsuperscript{e}) $\rightarrow$ sandhi obligatoire, prononcé \emph{ní hǎo}.
\item \zh{老师} : \emph{lǎo shī} (3\textsuperscript{e} + 1\textsuperscript{er}) $\rightarrow$ pas de sandhi (le 2\textsuperscript{e} caractère n'est pas un 3\textsuperscript{e} ton).
\item \zh{敬礼} : \emph{jìng lǐ} (4\textsuperscript{e} + 3\textsuperscript{e}) $\rightarrow$ pas de sandhi.
\item \zh{看法} : \emph{kàn fǎ} (4\textsuperscript{e} + 3\textsuperscript{e}) $\rightarrow$ pas de sandhi.
\item \zh{同意} : \emph{tóng yì} (2\textsuperscript{e} + 4\textsuperscript{e}) $\rightarrow$ pas de sandhi.
\end{itemize}
Seul \zh{你好} subit le sandhi du 3\textsuperscript{e} ton.

\textbf{2) Correction.} L'élève a prononcé \zh{卖} \emph{mài} (4\textsuperscript{e} ton, « vendre ») au lieu de \zh{买} \emph{mǎi} (3\textsuperscript{e} ton, « acheter »). La phrase correcte est : \zh{我买一封信。} \emph{Wǒ mǎi yì fēng xìn.} « J'achète une lettre. »
\textbf{Justifications :}
\begin{itemize}
\item Le 3\textsuperscript{e} ton de \zh{买} descend puis remonte ; le 4\textsuperscript{e} ton de \zh{卖} chute brutalement.
\item Le classificateur des lettres est \zh{封} (\emph{fēng}, 1\textsuperscript{er} ton) : \zh{一封信} \emph{yì fēng xìn} « une lettre ».
\item Règle du \zh{一} (\emph{yī}) : devant un ton 1, 2 ou 3, \zh{一} se prononce au 4\textsuperscript{e} ton \emph{yì} ; devant un ton 4, il se prononce au 2\textsuperscript{e} ton \emph{yí}. Ici \zh{封} est ton 1 $\rightarrow$ \emph{yì fēng}.
\end{itemize}
\textbf{Conclusion.} Un seul ton faux (\emph{mǎi} $\rightarrow$ \emph{mài}) inverse le sens de la phrase : à l'oral comme à l'écrit, vérifie toujours les tons des verbes et les classificateurs.
\end{exoresolu}

\courssub{Savoir-faire 4 : Lire et présenter une lettre formelle à voix haute}

\begin{methode}[Lire une lettre formelle à voix haute]
\begin{enumerate}
\item \textbf{Repérer la structure} avant de lire : appel (\zh{尊敬的老师：} « Monsieur le Professeur : »), corps (salutation \zh{您好！}, motif, demande avec \zh{我想请您……} « je voudrais vous demander de... »), formule finale (\zh{此致敬礼！} « Veuillez agréer mes salutations respectueuses »), signature et date à droite.
\item \textbf{Marquer les pauses} : une pause après l'appel, une après chaque phrase du corps, une voix plus lente et solennelle pour la formule de politesse.
\item \textbf{Respecter le vouvoiement} : \zh{您} (\emph{nín} « vous », politesse) et jamais \zh{你} dans une lettre formelle.
\item \textbf{Lire la date à la chinoise} : année + \zh{年} + mois + \zh{月} + jour + \zh{日} : \zh{二〇二五年五月十日} \emph{èr líng èr wǔ nián wǔ yuè shí rì} « le 10 mai 2025 ». Les chiffres de l'année se lisent chiffre par chiffre.
\end{enumerate}
\end{methode}

\begin{textebox}[Modèle de lettre formelle : demande de soutien]
尊敬的张老师：\\
您好！我是高一A班的学生。我想请您星期五下午给我辅导汉语。谢谢您！\\
此致\\
敬礼！\\
学生：李明\\
二〇二五年五月十日
\end{textebox}

\begin{exoresolu}[Lecture et questions sur une lettre]
\textbf{Énoncé.} Lis la lettre ci-dessus à voix haute, puis réponds aux questions.\\
Questions : 1) À qui la lettre est-elle adressée ? 2) Que demande l'élève ? 3) Relevez la formule de politesse finale et traduisez-la. 4) Où place-t-on la date et la signature ?
\tcblower
\textbf{1)} La lettre est adressée au professeur Zhang : \zh{尊敬的张老师} \emph{zūnjìng de Zhāng lǎoshī} « Respecté professeur Zhang ». L'appel se termine par deux-points chinois ：et se place en haut à gauche, collé à la marge.

\textbf{2)} L'élève demande un soutien en chinois vendredi après-midi : \zh{我想请您星期五下午给我辅导汉语。} \emph{Wǒ xiǎng qǐng nín xīngqīwǔ xiàwǔ gěi wǒ fǔdǎo Hànyǔ.} « Je voudrais vous demander de me donner du soutien en chinois vendredi après-midi. » Justifications : \zh{想} + verbe « vouloir » ; \zh{请您} « vous demander (poliment) » ; le complément de temps \zh{星期五下午} se place \emph{avant} le verbe ; \zh{给我} « pour moi » introduit le bénéficiaire.

\textbf{3)} La formule finale est \zh{此致敬礼！} \emph{Cǐ zhì jìnglǐ!} « Par la présente, je vous adresse mes salutations respectueuses » (équivalent de « Veuillez agréer... »). Remarque de mise en page : \zh{此致} seul sur une ligne, \zh{敬礼} sur la ligne suivante, collé à gauche, suivi d'un point d'exclamation.

\textbf{4)} La signature (\zh{学生：李明}) et la date (\zh{二〇二五年五月十日}) se placent en bas à droite, la date sous la signature.

\textbf{Conclusion.} Une lettre formelle chinoise = appel respectueux + \zh{您好} + demande claire avec \zh{您} + \zh{此致敬礼} + signature et date. Lire à voix haute en respectant ces blocs, c'est déjà moitié de la note.
\end{exoresolu}

\begin{experience}[Jeu de rôle : « Au secrétariat / Chez le proviseur »]
\textbf{Consigne.} Par groupes de 3. L'un joue l'élève, l'autre le proviseur (ou un professeur), le troisième observe et note les tons.
\begin{enumerate}
\item L'élève lit sa lettre (texte ci-dessus ou une variante : demande de stage, d'autorisation d'absence, de recommandation).
\item Le « proviseur » répond oralement : \zh{好的，没问题。} (\emph{Hǎo de, méi wèntí.} « D'accord, pas de problème. ») ou \zh{请稍等，我考虑一下。} (\emph{Qǐng shāo děng, wǒ kǎolǜ yíxià.} « Veuillez patienter, je vais réfléchir. »).
\item L'observateur vérifie : tons de \zh{您} / \zh{尊敬} / \zh{敬礼} / \zh{辅导} ; respect des pauses ; utilisation de \zh{您} et non \zh{你}.
\item Rotation des rôles.
\end{enumerate}
\textbf{Durée :} 15 min. \textbf{Production attendue :} chaque élève a lu une lettre complète avec tons corrects.
\end{experience}

\begin{savaistu}[Culture épistolaire : Cameroun et Chine]
\textbf{Au Cameroun,} les lettres administratives (demandes de stage, de bourse, lettres de motivation) suivent le modèle français : en-tête (expéditeur/destinateur), objet, formules « Monsieur le Directeur, je vous prie d'agréer... ». L'e-mail a largement remplacé le papier, mais la forme reste codifiée.

\textbf{En Chine,} la lettre formelle (\zh{公文} ou \zh{正式信函}) conserve une structure très ritualisée :
\begin{itemize}
\item L'appel \zh{尊敬的……:} (Respecté...) est impératif.
\item Le vouvoiement \zh{您} (\emph{nín}) s'écrit avec le radical « cœur » \zh{心} en bas (cœur respectueux), distinct de \zh{你} (homme \zh{亻}).
\item La formule \zh{此致敬礼} (litt. « par la présente, saluer respectueusement ») est immuable ; \zh{此致} est aligné à gauche sur une ligne, \zh{敬礼} sur la suivante, rentré de deux caractères (ou aligné gauche selon les normes actuelles).
\item La date s'écrit \zh{XXXX年XX月XX日} (année-mois-jour), souvent avec des chiffres chinois formels \zh{零一二三四五六七八九} pour l'année dans les documents très officiels.
\item \textbf{Différence clé :} Au Cameroun, on signe souvent à droite \emph{au-dessus} de son nom tapé ; en Chine, la signature manuscrite (\zh{签名}) est cruciale, placée sous le nom imprimé (\zh{姓名}), souvent accompagnée du cachet rouge (\zh{盖章}) pour les organismes.
\end{itemize}
\textbf{Point commun :} Dans les deux cultures, la lettre formelle engage la responsabilité de l'auteur ; le ton, la clarté et le respect des codes sont gages de sérieux.
\end{savaistu}

\begin{attention}[Les 5 erreurs qui coûtent des points au Bac]
\begin{enumerate}
\item \textbf{Tons négligés à l'oral.} Confondre \zh{买} \emph{mǎi} « acheter » et \zh{卖} \emph{mài} « vendre », ou oublier le sandhi de \zh{你好} (prononcé \emph{ní hǎo}). Un ton faux = un mot faux = incompréhension.
\item \textbf{Tutoyer dans une lettre formelle.} Écrire \zh{你} au lieu de \zh{您} (\emph{nín}) en s'adressant au professeur ou au proviseur : faute de registre éliminatoire. \zh{您} porte le radical cœur \zh{心} : le respect vient du cœur.
\item \textbf{Ordre des mots calqué du français.} Le complément de temps se place \emph{avant} le verbe : \zh{我星期五下午学习} et jamais \zh{*我学习星期五下午}. De même : lieu avant verbe \zh{我在学校写信} « j'écris une lettre à l'école ».
\item \textbf{Classificateur oublié ou faux.} Une lettre se compte avec \zh{封} : \zh{一封信} \emph{yì fēng xìn} « une lettre ». Ne jamais dire « \zh{一个信} » (\emph{yí gè xìn}). \zh{封} vient de « sceller une enveloppe ».
\item \textbf{Formules finales inventées ou mal placées.} Ne traduis pas « Cordialement » mot à mot : la formule consacrée est \zh{此致敬礼！}. Mise en page : \zh{此致} ligne à part, \zh{敬礼} ligne suivante. N'oublie ni la signature (\zh{学生：姓名}), ni la date complète avec \zh{年、月、日}.
\end{enumerate}
\end{attention}

\newpage

\coursec{Exercices}

\coursec{Leçon 41 — Expression orale : rédaction d'une lettre formelle}

\courssub{Découverte : La lettre formelle chinoise}

\begin{textebox}[Modèle de lettre formelle : Demande de stage]
\zh{尊敬的王经理：} \\
\zh{您好！} \\
\zh{我叫李华，是喀麦隆大学孔子学院的学生。我写信是想申请贵公司的暑期实习机会。} \\
\zh{我认为我有足够的中文基础，而且工作认真、刻苦。因为贵公司是中喀合资的知名企业，所以我非常希望能来贵公司学习。} \\
\zh{此致} \\
\zh{敬礼！} \\
\zh{署名：李华} \\
\zh{日期：二〇二四年十月十五日}
\end{textebox}

\begin{propriete}[Structure type d'une lettre formelle (5 parties)]
Une lettre formelle chinoise (\zh{公函} gōng hán / \zh{正式信函} zhèng shì xìn hán) suit un ordre immuable de haut en bas :
\begin{enumerate}
\item \textbf{称呼} (\zh{chēng hū}) : Appel. Aligné à gauche, suivi de deux points \zh{：}. Ex: \zh{尊敬的经理：} / \zh{尊敬的王总：}.
\item \textbf{问候语} (\zh{wèn hòu yǔ}) : Salutation. Nouvelle ligne, souvent indentée de deux caractères. Ex: \zh{您好！}.
\item \textbf{正文} (\zh{zhèng wén}) : Corps de la lettre. Paragraphes indentés de deux caractères. Contient l'objet, l'argumentation (\zh{因为……所以……}, \zh{首先……然后……最后……}).
\item \textbf{祝颂语} (\zh{zhù sòng yǔ}) : Formules de politesse finales. Deux lignes : \zh{此致} (aligné à gauche ou indenté) puis \zh{敬礼} (aligné à gauche, sur la ligne suivante). Autres formules : \zh{顺颂商祺} (affaires), \zh{祝您工作顺利} (courant).
\item \textbf{署名 与 日期} (\zh{shǔ míng yǔ rì qī}) : Signature (nom) et date. Alignés à droite. Format date : \zh{二〇二四年十月十五日}.
\end{enumerate}
\end{propriete}

\begin{definition}[Vocabulaire clé de la lettre formelle]
\begin{center}
\begin{tabularx}{\linewidth}{|X|X|X|X|}
\hline
\rowcolor{popblueL} \textbf{Caractères} & \textbf{Pinyin} & \textbf{Nature} & \textbf{Traduction} \\
\hline
\zh{尊敬的} & zūn jìng de & Adj. & Respecté (formule d'appel, précède titre/nom) \\
\hline
\zh{贵公司} & guì gōng sī & Nom (honorifique) & Votre société / entreprise \\
\hline
\zh{申请} & shēn qǐng & Verbe & Solliciter, postuler, demander formellement \\
\hline
\zh{实习} & shí xí & Verbe/Nom & Stage / Faire un stage \\
\hline
\zh{机会} & jī huì & Nom & Opportunité, chance \\
\hline
\zh{此致} & cǐ zhì & Locution & C'est tout / En terminant (début formule finale) \\
\hline
\zh{敬礼} & jìng lǐ & Verbe & Saluer respectueusement (fin formule finale) \\
\hline
\zh{署名} & shǔ míng & Nom/Verbe & Signature / Signer \\
\hline
\zh{打扰了} & dǎ rǎo le & Locution & Je vous ai dérangé (formule de congé oral/écrit) \\
\hline
\end{tabularx}
\end{center}
\end{definition}

\courssub{Formules fonctionnelles : Présenter, donner son avis, relancer}

\begin{propriete}[Présenter un exposé / une idée]
\begin{center}
\begin{tabularx}{\linewidth}{|X|X|X|}
\hline
\rowcolor{popgreenL} \textbf{Fonction} & \textbf{Structure / Expression} & \textbf{Exemple (Caractères + Pinyin + Traduction)} \\
\hline
Annoncer le sujet & \zh{今天我要谈谈……} (Jīn tiān wǒ yào tán tan…) & \zh{今天我要谈谈正式信函的格式。} (Aujourd'hui je vais parler du format des lettres formelles.) \\
\hline
Structurer (chronologique/logique) & \zh{首先……，然后……，最后……} (Shǒu xiān…, rán hòu…, zuì hòu…) & \zh{首先写称呼，然后写正文，最后写署名。} (D'abord l'appel, puis le corps, enfin la signature.) \\
\hline
Donner un exemple & \zh{例如……} (Lì rú…) / \zh{举个例子……} (Jǔ gè lì zi…) & \zh{例如：称呼写“尊敬的经理：”。} (Par ex. : l'appel s'écrit « Respecté Directeur : ».) \\
\hline
Conclure & \zh{总之……} (Zǒng zhī…) / \zh{以上就是我的发言。} (Yǐ shàng jiù shì wǒ de fā yán.) & \zh{总之，格式很重要。} (En résumé, le format est important.) \\
\hline
\end{tabularx}
\end{center}
\end{propriete}

\begin{propriete}[Exprimer son avis / Accord / Désaccord]
\begin{center}
\begin{tabularx}{\linewidth}{|X|X|X|}
\hline
\rowcolor{poporangeL} \textbf{Fonction} & \textbf{Expression} & \textbf{Exemple} \\
\hline
Opinion personnelle & \zh{我认为……} (Wǒ rèn wéi…) / \zh{我觉得……} (Wǒ jué de…) & \zh{我认为写信更正式。} (Je pense qu'écrire une lettre est plus formel.) \\
\hline
Accord & \zh{我同意……} (Wǒ tóng yì…) / \zh{赞成。} (Zàn chéng.) & \zh{我同意你的看法。} (Je suis d'accord avec ton avis.) \\
\hline
Désaccord poli & \zh{我不同意……} (Wǒ bù tóng yì…) / \zh{我不太赞同。} (Wǒ bù tài zàn tóng.) & \zh{我不太赞同，因为邮件更快。} (Je ne suis pas tout à fait d'accord, car l'email est plus rapide.) \\
\hline
Nuancer (Bien que... mais...) & \zh{虽然……，但是/可是……} (Suī rán…, dàn shì / kě shì…) & \zh{虽然发邮件快，但是写信更郑重。} (Bien que l'email soit rapide, la lettre est plus solennelle.) \\
\hline
Argumenter (Cause/Effet) & \zh{因为……，所以……} (Yīn wèi…, suǒ yǐ…) & \zh{因为是公文，所以要用正式语言。} (Car c'est un document officiel, il faut un langage formel.) \\
\hline
\end{tabularx}
\end{center}
\end{propriete}

\begin{propriete}[Relancer l'interlocuteur (Interaction orale)]
Pour maintenir l'échange et inviter l'autre à parler :
\begin{center}
\begin{tabularx}{\linewidth}{|X|X|X|}
\hline
\rowcolor{poppurpleL} \textbf{Expression} & \textbf{Pinyin} & \textbf{Nuance / Contexte} \\
\hline
\zh{你觉得呢？} & nǐ jué de ne? & Standard, neutre : « Qu'en penses-tu ? » \\
\hline
\zh{对吧？} & duì ba? & Recherche confirmation : « C'est ça, non ? / N'est-ce pas ? » \\
\hline
\zh{你有什么建议？} & nǐ yǒu shén me jiàn yì? & Ouvert : « Quelles suggestions as-tu ? » \\
\hline
\zh{请说。} & qǐng shuō. & Poli, formel : « Je vous en prie, parlez. » \\
\hline
\end{tabularx}
\end{center}
\end{propriete}

\courssub{Les tons à travailler : Règle du 3e ton (Sandhi)}

\begin{propriete}[Règle du changement de ton du 3e ton (三声变调)]
Lorsque \textbf{deux 3e tons} se suivent, \textbf{le premier devient un 2e ton} (montant).
\begin{itemize}
\item \zh{你好} : nǐ (3) + hǎo (3) $\rightarrow$ \textbf{ní hǎo} (2 + 3).
\item \zh{很好} : hěn (3) + hǎo (3) $\rightarrow$ \textbf{hén hǎo} (2 + 3).
\item \zh{马上} : mǎ (3) + shàng (4) $\rightarrow$ Pas de changement (2e syllabe n'est pas 3e ton).
\item \zh{请假} : qǐng (3) + jià (4) $\rightarrow$ Pas de changement.
\end{itemize}
\emph{Attention :} L'écriture en pinyin dictionnaire garde le 3e ton (nǐ hǎo), mais l'oral suit la règle. En classe, on pratique la prononciation réelle.
\end{propriete}

\begin{attention}[Piège fréquent : \zh{尊敬的} n'a PAS de 3e ton]
Dans le mot \zh{尊敬的} (zūn jìng de) :
\begin{itemize}
\item \zh{尊} (zūn) = 1er ton (haut plat).
\item \zh{敬} (jìng) = 4e ton (descendant).
\item \zh{的} (de) = ton neutre.
\end{itemize}
\textbf{Il n'y a aucun 3e ton, donc aucune sandhi ne s'applique.} Ne prononcez pas *zún jǐng de*. C'est une erreur très fréquente due à la confusion visuelle avec d'autres mots.
\end{attention}

\courssub{Dialogues modèles avec pinyin}

\begin{textebox}[Dialogue 1 : Appel téléphonique pour stage (Formel)]
\zh{A : 尊敬的王经理，您好！我是喀麦隆大学的学生李华。} (Zūn jìng de Wáng jīng lǐ, nín hǎo! Wǒ shì Kāmàilóng dà xué de xué sheng Lǐ Huá.) \\
\zh{B : 您好，李同学。请问有什么事？} (Nín hǎo, Lǐ tóng xué. Qǐng wèn yǒu shén me shì?) \\
\zh{A : 我想申请贵公司的暑期实习机会。我认为我的中文水平可以胜任。} (Wǒ xiǎng shēn qǐng guì gōng sī de shǔ qī shí xí jī huì. Wǒ rèn wéi wǒ de zhōng wén shuǐ píng kě yǐ shèng rèn.) \\
\zh{B : 很好。那请您把简历发到公司邮箱。我们会尽快回复。} (Hěn hǎo. Nà qǐng nín bǎ jiǎn lì fā dào gōng sī yóu xiāng. Wǒ men huì jǐn kuài huí fù.) \\
\zh{A : 好的，谢谢您！祝您工作顺利，打扰了。} (Hǎo de, xiè xie nín! Zhù nín gōng zuò shùn lì, dǎ rǎo le.) \\
\zh{B : 不客气，再见。} (Bú kè qi, zài jiàn.)
\end{textebox}

\begin{textebox}[Dialogue 2 : Débat "Lettre vs Email" (Interaction)]
\zh{甲 : 我认为写信比发邮件好。因为写信很正式，所以表示尊重。你觉得呢？} (Wǒ rèn wéi xiě xìn bǐ fā yóu jiàn hǎo. Yīn wèi xiě xìn hěn zhèng shì, suǒ yǐ biǎo shì zūn zhòng. Nǐ jué de ne?) \\
\zh{乙 : 我不同意。虽然写信郑重，但是发邮件更快速、方便、还环保。对吧？} (Wǒ bù tóng yì. Suī rán xiě xìn zhèng zhòng, dàn shì fā yóu jiàn gèng kuài sù, fāng biàn, hái huán bǎo. Duì ba?) \\
\zh{甲 : 一方面你说得对，但另一方面，正式场合必须写信。例如签合同。} (Yī fāng miàn nǐ shuō de duì, dàn lìng wài yī fāng miàn, zhèng shì chǎng hé bì xū xiě xìn. Lì rú qiān hé tong.) \\
\zh{乙 : 这点我同意。看来要看场合。} (Zhè diǎn wǒ tóng yì. Kàn lái yào kàn chǎng hé.)
\end{textebox}

\courssub{Je vérifie mes connaissances}

\begin{exercice}[QCM : Structure d'une lettre formelle]
Parmi les propositions suivantes, une seule réponse est exacte. Choisissez la bonne lettre (A, B, C ou D) et recopiez la phrase complète sur votre copie.

\begin{enumerate}
\item Dans une lettre formelle chinoise, la formule de politesse finale \zh{此致敬礼} (cǐ zhì jìng lǐ) se place :
      \begin{itemize}
      \item A. Au tout début, après l'en-tête.
      \item B. Juste avant la signature (\zh{署名} shǔ míng), alignée à gauche (sur deux lignes).
      \item C. Juste après la signature, alignée à droite.
      \item D. Au milieu du corps de la lettre (\zh{正文} zhèng wén).
      \end{itemize}

\item L'appel (\zh{称呼} chēng hū) \zh{尊敬的经理} (zūn jìng de jīng lǐ) s'écrit :
      \begin{itemize}
      \item A. Au centre de la ligne, suivi d'un point.
      \item B. En haut à gauche, suivi de deux points (\zh{：}).
      \item C. En bas à droite, précédé d'une virgule.
      \item D. N'importe où, tant qu'il est en gras.
      \end{itemize}

\item Pour \textbf{relancer} poliment son interlocuteur dans un échange oral sur la rédaction d'une lettre, on utilise de préférence :
      \begin{itemize}
      \item A. \zh{你说什么？} (Nǐ shuō shén me?) — Trop direct, familier.
      \item B. \zh{你觉得呢？} (Nǐ jué de ne?) — Correct, poli, ouvert.
      \item C. \zh{快点说。} (Kuài diǎn shuō.) — Impératif, impoli.
      \item D. \zh{不对。} (Bú duì.) — Négation brute, coupe le dialogue.
      \end{itemize}

\item La structure \zh{首先……然后……最后……} (shǒu xiān… rán hòu… zuì hòu…) sert à :
      \begin{itemize}
      \item A. Exprimer une hypothèse.
      \item B. Donner son avis personnel.
      \item C. Structurer un exposé ou une narration chronologique.
      \item D. Formuler une objection.
      \end{itemize}
\end{enumerate}
\end{exercice}

\begin{exercice}[Vrai ou Faux : Justifiez en chinois]
Déterminez si les affirmations sont vraies (V) ou fausses (F). Justifiez chaque réponse par une phrase courte en chinois (caractères + pinyin).

\begin{enumerate}
\item \textbf{Assertion :} Dans une lettre formelle, on écrit la date (\zh{日期} rì qī) avant l'appel (\zh{称呼}).
\item \textbf{Assertion :} La formule \zh{顺颂商祺} (shùn sòng shāng qí) s'utilise pour écrire à un ami proche.
\item \textbf{Assertion :} Pour \textbf{présenter} un exposé oral, \zh{我想说的是……} (wǒ xiǎng shuō de shì…) est plus approprié que \zh{喂，听我说} (wéi, tīng wǒ shuō).
\item \textbf{Assertion :} Le mot-outil \zh{了} (le) placé en fin de phrase indique toujours une action passée terminée.
\end{enumerate}
\end{exercice}

\begin{exercice}[Vocabulaire : Formules fonctionnelles]
Complétez le tableau ci-dessous en écrivant les caractères chinois correspondants aux définitions et au pinyin donnés. Attention aux tons.

\begin{center}
\begin{tabularx}{\linewidth}{|X|X|X|X|}
\hline
\rowcolor{popblueL} \textbf{Caractères} & \textbf{Pinyin} & \textbf{Nature} & \textbf{Traduction} \\
\hline
\zh{尊敬的} & zūn jìng de & Adj. & Respecté (formule d'appel) \\
\hline
\zh{申请} & shēn qǐng & Verbe & Solliciter, demander formellement (ex: stage) \\
\hline
\zh{贵公司} & guì gōng sī & Nom (politesse) & Votre compagnie / société \\
\hline
\zh{此致} & cǐ zhì & Locution & C'est tout / En terminant (début formule finale) \\
\hline
\zh{觉得} & jué de & Verbe & Penser / trouver que (avis) \\
\hline
\zh{同意} & tóng yì & Verbe & Être d'accord \\
\hline
\zh{因为……所以……} & yīn wèi… suǒ yǐ… & Conjonction & Parce que… donc… (argumentation) \\
\hline
\zh{署名} & shǔ míng & Nom/Verbe & Signature / Signer \\
\hline
\end{tabularx}
\end{center}
\end{exercice}

\begin{exercice}[Pinyin et Tons : Association Caractères / Sons]
Reliez chaque groupe de caractères à son pinyin correct (avec tons) en traçant un trait.

\begin{center}
\begin{tabular}{ll}
\textbf{Colonne A (Caractères)} & \textbf{Colonne B (Pinyin)} \\
\hline
1. 尊敬的 & a. cǐ zhì jìng lǐ \\
2. 此致敬礼 & b. zūn jìng de \\
3. 申请 & c. shēn qǐng \\
4. 贵公司 & d. guì gōng sī \\
5. 你觉得呢 & e. nǐ jué de ne \\
\hline
\end{tabular}
\end{center}

\textbf{Correspondances :} 1-b, 2-a, 3-c, 4-d, 5-e.

\textbf{Consigne supplémentaire :} Appliquez la règle du 3e ton sur les mots ci-dessous. Entourez la syllabe qui change de ton (devient 2e ton) et écrivez la prononciation réelle.
\begin{itemize}
\item \zh{很好} (hěn hǎo) $\rightarrow$ \textbf{hé} hǎo (2 + 3)
\item \zh{马上} (mǎ shàng) $\rightarrow$ Pas de changement (3 + 4)
\item \zh{请假} (qǐng jià) $\rightarrow$ Pas de changement (3 + 4)
\item \zh{老板} (lǎo bǎn) $\rightarrow$ \textbf{laó} bǎn (2 + 3)
\end{itemize}
\end{exercice}

\courssub{Je m'entraîne}

\begin{exercice}[Traduction : Thème (Français $\rightarrow$ Chinois)]
Traduisez les phrases suivantes en chinois (caractères simplifiés + pinyin). Utilisez le vocabulaire de la leçon.

\begin{enumerate}
\item Monsieur le Directeur respecté, je vous écris pour solliciter un stage dans votre entreprise.
\item Je pense que rédiger une lettre formelle est plus poli qu'envoyer un courriel.
\item D'abord, on écrit l'appel ; ensuite, le corps de la lettre ; enfin, la formule de politesse et la signature.
\item Je suis d'accord avec toi, cependant il faut respecter le format standard.
\item Parce que c'est un document officiel, donc on doit utiliser un langage formel.
\end{enumerate}
\end{exercice}

\begin{exercice}[Transformation : Nuancer son avis]
Transformez les phrases affirmatives en utilisant la structure demandée. Écrivez le résultat en caractères + pinyin.

\begin{enumerate}
\item \textbf{Phrase de base :} \zh{写信很好。} (Écrire une lettre c'est bien.)
      \begin{itemize}
      \item \textit{Transformez avec :} \zh{我认为……，但是……} (Je pense que… mais…)
      \end{itemize}
\item \textbf{Phrase de base :} \zh{发邮件很快。} (Envoyer un email c'est rapide.)
      \begin{itemize}
      \item \textit{Transformez avec :} \zh{虽然……，可是……} (Bien que… pourtant…)
      \end{itemize}
\item \textbf{Phrase de base :} \zh{格式很重要。} (Le format est important.)
      \begin{itemize}
      \item \textit{Transformez avec :} \zh{因为……所以……} (Parce que… donc…)
      \end{itemize}
\item \textbf{Phrase de base :} \zh{我们要写称呼。} (Nous devons écrire l'appel.)
      \begin{itemize}
      \item \textit{Transformez en question rhétorique pour \textbf{relancer} :} \zh{……，你觉得呢？}
      \end{itemize}
\end{enumerate}
\end{exercice}

\begin{exercice}[Texte à trous : Dialogue modèle - Solliciter un stage]
Complétez le dialogue suivant avec les expressions manquantes choisies dans la liste ci-dessous. Écrivez les \textbf{caractères chinois} dans les cases. Lisez ensuite le dialogue à haute voix en respectant les tons.

\begin{textebox}[Dialogue : Demande de stage (Complété)]
A : \zh{尊敬的王经理，您好！} \\
A : \zh{我叫李华，是喀麦隆大学的学生。我写信是想}\underline{\zh{申请}}\zh{贵公司的实习机会。} \\
A : \zh{我叫李华，是喀麦隆大学的学生。我写信是想}\underline{\zh{申请}}\zh{贵公司的实习机会。} \\
A : \zh{我知道。}\underline{\zh{我认为}}\zh{我有足够的中文基础，而且很努力。} \\
A : \zh{我知道。}\underline{\zh{我认为}}\zh{我有足够的中文基础，而且很努力。} \\
B : \zh{好，那请把简历发给我。\underline{\zh{祝您工作顺利}} } \\
A : \zh{好的，谢谢您！ \underline{\zh{打扰了}} 。}
\end{textebox}

\end{exercice}

\begin{attention}[Registre oral vs écrit : \zh{此致敬礼} vs \zh{祝您工作顺利}]
Dans l'exercice original, \zh{此致敬礼} était proposé pour l'oral (réplique B).
\textbf{Correction pédagogique :} \zh{此致敬礼} s'utilise \textbf{uniquement à l'écrit} (lettres, emails formels).
À l'oral (téléphone, entretien), on utilise : \zh{祝您工作顺利} (Bonne continuation), \zh{打扰了} (Je vous quitte / Excusez le dérangement), \zh{再见} (Au revoir).
Le dialogue corrigé ci-dessus reflète cet usage réel.
\end{attention}



\begin{exercice}[Remise en ordre : Les 5 parties d'une lettre formelle]
Les cinq éléments ci-dessous constituent une lettre formelle chinoise standard. Ils sont mélangés. Numérotez-les de 1 à 5 dans l'ordre logique d'apparition (de haut en bas de la page). Justifiez oralement l'ordre choisi en utilisant \zh{首先}, \zh{然后}, \zh{最后}.

\begin{center}
\begin{tabular}{|c|l|l|}
\hline
\textbf{Ordre} & \textbf{Caractères + Pinyin} & \textbf{Fonction} \\
\hline
1 & \zh{称呼} (chēng hū) & Appel (ex: \zh{尊敬的经理：}) \\
\hline
2 & \zh{问候语} (wèn hòu yǔ) & Formule de salutation initiale (ex: \zh{您好！}) \\
\hline
3 & \zh{正文} (zhèng wén) & Corps de la lettre \\
\hline
4 & \zh{祝颂语} (zhù sòng yǔ) & Formules de politesse finales (ex: \zh{此致 敬礼}) \\
\hline
5 & \zh{署名} (shǔ míng) & Signature (Nom + Date) \\
\hline
\end{tabular}
\end{center}
\end{exercice}

\courssub{Je me prépare au Bac}

\begin{exercice}[Compréhension de l'oral / Expression orale continue : Exposé structuré]
\textbf{Consigne :} Vous devez présenter oralement pendant \textbf{1 minute} la structure type d'une lettre formelle chinoise devant la classe. Vous utiliserez \textbf{obligatoirement} les connecteurs \zh{首先} (shǒu xiān), \zh{然后} (rán hòu), \zh{最后} (zuì hòu) pour structurer votre propos. Vous citerez les 5 parties techniques (voir exercice précédent) et donnerez un exemple de formule pour l'appel et pour la politesse finale.

\textbf{Grille d'évaluation indicative (20 pts) :}
\begin{itemize}
\item Respect du temps (1 min $\pm$ 10 s) : 4 pts
\item Utilisation correcte des 3 connecteurs imposés : 6 pts
\item Contenu exact (5 parties nommées + 2 exemples) : 6 pts
\item Prononciation (tons, liaison, intonation) : 4 pts
\end{itemize}

\textbf{Aide-mémoire (vocabulaire clé) :}
\zh{结构} (jié gòu - structure), \zh{格式} (gé shì - format), \zh{必不可少} (bì bù kě shǎo - indispensable), \zh{示例} (shì lì - exemple).
\end{exercice}

\begin{exercice}[Interaction orale : Jeu de rôle noté - Demande de stage]
\textbf{Consigne :} Par binôme. L'élève A joue le rôle d'un \textbf{étudiant de 1ère A} (Lycee de Yaoundé/Buea/Douala). L'élève B joue le \textbf{Directeur des Ressources Humaines} d'une entreprise sino-camerounaise (\zh{中喀合资公司} Zhōng Kā hé zī gōng sī).

\textbf{Scénario :} L'élève A appelle ou se présente pour \textbf{solliciter oralement un stage d'été} (\zh{暑期实习} shǔ qī shí xí). Il/Elle doit utiliser \textbf{au moins 5 fois} les formules de politesse et le vocabulaire suivant : \zh{尊敬的} (zūn jìng de), \zh{申请} (shēn qǐng), \zh{贵公司} (guì gōng sī), \zh{实习} (shí xí), \zh{机会} (jī huì), \zh{打扰了} (dǎ rǎo le - pour prendre congé).

\textbf{Déroulement suggéré :}
1. Salutations formelles + Présentation (Nom, École, Classe).
2. Objet de la venue / appel (Demande de stage).
3. Argumentation (Pourquoi cette entreprise ? Compétences : chinois, sérieux).
4. Réponse du Directeur (questions, accord de principe, demande de CV).
5. Remerciements + Formule de congé polie.

\textbf{Durée :} 2 à 3 minutes par binôme. Enregistrement audio possible pour l'évaluation.
\end{exercice}

\begin{exercice}[Débat argumenté : Lettre vs Email]
\textbf{Sujet :} \zh{写信比发邮件更好吗？} (Écrire une lettre est-il mieux qu'envoyer un email ?)

\textbf{Consigne :} En groupes de 4 (2 « Pour la lettre », 2 « Pour l'email »). Préparez 5 minutes. Débat de 10 minutes. Chaque élève doit produire \textbf{au moins deux arguments complets} en utilisant la structure \zh{因为……所以……} (yīn wèi… suǒ yǐ…) ou \zh{一方面……另一方面……} (yī fāng miàn… lìng wài yī fāng miàn…). Utilisez les formules de \textbf{relance} (\zh{你觉得呢？}, \zh{对吧？}) pour faire réagir les autres.

\textbf{Pistes de vocabulaire :}
\begin{itemize}
\item \textit{Pour la lettre :} \zh{正式} (zhèng shì - formel), \zh{尊重} (zūn zhòng - respect), \zh{保存} (bǎo cún - conservation/archives), \zh{郑重} (zhèng zhòng - solennel), \zh{签名} (qiān míng - signature manuscrite).
\item \textit{Pour l'email :} \zh{快速} (kuài sù - rapide), \zh{方便} (fāng biàn - pratique), \zh{环保} (huán bǎo - écologique / sans papier), \zh{随时随地} (suí shí suí dì - n'importe quand n'importe où), \zh{附件} (fù jiàn - pièce jointe).
\end{itemize}

\textbf{Exemple de tour de parole :}
\zh{学生甲：我认为写信更好。因为写信很正式，所以表示尊重。你觉得呢？} \\
\zh{学生乙：我不同意。因为发邮件很快速，而且环保，所以更方便。对吧？}
\end{exercice}

\begin{methode}[Réussir l'expression orale au Bac (Check-list)]
\begin{enumerate}
\item \textbf{Préparer des fiches mots-clés} (pas des phrases toutes faites) : Structure, Vocabulaire précis, Connecteurs.
\item \textbf{S'entraîner à voix haute} avec chronomètre (1 min exposé, 2-3 min dialogue).
\item \textbf{Soigner les tons} : 3e ton sandhi (\zh{很好} $\rightarrow$ \zh{hé hǎo}), tons neutres (\zh{了, 的, 呢, 吧}), distinction \zh{j/q/x} vs \zh{zh/ch/sh/r}.
\item \textbf{Utiliser les "béquilles" fonctionnelles} : \zh{那个…} (euh…), \zh{就是说…} (c'est-à-dire…), \zh{请允许我说完} (permettez-moi de finir).
\item \textbf{Regarder l'examinateur / le partenaire}, pas sa feuille.
\end{enumerate}
\end{methode}

\begin{savaistu}[Culture : La lettre en Chine aujourd'hui]
Bien que WeChat (\zh{微信} Wēi xìn) et l'email dominent la communication quotidienne et professionnelle en Chine, la \textbf{lettre formelle sur papier} (\zh{纸质公函} zhǐ zhì gōng hán) reste \textbf{obligatoire} pour :
\begin{itemize}
\item Les documents juridiques / contrats (\zh{合同} hé tong).
\item Les communications officielles entre administrations (\zh{公文} gōng wén).
\item Les demandes d'emploi / stage très formelles (grandes entreprises d'État, ambassades).
\item Les lettres de remerciement / condoléances officielles.
\end{itemize}
Au Cameroun, l'administration et les grandes entreprises (SNEC, CAMTEL, banques, entreprises chinoises comme \zh{中喀合资公司}) exigent encore souvent un courrier physique signé pour les dossiers officiels. Maîtriser le format \zh{称呼-正文-祝颂语-署名} est donc une compétence professionnelle transversale.
\end{savaistu}

\newpage

\coursec{Corrigés des exercices}

\coursec{Leçon 41 — Expression orale : rédaction d’une lettre formelle — Corrigés détaillés}

\begin{propriete}[Rappel : Structure type d'une lettre formelle chinoise]
Une lettre formelle (\zh{正式信函} zhèng shì xìn hán) respecte un ordre immuable, de haut en bas :
\begin{enumerate}
\item \zh{称呼} (chēng hū) — Appel : en haut à gauche, suivi des deux points pleine chasse \zh{：}. Ex. \zh{尊敬的王经理：}.
\item \zh{问候语} (wèn hòu yǔ) — Salutation : ligne suivante, ex. \zh{您好！}.
\item \zh{正文} (zhèng wén) — Corps de la lettre : objet, argumentation, connecteurs logiques (\zh{首先……然后……最后……}, \zh{因为……所以……}).
\item \zh{祝颂语} (zhù sòng yǔ) — Formules de politesse finales : \zh{此致} (cǐ zhì) sur une ligne, \zh{敬礼} (jìng lǐ) sur la ligne d'après, alignées à gauche.
\item \zh{署名} (shǔ míng) — Signature + \zh{日期} (rì qī) : alignés en bas à droite. Ex. \zh{李华 / 二〇二四年十月十五日}.
\end{enumerate}
\emph{Registre oral vs écrit :} \zh{此致敬礼} s'écrit seulement ; à l'oral on utilise \zh{祝您工作顺利}, \zh{打扰了}, \zh{再见}.
\end{propriete}

\begin{propriete}[Formules fonctionnelles clés pour l'oral (exposé, débat, relance)]
\begin{center}
\begin{tabularx}{\linewidth}{|X|X|X|X|}
\hline
\rowcolor{popblueL} \textbf{Fonction} & \textbf{Chinois} & \textbf{Pinyin} & \textbf{Français} \\
\hline
Ouvrir un exposé & \zh{今天我要谈谈……} & jīn tiān wǒ yào tán tan… & Aujourd'hui je vais parler de… \\
Structurer (étape 1) & \zh{首先……} & shǒu xiān… & D'abord… \\
Structurer (étape 2) & \zh{然后……} & rán hòu… & Ensuite… \\
Structurer (étape 3) & \zh{最后……} & zuì hòu… & Enfin… \\
Donner son avis & \zh{我认为…… / 我觉得……} & wǒ rèn wéi… / wǒ jué de… & Je pense que… / Je trouve que… \\
Argumenter (cause) & \zh{因为……所以……} & yīn wèi… suǒ yǐ… & Parce que… donc… (les deux s'emploient ensemble) \\
Concéder / Nuancer & \zh{虽然……但是/可是……} & suī rán… dàn shì / kě shì… & Bien que… cependant… \\
Relancer l'interlocuteur & \zh{你觉得呢？ / 对吧？} & nǐ jué de ne? / duì ba? & Qu'en penses-tu ? / N'est-ce pas ? \\
Exprimer accord/désaccord & \zh{我同意 / 我不同意} & wǒ tóng yì / wǒ bù tóng yì & Je suis d'accord / Je ne suis pas d'accord \\
Prendre congé (oral) & \zh{祝您工作顺利 / 打扰了} & zhù nín gōng zuò shùn lì / dǎ rǎo le & Bonne continuation / Je vous ai dérangé \\
\hline
\end{tabularx}
\end{center}
\end{propriete}

\begin{corrige}[Exercice 1 — QCM : Structure d'une lettre formelle]
\begin{enumerate}
\item \textbf{B.} La formule finale \zh{此致敬礼} (cǐ zhì jìng lǐ) se place \textbf{juste avant la signature}, alignée à gauche, sur deux lignes distinctes : \zh{此致} sur la première ligne, \zh{敬礼} sur la suivante. La signature (\zh{署名} shǔ míng) et la date (\zh{日期} rì qī) viennent après, alignées à droite. Ordre complet : \zh{称呼} $\rightarrow$ \zh{问候语} $\rightarrow$ \zh{正文} $\rightarrow$ \zh{此致敬礼} $\rightarrow$ \zh{署名、日期}.
\item \textbf{B.} L'appel (\zh{称呼} chēng hū) s'écrit \textbf{en haut à gauche}, sur la toute première ligne, suivi obligatoirement des deux points chinois pleine chasse \zh{：}. Exemple : \zh{尊敬的王经理：} (Zūn jìng de Wáng jīng lǐ :) « Respecté Directeur Wang : ». En chinois, on n'utilise \textbf{jamais} la virgule après l'appel dans une lettre formelle.
\item \textbf{B.} \zh{你觉得呢？} (Nǐ jué de ne?) « Qu'en penses-tu ? » est la formule de relance standard, polie et ouverte. \zh{你说什么？} est trop direct (familier), \zh{快点说} est un impératif impoli, \zh{不对} coupe brutalement le dialogue.
\item \textbf{C.} \zh{首先……然后……最后……} (shǒu xiān… rán hòu… zuì hòu…) « d'abord… ensuite… enfin… » sert à \textbf{structurer} un exposé ou une procédure chronologique. Exemple : \zh{首先写称呼，然后写正文，最后写署名。} (Shǒu xiān xiě chēng hū, rán hòu xiě zhèng wén, zuì hòu xiě shǔ míng.) « D'abord on écrit l'appel, ensuite le corps, enfin la signature. »
\end{enumerate}
\end{corrige}

\begin{attention}[Point de vigilance]
Ne pas confondre la place de \zh{此致敬礼} (avant la signature, \textbf{à gauche}) avec celle de la date (après la signature, \textbf{à droite}). C'est l'erreur n°1 en examen.
\end{attention}


\begin{corrige}[Exercice 2 — Vrai ou Faux : Conventions épistolaires et grammaire]
\begin{enumerate}
\item \textbf{Faux.} La date se place \textbf{à la fin} de la lettre, sous la signature, alignée à droite — jamais avant l'appel. Justification : \zh{日期写在最后，署名的下面。} (Rì qī xiě zài zuì hòu, shǔ míng de xià miàn.) « La date s'écrit à la fin, sous la signature. »
\item \textbf{Faux.} \zh{顺颂商祺} (shùn sòng shāng qí) est une formule de politesse \textbf{commerciale} (\zh{商} shāng = commerce) ; elle ne s'adresse pas à un ami proche. Justification : \zh{顺颂商祺用于商务信函，不用于朋友。} (Shùn sòng shāng qí yòng yú shāng wù xìn hán, bù yòng yú péng you.) « \zh{顺颂商祺} s'utilise dans la correspondance d'affaires, pas entre amis. »
\item \textbf{Vrai.} \zh{我想说的是……} (wǒ xiǎng shuō de shì…) est une ouverture d'exposé appropriée, formelle ; \zh{喂，听我说} (wéi, tīng wǒ shuō) est familier et impoli. Justification : \zh{“我想说的是……”很正式；“喂，听我说”很不礼貌。} (« … » hěn zhèng shì ; « … » hěn bù lǐ mào.)
\item \textbf{Faux.} Le \zh{了} (le) en fin de phrase marque surtout un \textbf{changement de situation} ou un fait nouveau (aspect perfectif / modal), pas seulement le passé. Justification : \zh{句尾的“了”表示情况的变化，不只表示过去。} (Jù wěi de “le” biǎo shì qíng kuàng de biàn huà, bù zhǐ biǎo shì guò qù.) Exemple : \zh{下雨了。} (Xià yǔ le.) « Il se met à pleuvoir / Il pleut (maintenant). »
\end{enumerate}
\end{corrige}

\begin{corrige}[Exercice 3 — Vocabulaire : Formules fonctionnelles (Tableau complété)]
\begin{center}
\begin{tabularx}{\linewidth}{|X|X|X|X|}
\hline
\rowcolor{popblueL} \textbf{Caractères} & \textbf{Pinyin} & \textbf{Nature} & \textbf{Traduction} \\
\hline
\zh{尊敬的} & zūn jìng de & Adj. (formule) & Respecté (appel formel) \\
\hline
\zh{申请} & shēn qǐng & Verbe (trans. direct) & Solliciter, demander formellement \\
\hline
\zh{贵公司} & guì gōng sī & Nom (honorifique) & Votre société (honorifique) \\
\hline
\zh{此致} & cǐ zhì & Locution (écrit) & En terminant (début formule finale) \\
\hline
\zh{觉得} & jué de & Verbe & Penser, trouver que (avis subjectif) \\
\hline
\zh{同意} & tóng yì & Verbe & Être d'accord (avec qqn / qqch) \\
\hline
\zh{因为……所以……} & yīn wèi… suǒ yǐ… & Conjonctions corrélatives & Parce que… donc… (paire inséparable) \\
\hline
\zh{署名} & shǔ míng & Nom / Verbe & Signature / Signer \\
\hline
\end{tabularx}
\end{center}
\end{corrige}

\begin{attention}[Points de vigilance lexicaux]
\begin{itemize}
\item \zh{贵} (guì, 4e ton) est un \textbf{préfixe honorifique} : \zh{贵公司} « votre honorable société », \zh{贵姓} « votre nom honorable ». Ne pas confondre avec l'adjectif « cher / coûteux » (\zh{很贵} hěn guì).
\item \zh{申请} (shēn qǐng) est un verbe \textbf{transitif direct} : \zh{申请实习机会} « solliciter une opportunité de stage », \textbf{sans préposition} (pas de *申请给*).
\item \zh{觉得} (jué de) : le deuxième caractère \zh{得} se prononce au \textbf{ton neutre} (de) ; ne pas écrire *jué dé* ni prononcer *jué dé*.
\item \zh{因为……所以……} : en chinois, les deux connecteurs \textbf{coexistent} dans la phrase complexe (structure standard), contrairement au français où « parce que… donc… » est fautif.
\end{itemize}
\end{attention}


\begin{corrige}[Exercice 4 — Pinyin et Tons : Association et Sandhi]
\textbf{Correspondances (Caractères $\rightarrow$ Pinyin) :}
\begin{enumerate}
\item \zh{尊敬的} — \textbf{b.} zūn jìng de (1 + 4 + neutre ; \textbf{aucun 3e ton}, donc aucune sandhi).
\item \zh{此致敬礼} — \textbf{a.} cǐ zhì jìng lǐ (3 + 4 + 4 + 3).
\item \zh{申请} — \textbf{c.} shēn qǐng (1 + 3).
\item \zh{贵公司} — \textbf{d.} guì gōng sī (4 + 1 + 1).
\item \zh{你觉得呢} — \textbf{e.} nǐ jué de ne (3 + 2 + neutre + neutre).
\end{enumerate}

\textbf{Application de la règle du 3e ton (Sandhi) — Rappel :} Deux 3e tons consécutifs $\rightarrow$ le premier devient 2e ton (montant). L'écriture pinyin \textbf{conserve} le 3e ton dictionnaire.
\begin{center}
\begin{tabularx}{\linewidth}{|X|X|X|X|}
\hline
\rowcolor{poporangeL} \textbf{Mot} & \textbf{Pinyin dictionnaire} & \textbf{Pinyin prononcé (sandhi)} & \textbf{Syllabe modifiée} \\
\hline
\zh{很好} & hěn hǎo (3+3) & \textbf{hén hǎo} (2+3) & \zh{很} hěn $\rightarrow$ hén \\
\hline
\zh{老板} & lǎo bǎn (3+3) & \textbf{láo bǎn} (2+3) & \zh{老} lǎo $\rightarrow$ láo \\
\hline
\zh{马上} & mǎ shàng (3+4) & mǎ shàng (3+4) & \textbf{Pas de changement} (2e syllabe $\neq$ 3e ton) \\
\hline
\zh{请假} & qǐng jià (3+4) & qǐng jià (3+4) & \textbf{Pas de changement} \\
\hline
\end{tabularx}
\end{center}
\emph{Note :} \zh{你} (nǐ) devant \zh{好} (hǎo) devient \textbf{ní hǎo} ; devant \zh{觉得} (jué de, 2e ton) \textbf{pas de sandhi} (nǐ jué de).
\end{corrige}

\begin{corrige}[Exercice 5 — Traduction : Thème (Français $\rightarrow$ Chinois)]
\begin{enumerate}
\item \zh{尊敬的经理：我写信是想申请贵公司的实习机会。} \\
(Zūn jìng de jīng lǐ : wǒ xiě xìn shì xiǎng shēn qǐng guì gōng sī de shí xí jī huì.) \\
\emph{Structure :} \zh{我写信是想……} « Je vous écris dans l'intention de… » ; \zh{贵公司} honorifique pour « votre société ».
\item \zh{我认为写正式信比发邮件更客气。} \quad \emph{(Variante : \zh{更正式} gèng zhèng shì « plus formel »)} \\
(Wǒ rèn wéi xiě zhèng shì xìn bǐ fā yóu jiàn gèng kè qi.) \\
\emph{Comparatif :} Sujet A + \zh{比} bǐ + Sujet B + \zh{更} gèng + Adjectif.
\item \zh{首先写称呼，然后写正文，最后写祝颂语和署名。} \\
(Shǒu xiān xiě chēng hū, rán hòu xiě zhèng wén, zuì hòu xiě zhù sòng yǔ hé shǔ míng.) \\
\emph{Connecteurs imposés :} \zh{首先……然后……最后……}. \zh{祝颂语} inclut \zh{此致敬礼}.
\item \zh{我同意你的看法，但是必须遵守标准格式。} \\
(Wǒ tóng yì nǐ de kàn fǎ, dàn shì bì xū zūn shǒu biāo zhǔn gé shì.) \\
\zh{同意} tóng yì + [nom / proposition] ; \zh{必须} bì xū « devoir, falloir » (obligation forte).
\item \zh{因为是公文，所以必须使用正式语言。} \\
(Yīn wèi shì gōng wén, suǒ yǐ bì xū shǐ yòng zhèng shì yǔ yán.) \\
\textbf{Point clé :} En chinois, \zh{因为……所以……} \textbf{doivent} apparaître ensemble pour marquer la relation cause-effet. Ne pas traduire mot à mot le français « Parce que c'est un document officiel, il faut… » en omettant \zh{所以}.
\end{enumerate}
\end{corrige}

\begin{corrige}[Exercice 6 — Transformation : Nuancer son avis (Structures complexes)]
\begin{enumerate}
\item \zh{我认为写信很好，但是发邮件也很方便。} \\
(Wǒ rèn wéi xiě xìn hěn hǎo, dàn shì fā yóu jiàn yě hěn fāng biàn.) \\
\emph{Structure :} Avis (\zh{我认为}) + Concession (\zh{但是} / \zh{可是}) + \zh{也} yě « aussi ».
\item \zh{虽然发邮件很快，可是写信更正式。} \\
(Suī rán fā yóu jiàn hěn kuài, kě shì xiě xìn gèng zhèng shì.) \\
\emph{Structure :} \zh{虽然……可是……} (paire corrélative obligatoire). \zh{虽然} en tête de proposition subordonnée, \zh{可是} / \zh{但是} en tête de principale.
\item \zh{因为格式很重要，所以我们要仔细检查。} \\
(Yīn wèi gé shì hěn zhòng yào, suǒ yǐ wǒ men yào zǐ xì jiǎn chá.) \\
\emph{Structure :} \zh{因为……所以……} (paire corrélative). \zh{仔细} zǐ xì « soigneusement, en détail ».
\item \zh{我们要写称呼，你觉得呢？} \\
(Wǒ men yào xiě chēng hū, nǐ jué de ne?) \\
\emph{Relance :} La particule finale \zh{呢} (ne) transforme l'affirmation en question ouverte / relance interactive.
\end{enumerate}
\end{corrige}

\begin{attention}[Piège francophone : paires de connecteurs]
En français, on dit « Bien que ce soit rapide, c'est formel » (pas de second mot). En chinois, \zh{虽然} \textbf{appelle obligatoirement} \zh{可是} ou \zh{但是} ; \zh{因为} appelle \zh{所以}. L'omission du second élément rend la phrase incomplète / incorrecte.
\end{attention}


\begin{corrige}[Exercice 7 — Texte à trous : Dialogue « Demande de stage » (Registre oral formel)]
\textbf{Réponses dans l'ordre des blancs :}
\begin{enumerate}
\item \zh{申请} (shēn qǐng) — Verbe transitif direct : \zh{申请贵公司的实习机会}. Pas de préposition.
\item \zh{你觉得呢} (nǐ jué de ne) — Relance du directeur : « Qu'en pensez-vous ? » (interaction orale).
\item \zh{我认为} (wǒ rèn wéi) — Introduction de l'argumentation de l'étudiant : \zh{我认为我有足够的中文基础}.
\item \zh{祝您工作顺利} (zhù nín gōng zuò shùn lì) — Formule de congé \textbf{orale} polie (supérieur $\rightarrow$ subalterne / pair) : « Bonne continuation dans votre travail. »
\item \zh{打扰了} (dǎ rǎo le) — Formule de politesse de l'étudiant pour prendre congé : « Je vous ai dérangé / Excusez le dérangement. »
\end{enumerate}

\textbf{Dialogue complet à lire à voix haute (modèle de prononciation) :}
\begin{textebox}[Dialogue : Demande de stage — Oral formel]
\zh{A：尊敬的王经理，您好！} \\
\zh{B：您好。请问有什么事？} \\
\zh{A：我叫李华，是雅温得中学一年级的学生。我写信是想申请贵公司的实习机会。} \\
\zh{B：哦？你觉得呢？我们通常在三月招聘实习生。} \\
\zh{A：我知道。我认为我有足够的中文基础，而且很努力。} \\
\zh{B：好，那请把简历发给我。祝您工作顺利。} \\
\zh{A：好的，谢谢您！打扰了。}
\end{textebox}

\end{corrige}

\begin{attention}[Registre oral vs écrit \& Tons]
\begin{itemize}
\item \zh{此致敬礼} est \textbf{réservé à l'écrit} ; à l'oral : \zh{祝您工作顺利}, \zh{打扰了}, \zh{再见}.
\item \zh{您好} se prononce \textbf{nín hǎo} (\zh{您} nín est déjà 2e ton, pas de sandhi).
\item \zh{很好} se prononce \textbf{hén hǎo} (sandhi 3+3 $\rightarrow$ 2+3).
\item \zh{雅温得} (Yǎ wēn dé) : 3 + 1 + 2, pas de sandhi.
\end{itemize}
\end{attention}


\begin{corrige}[Exercice 8 — Remise en ordre : Les 5 parties d'une lettre formelle]
\textbf{Ordre correct (de haut en bas de la page) :}
\begin{enumerate}
\item \zh{称呼} (chēng hū) — Appel, en haut à gauche + \zh{：} (ex. \zh{尊敬的经理：}).
\item \zh{问候语} (wèn hòu yǔ) — Salutation, ligne suivante (ex. \zh{您好！}).
\item \zh{正文} (zhèng wén) — Corps de la lettre (objet + argumentation).
\item \zh{祝颂语} (zhù sòng yǔ) — Formules finales : \zh{此致} puis \zh{敬礼} sur deux lignes, à gauche.
\item \zh{署名} (shǔ míng) — Signature puis date, alignées à droite (ex. \zh{李华 / 二〇二四年十月十五日}).
\end{enumerate}

\textbf{Justification orale modèle (à réciter) :}
\zh{首先写称呼，然后写问候语，然后写正文，然后写祝颂语，最后写署名和日期。} \\
(Shǒu xiān xiě chēng hū, rán hòu xiě wèn hòu yǔ, rán hòu xiě zhèng wén, rán hòu xiě zhù sòng yǔ, zuì hòu xiě shǔ míng hé rì qī.)
\end{corrige}

\begin{corrige}[Exercice 9 — Expression orale continue : Exposé structuré (Type Bac — 1 min)]
\textbf{Proposition d'exposé modèle :}
\begin{textebox}[Exposé modèle : Structure de la lettre formelle]
\zh{今天我要谈谈中文正式信函的结构。} (Jīn tiān wǒ yào tán tan Zhōng wén zhèng shì xìn hán de jié gòu.) \\
\zh{一封正式信有五个必不可少的部分。} (Yī fēng zhèng shì xìn yǒu wǔ gè bì bù kě shǎo de bù fen.) \\
\zh{首先，在左上角写称呼，例如：“尊敬的经理：”。} (Shǒu xiān, zài zuǒ shàng jiǎo xiě chēng hū, lì rú : “Zūn jìng de jīng lǐ :”.) \\
\zh{然后写问候语：“您好！”。} (Rán hòu xiě wèn hòu yǔ : “Nín hǎo!”.) \\
\zh{然后写正文，说明写信的目的。} (Rán hòu xiě zhèng wén, shuō míng xiě xìn de mù dì.) \\
\zh{然后写祝颂语：“此致敬礼！”。} (Rán hòu xiě zhù sòng yǔ : “Cǐ zhì jìng lǐ!”.) \\
\zh{最后，在右下角写署名和日期。} (Zuì hòu, zài yòu xià jiǎo xiě shǔ míng hé rì qī.) \\
\zh{总之，正式信的格式非常重要。以上就是我的发言，谢谢大家！} (Zǒng zhī, zhèng shì xìn de gé shì fēi cháng zhòng yào. Yǐ shàng jiù shì wǒ de fā yán, xiè xie dà jiā!)
\end{textebox}

\textbf{Barème indicatif (20 pts) :}
\begin{itemize}
\item Respect du temps (1 min $\pm$ 10 s) : 4 pts.
\item Utilisation correcte des 3 connecteurs \zh{首先 / 然后 / 最后} : 6 pts (2 pts chacun).
\item Contenu exact : 5 parties nommées (4 pts) + 2 exemples de formules (2 pts) = 6 pts.
\item Prononciation : tons corrects (sandhi \zh{很好} hén hǎo, \zh{左上} zuǒ shàng 3+4 pas de sandhi), intonation, fluidité : 4 pts.
\end{itemize}
\end{corrige}

\begin{attention}[Consignes d'examen]
Ne pas lire sa fiche ; annoncer le sujet dès la première phrase avec \zh{今天我要谈谈……} ; soigner l'articulation des connecteurs (\zh{首先}… \zh{然后}… \zh{最后}…).
\end{attention}


\begin{corrige}[Exercice 10 — Interaction orale : Jeu de rôle « Demande de stage » (2 pers.)]
\textbf{Dialogue modèle (rédaction attendue) :}
\begin{textebox}[Jeu de rôle : Étudiant (A) / Directeur (B)]
\zh{A（学生）：尊敬的王经理，您好！我叫李华，是雅温得中学一年级的学生。} (Zūn jìng de Wáng jīng lǐ, nín hǎo! Wǒ jiào Lǐ Huá, shì Yǎ wēn dé zhōng xué yī nián jí de xué sheng.) \\
\zh{B（经理）：您好，李同学。请坐。请问有什么事？} (Nín hǎo, Lǐ tóng xué. Qǐng zuò. Qǐng wèn yǒu shén me shì?) \\
\zh{A：我想申请贵公司的暑期实习机会。} (Wǒ xiǎng shēn qǐng guì gōng sī de shǔ qī shí xí jī huì.) \\
\zh{B：哦？你为什么选择我们公司？} (Ó? Nǐ wèi shén me xuǎn zé wǒ men gōng sī?) \\
\zh{A：因为贵公司是中喀合资的知名企业，所以我很想来学习。我认为我的中文基础不错，而且工作认真。} (Yīn wèi guì gōng sī shì Zhōng Kā hé zī de zhī míng qǐ yè, suǒ yǐ wǒ hěn xiǎng lái xué xí. Wǒ rèn wéi wǒ de Zhōng wén jī chǔ bú cuò, ér qiě gōng zuò rèn zhēn.) \\
\zh{B：很好。请把你的简历发给我们，我们会尽快回复。} (Hěn hǎo. Qǐng bǎ nǐ de jiǎn lì fā gěi wǒ men, wǒ men huì jǐn kuài huí fù.) \\
\zh{A：好的，非常感谢您！祝您工作顺利，打扰了！} (Hǎo de, fēi cháng gǎn xiè nín! Zhù nín gōng zuò shùn lì, dǎ rǎo le!) \\
\zh{B：不客气，再见！} (Bú kè qi, zài jiàn!)
\end{textebox}

\textbf{Barème indicatif (20 pts) :}
\begin{itemize}
\item Respect du scénario et des 5 étapes (salutation, demande, justification, conclusion, congé) : 5 pts.
\item Vocabulaire imposé utilisé au moins 5 fois (\zh{尊敬的}, \zh{申请}, \zh{贵公司}, \zh{实习}, \zh{机会}, \zh{打扰了}, \zh{因为……所以……}, \zh{我认为}) : 5 pts.
\item Formules de politesse et registre formel adapté (\zh{您}, \zh{贵公司}, pas de \zh{此致敬礼} à l'oral) : 4 pts.
\item Argumentation structurée (\zh{因为……所以……}, \zh{我认为……}, \zh{而且}) : 3 pts.
\item Prononciation et tons (sandhi \zh{很好} hén hǎo, \zh{雅温得} Yǎ wēn dé) : 3 pts.
\end{itemize}
\end{corrige}

\begin{attention}[Points de vigilance jeu de rôle]
\begin{itemize}
\item Saluer avec \zh{您} (nín) et non \zh{you} face au directeur.
\item Ne \textbf{jamais} dire \zh{此致敬礼} à l'oral.
\item Terminer par \zh{打扰了} ou \zh{再见} (pas \zh{拜拜} trop familier).
\item \zh{中喀合资} (Zhōng Kā hé zī) : vocabulaire contextuel Cameroun-Chine.
\end{itemize}
\end{attention}


\begin{corrige}[Exercice 11 — Débat argumenté : Lettre formelle vs Email]
\textbf{Arguments modèles — Camp « Pour la lettre formelle » :}
\begin{itemize}
\item \zh{我认为写信更好。因为写信很正式，所以表示对对方的尊重。} (Wǒ rèn wéi xiě xìn gèng hǎo. Yīn wèi xiě xìn hěn zhèng shì, suǒ yǐ biǎo shì duì duì fāng de zūn zhòng.)
\item \zh{一方面，纸质的信可以长期保存；另一方面，手写的签名有法律效力。} (Yī fāng miàn, zhǐ zhì de xìn kě yǐ cháng qī bǎo cún; lìng wài yī fāng miàn, shǒu xiě de qiān míng yǒu fǎ lǜ xiào lì.) \emph{Structure : \zh{一方面……另一方面……}}.
\end{itemize}

\textbf{Arguments modèles — Camp « Pour l'email » :}
\begin{itemize}
\item \zh{我不同意。因为发邮件很快速，所以非常方便。} (Wǒ bù tóng yì. Yīn wèi fā yóu jiàn hěn kuài sù, suǒ yǐ fēi cháng fāng biàn.)
\item \zh{虽然写信很郑重，可是发邮件更环保，还可以随时随地发附件。} (Suī rán xiě xìn hěn zhèng zhòng, kě shì fā yóu jiàn gèng huán bǎo, hái kě yǐ suí shí suí dì fā fù jiàn.) \emph{Structure : \zh{虽然……可是……}}.
\end{itemize}

\textbf{Conclusion de synthèse possible (animateur) :}
\zh{看来要看场合：正式场合必须写信，日常交流可以发邮件。} (Kàn lái yào kàn chǎng hé : zhèng shì chǎng hé bì xū xiě xìn, rì cháng jiāo liú kě yǐ fā yóu jiàn.)

\textbf{Barème indicatif (20 pts) :}
\begin{itemize}
\item Au moins deux arguments complets par élève : 6 pts.
\item Structures imposées correctement employées (\zh{因为……所以……}, \zh{虽然……可是……}, \zh{一方面……另一方面……}) : 6 pts.
\item Relances et interaction (\zh{你觉得呢？}, \zh{对吧？}, \zh{我不同意, 因为……}) : 4 pts.
\item Prononciation, tons, fluidité : 4 pts.
\end{itemize}
\end{corrige}

\begin{attention}[Éthique du débat \& Phonétique]
\begin{itemize}
\item Ne pas opposer brutalement (\zh{不对！}, \zh{你错了}) ; nuancer : \zh{我不太赞同，因为……} (Wǒ bù tài zàn tóng, yīn wèi…).
\item Garder la paire \zh{虽然……可是……} complète.
\item Tons : \zh{环保} (huán bǎo, 2+3) ; \zh{保存} (bǎo cún, 3+2 : \textbf{pas de sandhi} car 2e syllabe = 2e ton).
\end{itemize}
\end{attention}


\newpage

\coursec{Fiche bilan}

\begin{popfigure}
\begin{tikzpicture}[mindmap, grow cyclic, every node/.style={concept, font=\small},
  root concept/.append style={concept color=popblue, font=\bfseries},
  level 1/.append style={level distance=3.2cm, sibling angle=60},
  level 2/.append style={level distance=2.2cm, font=\scriptsize}]
\node[root concept]{Lettre formelle 正式信函}
  child[concept color=poporange]{ node {Structure 结构}
    child { node[align=center]{称呼 chēnghu\\ salutation} }
    child { node[align=center]{正文 zhèngwén\\ corps} }
    child { node[align=center]{此致敬礼 cǐ zhì jìnglǐ\\ formule finale} }
  }
  child[concept color=popgreen]{ node {Présenter 介绍}
    child { node[align=center]{我是… Wǒ shì…\\ je suis…} }
    child { node[align=center]{我想… Wǒ xiǎng…\\ je voudrais…} }
  }
  child[concept color=poppurple]{ node {Donner son avis 看法}
    child { node[align=center]{我认为… Wǒ rènwéi…\\ je pense que…} }
    child { node[align=center]{我觉得… Wǒ juéde…\\ il me semble…} }
  }
  child[concept color=poppink]{ node {Relancer 请求}
    child { node[align=center]{请您… Qǐng nín…\\ veuillez…} }
    child { node {希望收到回信 xīwàng shōudào huíxìn} }
  }
  child[concept color=popgold]{ node {Les tons 声调}
    child { node[align=center]{2\textsuperscript{e} + 3\textsuperscript{e} ton\\ nǐ hǎo, qǐngwèn} }
    child { node[align=center]{ton neutre\\ 吗 ma, 的 de} }
  }
  child[concept color=popblue!60]{ node {Débat 讨论}
    child { node[align=center]{同意 tóngyì / 反对 fǎnduì\\ d'accord / contre} }
  };
\end{tikzpicture}
\legende{Carte mentale de la leçon : rédiger et présenter une lettre formelle en chinois.}
\end{popfigure}

\begin{bilan}
\textbf{1. Lexique clé à connaître par cœur}

\begin{center}
\begin{tabularx}{\linewidth}{|X|X|X|X|}
\hline
\rowcolor{popblueL} Caractères & Pinyin & Nature & Traduction \\
\hline
信 & xìn & nom & lettre \\
\hline
写信 & xiě xìn & verbe + compl. & écrire une lettre \\
\hline
尊敬的 & zūnjìng de & adj. & honorable, cher (formel) \\
\hline
称呼 & chēnghu & nom & formule d'appel \\
\hline
正文 & zhèngwén & nom & corps de la lettre \\
\hline
申请 & shēnqǐng & verbe / nom & demander ; candidature \\
\hline
我认为 & wǒ rènwéi & expression & je pense que \\
\hline
请您 & qǐng nín & expression & veuillez (poli) \\
\hline
回信 & huíxìn & nom / verbe & réponse ; répondre \\
\hline
此致敬礼 & cǐ zhì jìnglǐ & formule & veuillez agréer mes salutations \\
\hline
日期 & rìqī & nom & date \\
\hline
签名 & qiānmíng & nom / verbe & signature ; signer \\
\hline
\end{tabularx}
\end{center}

\textbf{2. Structures à connaître par cœur}

\begin{center}
\begin{tabularx}{\linewidth}{|X|X|X|}
\hline
\rowcolor{popblueL} Structure & Exemple (caractères + pinyin) & Traduction \\
\hline
尊敬的 + nom + ： & 尊敬的王老师： Zūnjìng de Wáng lǎoshī: & Cher Monsieur Wang, \\
\hline
我是 + identité & 我是雅温得中学的学生。 Wǒ shì Yǎwēndé zhōngxué de xuésheng. & Je suis élève du lycée de Yaoundé. \\
\hline
我想 / 我要 + verbe & 我想申请奖学金。 Wǒ xiǎng shēnqǐng jiǎngxuéjīn. & Je voudrais demander une bourse. \\
\hline
我认为 + phrase & 我认为中文很有用。 Wǒ rènwéi Zhōngwén hěn yǒuyòng. & Je pense que le chinois est très utile. \\
\hline
请您 + verbe & 请您给我回信。 Qǐng nín gěi wǒ huíxìn. & Veuillez me répondre. \\
\hline
希望 + phrase & 希望早日收到您的回信。 Xīwàng zǎorì shōudào nín de huíxìn. & J'espère recevoir bientôt votre réponse. \\
\hline
此致 / 敬礼 (fin) & 此致敬礼！ Cǐ zhì jìnglǐ! & Veuillez agréer mes salutations. \\
\hline
\end{tabularx}
\end{center}

\textbf{3. Méthode express : rédiger une lettre formelle en 5 étapes}
\begin{enumerate}
\item En haut : la \cle{formule d'appel} 尊敬的 + nom/titre + deux-points ： (jamais de virgule).
\item Corps : se \cle{présenter} (我是…), dire le \cle{but} (我想…), \cle{argumenter} (我认为…).
\item \cle{Relancer} poliment : 请您… / 希望收到您的回信。
\item Formule finale : 此致敬礼！ puis la date (日期) et la signature (签名) en bas à droite.
\item À l'oral : lire à voix haute en soignant les tons, surtout 2\textsuperscript{e} + 3\textsuperscript{e} ton (qǐngwèn, nǐhǎo) et le ton neutre (吗, 的).
\end{enumerate}

\textbf{4. Pour le débat guidé}
\begin{itemize}
\item Donner son avis : 我认为… / 我觉得… / 对我来说… (duì wǒ láishuō, pour moi).
\item Être d'accord : 我同意你的看法。 Wǒ tóngyì nǐ de kànfǎ. (Je suis d'accord avec toi.)
\item Ne pas être d'accord : 我不同意，因为… Wǒ bù tóngyì, yīnwèi… (Je ne suis pas d'accord, parce que…)
\item Relancer la parole : 你怎么看？ Nǐ zěnme kàn? (Qu'en penses-tu ?)
\end{itemize}
\end{bilan}

\begin{aretenir}[Checklist avant le Bac]
\begin{itemize}
\item[$\square$] Je sais écrire la formule d'appel 尊敬的 + nom + ： sans faute.
\item[$\square$] Je connais la structure d'une lettre formelle : appel → corps → 此致敬礼 → date et signature.
\item[$\square$] Je sais me présenter en une phrase : 我是 + identité.
\item[$\square$] Je sais exprimer le but de ma lettre avec 我想 / 我要 + verbe.
\item[$\square$] Je sais donner mon avis avec 我认为… et 我觉得…
\item[$\square$] Je sais relancer poliment avec 请您… et 希望收到您的回信。
\item[$\square$] Je place correctement la date (年 / 月 / 日) et la signature en bas à droite.
\item[$\square$] Je prononce correctement le 3\textsuperscript{e} ton et la suite 2\textsuperscript{e} + 3\textsuperscript{e} ton (qǐngwèn, lǎoshī… non : lǎoshī = 1\textsuperscript{er} ton !).
\item[$\square$] Je n'oublie pas le ton neutre sur 吗, 的, 了.
\item[$\square$] Je sais dire « d'accord / pas d'accord » : 我同意 / 我不同意，因为…
\item[$\square$] Je peux lire ma lettre à voix haute, lentement, avec les bons tons.
\item[$\square$] J'ai mémorisé au moins dix mots du lexique de la lettre formelle.
\end{itemize}
\end{aretenir}

\findecours

\end{document}
